\documentclass[a4paper,11pt,reqno]{amsart}
\usepackage[utf8]{inputenc}
\usepackage{enumerate}
\usepackage{amsmath, amssymb,amsthm}
\usepackage{mathtools}
\usepackage{leftidx}
\usepackage{setspace}
\numberwithin{equation}{section}
%\usepackage{mathbbol}
\setcounter{tocdepth}{1} %2=show subsections, 1= show only sections
\usepackage[left=3cm, right=3cm]{geometry}
\usepackage{hyperref}
\usepackage{xcolor}
%\definecolor{name}{model}{color-spec}
\definecolor{my-black}{rgb}{0,0,0}
\definecolor{my-blue}{rgb}{0,0,0.8}
\definecolor{my-red}{rgb}{0.8,0,0} 
\definecolor{my-green}{rgb}{0,0.5,0}
\hypersetup{colorlinks,
    linkcolor	= {my-blue},
    citecolor	= {my-red},
    urlcolor		= {my-black}
}
\usepackage{graphicx}
\usepackage{tikz-cd}
\onehalfspacing
%Javier's suggestion to avoid overfull boxes
%\overfullrule=10cm better not to use because ArXiv might mess up things
%\usepackage{microtype}

\theoremstyle{plain} %italic

\newtheorem{lemma}{Lemma}[section]
\newtheorem{theorem}{Theorem}
\newtheorem{corollary}{Corollary}[section]
\newtheorem{proposition}{Proposition}[section]
\newtheorem{assumption}{Assumption}[section]
\newtheorem{problem}{Problem}[section]
\newtheorem{consequence}{Consequence}[section]
\newtheorem*{theorem*}{Theorem} %to avoid numbering


\theoremstyle{definition} %non-italic
\newtheorem{remark}{Remark}[section]

\newtheorem{definition}{Definition}[section]
\newtheorem{notation}{Notation}[section]
\newtheorem{conjecture}{Conjecture}[section]
\newtheorem{example}{Example}[section]
\newtheorem{claim}{Claim}[section]

\theoremstyle{remark}
\newtheorem*{remark-non}{Remark}

\DeclareMathOperator{\id}{id}
\DeclareMathOperator{\real}{Re}
\DeclareMathOperator{\imag}{Im}
\DeclareMathOperator{\image}{im}
\DeclareMathOperator{\supp}{supp}
\DeclareMathOperator{\vol}{vol}
\DeclareMathOperator{\inv}{inv}
\DeclareMathOperator{\Tr}{Tr}
\DeclareMathOperator{\ord}{ord}
\DeclareMathOperator{\Hom}{Hom}
\DeclareMathOperator{\Aut}{Aut}
\DeclareMathOperator{\End}{End}
\DeclareMathOperator{\Hol}{Hol}
\DeclareMathOperator{\GL}{GL}
\DeclareMathOperator{\SL}{SL}
\DeclareMathOperator{\PGL}{PGL}
\DeclareMathOperator{\PSL}{PSL}
\DeclareMathOperator{\Spur}{Sp}
\DeclareMathOperator{\PSO}{PSO}
\DeclareMathOperator{\SO}{SO}
\DeclareMathOperator{\Or}{O}
\DeclareMathOperator{\Res}{Res}
\DeclareMathOperator{\diam}{diam}
\DeclareMathOperator{\sgn}{sgn}
\DeclareMathOperator{\sinc}{sinc}
\DeclareMathOperator{\disc}{disc}


\newcommand{\R}{\mathbb{R}}
\renewcommand{\P}{\mathbb{P}}
\newcommand{\C}{\mathbb{C}}
\newcommand{\N}{\mathbb{N}}
\newcommand{\Q}{\mathbb{Q}}
\newcommand{\Z}{\mathbb{Z}}
\newcommand{\F}{\mathbb{F}}
\newcommand{\g}{\mathfrak{g}}
\renewcommand{\o}{\mathfrak{o}}
\renewcommand{\d}{\mathfrak{d}}
\renewcommand{\a}{\mathfrak{a}}
\renewcommand{\b}{\mathfrak{b}}
\newcommand{\p}{\mathfrak{p}}
\newcommand{\I}{\mathbf{I}}

\renewcommand{\sl}{\mathfrak{sl}}
\newcommand{\h}{\mathfrak{h}}
\renewcommand{\H}{\mathbb{H}}
\textwidth 16cm
\oddsidemargin  0cm \evensidemargin 0cm \parindent0.7em
\textheight 23cm
\allowdisplaybreaks


\title[Time-frequency localization]{Eigenfunctions, Free Boundaries, and Time-Frequency Localization}
\author{Jo{\~a}o P.G. Ramos}



\begin{document}
\begin{abstract}
We study Gaussian time--frequency localization operators through their Bargmann--Fock representation and develop an inverse and free-boundary approach to their spectral theory.

We develop a perturbative inverse theory showing that functions sufficiently close to the ground-state Hermite function can be realized as eigenfunctions of localization operators. As a byproduct, we recover the classical disk characterization of Abreu and D\"orfler in the simply connected setting by the same free-boundary method. As an application we prove that the stability estimate of Gomez-Guerra, Ramos and Tilli is sharp.

In the second part of the paper we apply these inverse results to Faber--Krahn-type problems for Gaussian time--frequency concentration. We prove that local maximizers of the first localization eigenvalue under a measure constraint are necessarily disks, and we establish existence of global maximizers via a concentration--compactness argument in the Bargmann--Fock space, recovering the Nicola--Tilli theorem from this perspective.
\end{abstract}
\maketitle 

\section{Introduction}

The study of time--frequency concentration lies at the heart of modern harmonic analysis, with deep connections to time--frequency analysis, signal processing, mathematical physics, and phase--space methods; see for instance \cite{Grochenig,Mallat,Folland}. A central object in this theory is the \emph{short-time Fourier transform} (STFT), defined for a function $f \in L^2(\mathbb{R})$ and a window $\varphi \in L^2(\mathbb{R})$ by
\[
V_\varphi f(x,\omega) = \int_{\mathbb{R}} f(t)\,\varphi(x-t)\,e^{-2\pi i t \omega}\,dt.
\]
This transform measures the joint time--frequency content of $f$, and its squared modulus $|V_\varphi f|^2$ provides a phase--space energy density. Understanding how this energy can be concentrated in regions of the time--frequency plane is a fundamental problem, with both theoretical and applied significance.

In applications ranging from signal recovery to uncertainty principles, one is often interested in maximizing the quantity
\[
\int_\Omega |V_\varphi f(x,\omega)|^2\,dx\,d\omega
\]
over functions $f$ of unit norm and sets $\Omega$ of prescribed measure. This leads naturally to the study of \emph{time--frequency localization operators}, introduced by Daubechies in \cite{Daubechies} and further developed in \cite{Daubechies-Paul,Ramanathan-Topiwala-Weyl,Ramanathan-Topiwala-Spectrogram,Cordero-Grochenig,Boggiatto-Cordero-Grochenig,Bayer-Grochenig,Kisil-CrossToeplitz,Nicola-Tilli-2,Ramos-time-frequency}. For general windows, obtaining sharp concentration bounds is difficult, and Lieb's uncertainty principle \cite{Lieb} remains a basic benchmark; see also \cite{abreu2021donoho} for related large-sieve phenomena in modulation and polyanalytic Fock spaces. In the Gaussian setting, however, these operators enjoy a particularly rigid structure: via the Bargmann transform \cite{Bargmann}, they can be realized as integral operators on the Bargmann--Fock space, and the problem becomes one of complex analysis in weighted spaces; see \cite{Berezin,Berger-Coburn-Symbol,Berger-Coburn,Berger-Coburn-Heat,Zhu,Grochenig,Folland}. In particular, the STFT becomes intertwined with an entire function whose weighted modulus captures the concentration properties of the signal, allowing geometric questions to be translated into analytic ones.

A central theme in this area is the interplay between \emph{spectral properties} of localization operators and the \emph{geometry} of the underlying domains. In the Gaussian case, Nicola and Tilli proved the sharp Faber--Krahn inequality for the STFT, showing that disks maximize concentration among all sets of fixed measure \cite{Nicola-Tilli}; see also \cite{Nicola-Tilli-2,Bastianoni-Cordero-Nicola,Bastianoni-Teofanov} for complementary spectral and regularity results on localization operators. The stability theory for this variational problem was subsequently developed in \cite{Gomez-Guerra-Ramos-Tilli}, where quantitative control of both the first eigenfunction and the Fraenkel asymmetry is obtained from the spectral deficit. Related concentration questions for other transforms, including wavelet transforms, have been investigated in \cite{Ramos-Tilli}.

From a complementary perspective, one may ask inverse questions: to what extent does spectral information determine the geometry of the localization domain? Conversely, given a function, can one construct a domain for which it is an eigenfunction of the associated localization operator? The first systematic result in this direction is due to Abreu and D\"orfler \cite{Abreu-Doerfler}, who showed that if a Hermite function is an eigenfunction of a localization operator on a simply connected domain, then the domain must be a disk. This inverse viewpoint is closely related to weighted quadrature identities and free-boundary problems in potential theory; compare \cite{Aharonov-Schiffer-Zalcman,Cherednichenko,Gustafsson,Gustafsson-Shahgholian,Shahgholian,Karp}. It is also closely tied to the special role of Gaussian decay and Hermite structures in the Bargmann model; see \cite{Kulikov-Oliveira-Ramos,Goncalves-Oliveira-Steinerberger,Thangavelu}.

The guiding principle of the present paper is that, for Gaussian localization operators, eigenfunctions should be regarded as geometric data. More precisely, the eigenfunction should determine the domain in rigid situations, and the possibility of prescribing the eigenfunction should in turn have strong consequences for nearby variational problems. The three main threads of the paper are organized around this idea:
\begin{enumerate}
\item rigid eigenfunctions force rigid domains;
\item near the Gaussian, one can prescribe the eigenfunction and construct the corresponding localization domain;
\item once such an eigenfunction-to-domain mechanism is available, it can be used to analyze sharpness and extremizers in concentration inequalities.
\end{enumerate}

\medskip

The present work contributes to this line of research by developing a unified inverse and free-boundary framework for Gaussian time--frequency localization operators. Our approach builds on and extends the rigidity result of Abreu and D\"orfler \cite{Abreu-Doerfler}, the variational theory of Nicola and Tilli \cite{Nicola-Tilli,Nicola-Tilli-2}, and the quantitative stability results of G\'omez, Guerra, Ramos, and Tilli \cite{Gomez-Guerra-Ramos-Tilli}. At the technical level, our arguments also draw on ideas from inverse potential problems and quadrature-domain theory, in the spirit of \cite{Isakov,Cherednichenko,Gustafsson,Gustafsson-Shahgholian,Shahgholian}. In this way, the paper places localization operators squarely within the broader free-boundary and inverse-potential framework suggested by the Gaussian Bargmann model.

Despite these advances, several fundamental questions have remained open. In particular, the construction of localization domains associated with prescribed functions and the role of eigenfunction rigidity in variational concentration problems have not been fully understood. The results of the present paper aim to address these questions in a unified manner.

\medskip

Our first main result is the constructive half of this picture. It shows that, near the Gaussian, one can indeed prescribe the eigenfunction and realize it by a suitable localization domain.

\begin{theorem}\label{thm:perturb}
Let $\lambda\in (0,1)$ and $N\in\mathbb N$. There exists $\varepsilon_0(\lambda,N)>0$ such that if
\[
f_0 = h_0 + \sum_{k=1}^N a_k h_k
\]
satisfies $\max_k |a_k|<\varepsilon_0$, then there exists a domain $U_\lambda$ such that
\[
\mathcal L_{U_\lambda} f_0 = \lambda f_0 .
\]
\end{theorem}

This theorem provides, to the best of our knowledge, the first systematic construction of localization domains associated with nontrivial linear combinations of Hermite functions. In particular, it supplies the inverse counterpart to the rigidity phenomenon of Abreu and D\"orfler: rather than deducing the domain from a rigid eigenfunction, we construct domains from prescribed near-Gaussian eigenfunctions.

The proof proceeds in several stages, combining ideas from free-boundary problems and complex analysis. First, we reformulate the eigenvalue equation in the Bargmann--Fock space as a nonlinear integral identity involving the squared modulus of a polynomial. By applying a suitable differential operator, we reduce the problem to constructing a domain satisfying a weighted identity reminiscent of a quadrature condition.

We then construct an approximate solution by solving a free-boundary problem of obstacle type. This is achieved via a perturbative theorem of Isakov, which guarantees the existence of a domain whose associated potential satisfies a prescribed equation under small perturbations of the data. This step produces a domain for which the desired identity holds up to a controlled error.

In a second step, we refine this construction using a weighted quadrature-domain theorem. By viewing the problem through the lens of complex analysis, we interpret the identity as a condition on a weighted Cauchy transform. A perturbative implicit-function argument then allows us to deform the domain so that the identity holds exactly. Finally, we show that the resulting function satisfies the original eigenvalue equation by analyzing an associated differential operator and ruling out nontrivial solutions via growth estimates. The combination of these steps yields a robust local inverse theory.

\medskip

A noteworthy aspect of this construction is that it allows us to go beyond monomial-type eigenfunctions. Indeed, previous constructions in the literature were largely restricted to highly symmetric situations. Our result shows that a much broader class of functions can be realized spectrally, opening the door to a richer inverse theory and suggesting new connections with nonlinear potential theory.

\medskip

The same free-boundary framework also yields the rigid half of the correspondence, namely a short recovery of the classical theorem of Abreu and Dörfler.

More precisely, we provide two new proofs of the Abreu--Dörfler rigidity result based on our
methods: one through the local radiality mechanism for the logarithmic potential, and a second
through a holomorphic auxiliary function attached to the associated Poisson problem.

\begin{theorem}\label{thm:eig-loc-U}
Let $U$ be a simply connected bounded domain. If a Hermite function is an eigenfunction of the localization operator associated with $U$, then $U$ is a disk.
\end{theorem}

The proof proceeds by transporting the eigenvalue equation to the Bargmann--Fock space, where Hermite functions correspond to monomials and the localization operator becomes an integral operator with reproducing kernel structure. In this setting, the eigenvalue equation can be interpreted as a functional identity involving exponential generating functions. By differentiating appropriately and invoking a Paley--Wiener type argument, we reduce the problem to the vanishing of a Fourier transform along complex lines.

This allows us to apply a divisibility lemma of Ehrenpreis--Malgrange type, which yields the existence of a logarithmic potential solving a Poisson equation with source supported on the domain. The key insight is that the spectral assumption imposes strong constraints on this potential: it must exhibit a form of radial behavior on the boundary. In the simply connected setting, this forces the domain to be a disk, thereby recovering the Abreu--Dörfler rigidity theorem from the same inverse/free-boundary perspective.

\medskip

We now turn to consequences of this inverse point of view. Our first application concerns the quantitative stability theorem of G\'omez, Guerra, Ramos, and Tilli for the Faber--Krahn inequality for the Gaussian STFT. In the formulation most relevant here, if \(U\subset \R^2\) is a measurable set of finite positive measure, \(\lambda_1(U)\) denotes the first eigenvalue of the associated localization operator, and \(f_U\) is a corresponding \(L^2\)-normalized first eigenfunction, then Corollary~1.4 of \cite{Gomez-Guerra-Ramos-Tilli} gives the bounds
\[
\min_{z_0\in\C,\ |c|=1}\|f_U-c\phi_{z_0}\|_{L^2(\R)}
\leq
C e^{|U|/2}
\left(
1-\frac{\lambda_1(U)}{1-e^{-|U|}}
\right)^{1/2},
\]
and
\[
\mathcal A(U)
\leq
K(|U|)
\left(
1-\frac{\lambda_1(U)}{1-e^{-|U|}}
\right)^{1/2},
\]
where \(\phi_{z_0}\) ranges over the Gaussian optimizers and \(\mathcal A(U)\) denotes the Fraenkel asymmetry of \(U\).

Our perturbative inverse theorem provides the missing input needed to test this stability estimate on genuine spectral data. Starting from a prescribed perturbation of the Gaussian, we construct an exact eigenpair \((U_\varepsilon,f_\varepsilon)\) for the localization operator, with \(U_\varepsilon\) a smooth perturbation of the disk and \(f_\varepsilon\) a corresponding first eigenfunction. Thus the constructive direction of the paper produces explicit near-extremizers, and allows us to compare the size of the spectral deficit with the actual geometric displacement of the domain from the class of disks. The resulting family shows that the square-root dependence in Corollary~1.4 is the best possible one.

\begin{theorem} \label{thm:GGRT-sharp}
The exponent \(1/2\) in Corollary~1.4 of \cite{Gomez-Guerra-Ramos-Tilli} is optimal. More precisely, the powers of the deficit
\[
1-\frac{\lambda_1(U)}{1-e^{-|U|}}
\]
appearing in both the eigenfunction estimate and the Fraenkel-asymmetry estimate there cannot be replaced by \(\beta\) for any \(\beta>1/2\).
\end{theorem}

The argument is based on constructing a carefully chosen family of perturbations of the disk using the inverse theorem above. Starting from a polynomial perturbation of the ground-state Hermite function, we obtain a corresponding family of domains whose geometry deviates from that of a disk at first order. We then perform a detailed asymptotic analysis of the eigenvalue equation, computing its first variation via a transport formula and identifying the leading Fourier modes of the boundary deformation.

This analysis shows that the deformation is governed by second-order harmonics, which cannot be eliminated by translations. By comparing the perturbed domains with arbitrary disks, we obtain lower bounds on their symmetric difference that match the upper bounds in the stability estimate. The proof combines spectral perturbation theory, geometric measure estimates, and harmonic analysis on the circle, and illustrates how inverse constructions can be used to probe the sharpness of variational inequalities.

\medskip

We next connect this eigenfunction-to-domain philosophy to variational problems. Our fourth main result concerns the structure of local maximizers for Gaussian time--frequency concentration.

\begin{theorem}
Any local maximizer of $\lambda_{1,\varphi}$ under a measure constraint is, up to translation, a disk.
\end{theorem}

The proof relies on combining variational arguments with the inverse spectral perspective developed earlier. We first show that any local maximizer must correspond to a superlevel set of the density associated with its first eigenfunction in the Bargmann--Fock space. This reduces the geometric problem to a functional one.

We then analyze the spectral properties of the localization operator at a maximizer, showing that the first eigenvalue is simple and that derivatives of the eigenfunction remain eigenfunctions. This strong rigidity property forces the maximizing eigenfunction into a Gaussian form. Once this is known, the corresponding domain is determined as well, and one concludes that the only admissible configuration is a disk. The argument thus shows that extremizers are governed by the same rigidity mechanism that underlies the inverse theory.

\medskip

Finally, we address the global variational problem and recover the Nicola--Tilli theorem from the same perspective.

\begin{theorem}[New proof of Nicola--Tilli]
Disks maximize Gaussian time--frequency concentration among all sets of fixed measure.
\end{theorem}

Our proof is based on a concentration--compactness argument in the Bargmann--Fock space. We consider maximizing sequences and analyze their possible limiting behavior, ruling out vanishing and dichotomy scenarios using the structure of the Fock space and the invariance properties of the STFT. This yields the existence of a maximizer. Related concentration--compactness methods have recently been used in nearby time--frequency concentration problems, notably by Nicola, Romero, and Trapasso for ambiguity-function type concentration functionals \cite{Nicola-Romero-Trapasso}; see also \cite{Marceca-Romero-Speckbacher,Samuelsen} for recent work on Fourier concentration operators, and \cite{Fulsche-Hagger} for closely related developments in the polyanalytic Fock-space setting.

We then combine this existence result with the rigidity of local maximizers established earlier. Since any global maximizer is, in particular, a local maximizer, its eigenfunction is rigid and therefore determines the maximizing domain. This provides a new conceptual route to the Nicola--Tilli theorem, avoiding rearrangement inequalities and instead relying on inverse and free-boundary methods. The argument highlights the power of treating eigenfunctions as geometric data.

The paper is organized as follows. In Section~2 we review the Bargmann--Fock representation and establish several auxiliary results, including a key divisibility lemma and perturbative free-boundary theorems. Section~3 is devoted to the perturbative inverse construction. Section~4 explains how the same machinery recovers the classical theorem of Abreu and Dörfler in the simply connected case. Section~5 contains the sharpness result, while the remaining sections are devoted to variational applications, culminating in the new proof of the Nicola--Tilli theorem.


\section{Preliminaries}

\subsection{Properties of the Bargmann-Fock space} We briefly recall the connection between the Gaussian short-time Fourier transform
(STFT) and the Bargmann--Fock space, which will be used systematically in what follows.

\medskip

Let $\phi(x) = 2^{1/4} e^{-\pi x^2}$ be the normalized Gaussian. For $f \in L^2(\mathbb{R})$, its Gaussian STFT is
\begin{equation}\label{eq:STFT}
V_\phi f(x,\omega)
= \int_{\mathbb{R}} f(t)\, \phi(x-t)\, e^{-2\pi i t \omega}\, dt.
\end{equation}
The Bargmann transform
$\mathcal{B} : L^2(\mathbb{R}) \to \mathcal{F}^2(\mathbb{C})$ is defined by
\begin{equation}\label{eq:Bargmann}
(\mathcal{B}f)(z)
= 2^{1/4} \int_{\mathbb{R}} f(t)\, e^{2\pi t z - \pi t^2 - \frac{\pi}{2} z^2}\, dt,
\qquad z \in \mathbb{C}.
\end{equation}
It is well known that $\mathcal{B}$ is a unitary isomorphism onto the Bargmann--Fock
space $\mathcal{F}^2(\mathbb{C})$, endowed with the norm
\[
\|F\|_{\mathcal{F}^2}^2
= \int_{\mathbb{C}} |F(z)|^2 e^{-\pi |z|^2} \, dA(z).
\]
A direct computation shows that, for $z = x + i\omega$,
\begin{equation}\label{eq:STFT-B}
V_\phi f(x,-\omega)
= e^{\pi i x \omega}\, (\mathcal{B}f)(z)\, e^{-\frac{\pi}{2}|z|^2}.
\end{equation}
In particular, setting $F = \mathcal{B}f$, one obtains
\begin{equation}\label{eq:energy}
\int_{\mathbb{R}^2} |V_\phi f(x,\omega)|^2 \, dx\, d\omega
= \int_{\mathbb{C}} |F(z)|^2 e^{-\pi |z|^2} \, dA(z)
= \|F\|_{\mathcal{F}^2}^2.
\end{equation}

\medskip

\noindent
\subsubsection*{Weyl operators and invariance.}
For $a \in \mathbb{C}$, define the operator
\begin{equation}\label{eq:Ta}
(T_a F)(z) := e^{\pi a z - \frac{\pi}{2}|a|^2} F(z - a).
\end{equation}
Then $T_a$ is unitary on $\mathcal{F}^2(\mathbb{C})$. Indeed,
\begin{align*}
\|T_aF\|_{\mathcal F^2}^2
&=
\int_{\C} e^{2\pi \Re(az)-\pi |a|^2}\,|F(z-a)|^2 e^{-\pi |z|^2}\,dA(z) \\
&=
\int_{\C} |F(z-a)|^2 e^{-\pi |z-a|^2}\,dA(z)
=
\|F\|_{\mathcal F^2}^2,
\end{align*}
because
\[
2\Re(az)-|a|^2-|z|^2=-|z-a|^2.
\]
Moreover,
\begin{equation}\label{eq:u-translation-global}
u_{T_a F}(z) = u_F(z - a).
\end{equation}
Indeed,
\[
u_{T_aF}(z)
=
|T_aF(z)|^2e^{-\pi |z|^2}
=
e^{2\pi\Re(az)-\pi|a|^2}|F(z-a)|^2e^{-\pi |z|^2}
=
|F(z-a)|^2e^{-\pi |z-a|^2}.
\]
In particular, the functionals $I_F$ defined by, 
for \(s>0\),
\[
I_F(s):=\sup_{|\Omega|=s}\int_\Omega u_F\,dA,
\]
satisfy
\begin{equation}\label{eq:I-translation-global}
I_{T_a F}(s) = I_F(s), \qquad a \in \mathbb{C},
\end{equation}
reflecting the invariance of the Gaussian STFT under time--frequency shifts.

\medskip

\noindent
\subsubsection*{Connection with localization.}
If $F = \mathcal{B}f$, then for every measurable $\Omega \subset \mathbb{C}$,
\[
\int_\Omega u_F \, dA
= \int_{\widetilde{\Omega}} |V_\phi f(x,\omega)|^2 \, dx\, d\omega,
\]
where $\widetilde{\Omega} = \{(x,\omega) : (x,-\omega) \in \Omega\}$.
Thus, concentration properties of the STFT are naturally encoded by the function
\begin{equation}\label{eq:uF}
u_F(z) := |F(z)|^2 e^{-\pi |z|^2}, \qquad z \in \mathbb{C},
\end{equation}
which is the main object of interest in the Bargmann--Fock setting, in terms of the
quantities $I_F(s)$.

\medskip

We next collect a few elementary facts that will be used repeatedly.

\begin{lemma}\label{lem:Fock-basic}
Let $F,G \in \mathcal{F}^2(\mathbb{C})$. Then:
\begin{enumerate}
\item[(i)] The space $\mathcal{F}^2(\mathbb{C})$ is a reproducing kernel Hilbert space with kernel
\[
K_w(z) = e^{\pi z \overline{w}}, \qquad z,w \in \mathbb{C}.
\]

\item[(ii)] (Sharp pointwise bound) For every $z \in \mathbb{C}$,
\begin{equation}\label{eq:pointwise}
u_F(z) \le \|F\|_{\mathcal{F}^2}^2,
\end{equation}
and equality holds at some point if and only if $F$ is proportional to $F_w$
for some $w \in \mathbb{C}$.

\item[(iii)] $u_F(z) \to 0$ as $|z| \to \infty$.

\item[(iv)] One has the estimate
\[
\|u_F - u_G\|_{L^1(\mathbb{C})}
\le \big( \|F\|_{\mathcal{F}^2} + \|G\|_{\mathcal{F}^2} \big)
\|F - G\|_{\mathcal{F}^2}.
\]
\end{enumerate}
\end{lemma}

\begin{proof}
(i) For \(w\in \C\), the function \(K_w(z)=e^{\pi z\overline w}\) belongs to
\(\mathcal F^2(\C)\), and
\[
\|K_w\|_{\mathcal F^2}^2
=
\int_{\C} e^{2\pi \Re(z\overline w)-\pi |z|^2}\,dA(z)
=
e^{\pi |w|^2},
\]
after completing the square. Thus point evaluation is continuous, and for every
\(F\in \mathcal F^2(\C)\) one has the reproducing identity
\[
F(w)=\langle F,K_w\rangle_{\mathcal F^2}.
\]
In particular \(F\) is entire.

(ii) is a consequence of the reproducing kernel identity and Cauchy--Schwarz,
as in \eqref{eq:pointwise}.

(iii) Fix \(\varepsilon>0\), and choose a polynomial \(P\) such that
\[
\|F-P\|_{\mathcal F^2}<\varepsilon
\]
using the density of polynomials in \(\mathcal F^2(\C)\). By the pointwise bound from (ii),
\[
|F(z)-P(z)|e^{-\frac{\pi}{2}|z|^2}\le \varepsilon
\qquad\text{for every }z\in \C.
\]
Since \(P\) is a polynomial,
\[
P(z)e^{-\frac{\pi}{2}|z|^2}\to 0
\qquad\text{as }|z|\to\infty.
\]
Hence
\[
\limsup_{|z|\to\infty}|F(z)|e^{-\frac{\pi}{2}|z|^2}\le \varepsilon.
\]
As \(\varepsilon>0\) is arbitrary, we conclude that
\[
|F(z)|e^{-\frac{\pi}{2}|z|^2}\to 0,
\]
that is, \(u_F(z)\to 0\) as \(|z|\to\infty\).

(iv) Writing
\[
|u_F - u_G|
= \big||F|^2 - |G|^2\big| e^{-\pi|z|^2}
\le |F-G|(|F|+|G|) e^{-\pi|z|^2},
\]
the claim follows by Cauchy--Schwarz.
\end{proof}

\medskip

We now note the following crucial fact about the Bargmann transform, which is that it transforms the basis of normalized Hermite functions onto normalized monomials, yielding thus a natural orthonormal basis on $\mathcal{F}^2(\C).$

\begin{lemma}\label{lem:monomials}
The functions
\[
e_k(z) := b_k z^k, \qquad k \ge 0,
\]
form an orthonormal basis of $\mathcal{F}^2(\mathbb{C})$, where
\[
b_k = \left( \frac{\pi^k}{k!} \right)^{1/2}.
\]
In particular,
\[
\langle e_k, e_m \rangle_{\mathcal{F}^2} = \delta_{k,m}.
\]
\end{lemma}

\begin{proof}
A direct computation in polar coordinates gives
\[
\int_{\mathbb{C}} |z|^{2k} e^{-\pi |z|^2} dA(z)
= 2\pi \int_0^\infty r^{2k+1} e^{-\pi r^2} dr
= \frac{k!}{\pi^k}.
\]
Thus $\|z^k\|_{\mathcal{F}^2}^2 = \frac{k!}{\pi^k}$, which yields the normalization.
Orthogonality follows from rotational invariance.
Completeness is standard.
\end{proof}

\medskip

\begin{lemma}\label{lem:growth}
Let $F \in \mathcal{F}^2(\mathbb{C})$. Then:
\begin{enumerate}
\item[(i)] $F$ admits the expansion
\[
F(z) = \sum_{k=0}^\infty a_k e_k(z),
\qquad \sum_{k=0}^\infty |a_k|^2 = \|F\|_{\mathcal{F}^2}^2.
\]

\item[(ii)] (Growth bound) For every $z \in \mathbb{C}$,
\[
|F(z)| \le \|F\|_{\mathcal{F}^2} e^{\frac{\pi}{2}|z|^2}.
\]

\item[(iii)] (Derivative bounds) For every $z \in \mathbb{C}$,
\[
|F'(z)| \le C(1+|z|)\, \|F\|_{\mathcal{F}^2} e^{\frac{\pi}{2}|z|^2}.
\]
\end{enumerate}
\end{lemma}

\begin{proof}
(i) follows from the orthonormal basis in Lemma \ref{lem:monomials}.

(ii) follows from the reproducing kernel estimate:
\[
|F(z)| = |\langle F, K_z \rangle|
\le \|F\|_{\mathcal{F}^2} \|K_z\|_{\mathcal{F}^2}
= \|F\|_{\mathcal{F}^2} e^{\frac{\pi}{2}|z|^2}.
\]

(iii) Since
\[
F'(z)=\pi \langle F,\overline{w}K_z(w)\rangle_{\mathcal F^2},
\]
Cauchy--Schwarz gives
\[
|F'(z)|
\le
\pi \|F\|_{\mathcal F^2}\,\|\overline{w}K_z\|_{\mathcal F^2}.
\]
Now
\[
\|\overline{w}K_z\|_{\mathcal F^2}^2
=
\int_{\C} |w|^2 e^{2\pi\Re(w\overline z)-\pi |w|^2}\,dA(w)
=
e^{\pi |z|^2}\int_{\C} |w+z|^2 e^{-\pi |w|^2}\,dA(w),
\]
after the change of variables \(w\mapsto w+z\). The last integral is bounded by
\(C(1+|z|^2)\), so
\[
|F'(z)| \le C(1+|z|)\,\|F\|_{\mathcal F^2}e^{\frac{\pi}{2}|z|^2}.
\]
\end{proof}

\medskip

\begin{lemma}\label{lem:max}
Let $F \in \mathcal{F}^2(\mathbb{C})$ and set
\[
T := \sup_{z \in \mathbb{C}} u_F(z).
\]
Then $0 \le T \le \|F\|_{\mathcal{F}^2}^2$, and $T = \|F\|_{\mathcal{F}^2}^2$
if and only if $F$ is proportional to $F_w$ for some $w \in \mathbb{C}$.
\end{lemma}

\begin{proof}
The bound follows from Lemma \ref{lem:Fock-basic}(ii). The characterization of
equality is the same as in the reproducing kernel inequality.
    \end{proof}

\medskip

    We proceed with a few elementary facts.

\begin{lemma}\label{lem:basic-properties-global}
Let \(F,G\in \mathcal F^2(\C)\). Then:
\begin{enumerate}
    \item[(i)] \(I_F(s)\) is attained by a superlevel set of \(u_F\), and
    \[
    0\le I_F(s)\le \|F\|_{\mathcal F^2}^2.
    \]
    \item[(ii)] For every \(c\in \C\),
    \[
    I_{cF}(s)=|c|^2 I_F(s).
    \]
    \item[(iii)] One has the Lipschitz bound
    \[
    |I_F(s)-I_G(s)|
    \le
    \|u_F-u_G\|_{L^1(\C)}
    \le
    \big(\|F\|_{\mathcal F^2}+\|G\|_{\mathcal F^2}\big)\,
    \|F-G\|_{\mathcal F^2}.
    \]
\end{enumerate}
\end{lemma}

\begin{proof}
The identity
\[
\int_\C u_F\,dA=\int_\C |F(z)|^2e^{-\pi |z|^2}\,dA(z)=\|F\|_{\mathcal F^2}^2
\]
is immediate from the definition of the Fock norm. Since \(F\in\mathcal F^2\),
the standard reproducing-kernel estimate implies that \(u_F\) is bounded and
continuous, and in fact \(u_F(z)\to 0\) as \(|z|\to \infty\); thus
\(u_F\in C_0(\C)\cap L^1(\C)\).

The existence of a maximizing set for \(I_F(s)\), given by a superlevel set of
\(u_F\), is the classical bathtub principle. The bound
\[
0\le I_F(s)\le \int_\C u_F\,dA=\|F\|_{\mathcal F^2}^2
\]
is immediate. The scaling identity in (ii) is obvious.

For (iii), for every measurable \(\Omega\subset \C\) with \(|\Omega|=s\),
\[
\left|\int_\Omega u_F\,dA-\int_\Omega u_G\,dA\right|
\le \int_\C |u_F-u_G|\,dA.
\]
Taking suprema over \(|\Omega|=s\) gives
\[
|I_F(s)-I_G(s)|\le \|u_F-u_G\|_{L^1(\C)}.
\]
Finally,
\begin{align*}
\|u_F-u_G\|_{L^1(\C)}
&=
\int_\C \big||F|^2-|G|^2\big|\,e^{-\pi |z|^2}\,dA \\
&\le
\int_\C |F-G|(|F|+|G|)e^{-\pi |z|^2}\,dA \\
&\le
\|F-G\|_{\mathcal F^2}\big(\|F\|_{\mathcal F^2}+\|G\|_{\mathcal F^2}\big)
\end{align*}
by Cauchy--Schwarz.
\end{proof}


\subsection{From time-frequency analysis to free boundary problems} We now use the following divisibility lemma of Ehrenpreis--Malgrange flavour.

In order to state it, we say that a compactly supported distribution $\omega$ is a \emph{singular-continuous free measure} if it defines a Radon measure on $\R^2$ whose Lebesgue decomposition has only absolutely continuous and pure point parts, and whose pure point part consists of finitely many atoms.

\begin{lemma}\label{lemma:fourier-to-fboundary}
Let $\omega\in \mathcal{S}'(\R^2)$ be a compactly supported singular-continuous free distribution such that
\[
\widehat{\omega}(w,\pm iw)=0,\qquad \forall w\in\C.
\]
Then there exist a finite set $P\subset \R^2$ and a compactly supported function $u\in L^2(\R^2)$ such that
\[
\Delta u=\omega
\]
in the sense of distributions, and
\[
u\in C(\R^2\setminus P).
\]
\end{lemma}

\begin{proof}
By the Paley--Wiener theorem, $\widehat{\omega}(\xi,\eta)$ is an entire function on $\C^2$ of exponential type. Hence it admits a power series expansion
\[
\widehat{\omega}(\xi,\eta)=\sum_{k,l\ge 0} a_{k,l}\,\xi^k\eta^l.
\]
Using the vanishing hypotheses on the complex lines $\eta=\pm i\xi$, we get
\[
0=\widehat{\omega}(w,\pm iw)
   =\sum_{n\ge 0}\left(\sum_{k=0}^n a_{k,n-k}(\pm i)^{\,n-k}\right)w^n,
\]
and therefore, by comparing coefficients,
\[
\sum_{k=0}^n a_{k,n-k}(\pm i)^{\,n-k}=0,\qquad \forall n\ge 0.
\]
For each $n\ge 0$, define
\[
p_n(t):=\sum_{k=0}^n a_{k,n-k} t^k.
\]
Then $p_n(\pm i)=0$, so $p_n(t)$ is divisible by $t^2+1$. Thus we may write
\[
p_n(t)=(t^2+1)q_n(t),
\]
where $q_n$ is a polynomial of degree at most $n-2$. Reassembling the homogeneous pieces, it follows that
\[
\widehat{\omega}(\xi,\eta)=(\xi^2+\eta^2)F(\xi,\eta)
\]
for some entire function $F$ on $\C^2$.

To justify the reassembly near the origin, write
\[
\widehat\omega(\xi,\eta)=\sum_{n\ge 0} H_n(\xi,\eta),
\]
where \(H_n\) is the homogeneous polynomial of degree \(n\),
\[
H_n(\xi,\eta):=\sum_{k=0}^n a_{k,n-k}\,\xi^k\eta^{\,n-k}.
\]
The argument above shows that each \(H_n\) vanishes on the two lines
\(\eta=\pm i\xi\), hence is divisible by \(\xi^2+\eta^2\). Thus
\[
H_n(\xi,\eta)=(\xi^2+\eta^2)G_{n-2}(\xi,\eta)
\]
for some homogeneous polynomial \(G_{n-2}\) of degree \(n-2\) (with
\(G_{n-2}\equiv 0\) for \(n<2\)). Summing these homogeneous pieces gives an
entire function
\[
F(\xi,\eta):=\sum_{n\ge 2} G_{n-2}(\xi,\eta)
\]
and the claimed factorization follows term by term.

Again by Paley--Wiener, there exist constants $C,c>0$ such that
\[
|\widehat{\omega}(\xi,\eta)|\le C e^{c(|\Im \xi|+|\Im \eta|)}.
\]
Moreover, since \(\widehat\omega\) has a zero of order at least two at the
origin, the quotient \(F\) is entire there, and away from a fixed neighborhood
of the origin we may divide by \(\xi^2+\eta^2\). Therefore
\[
|F(\xi,\eta)|\le C'\frac{e^{c(|\Im \xi|+|\Im \eta|)}}{1+|\xi|^2+|\eta|^2},
\]
after enlarging \(C'\) if necessary to absorb the bounded region near the
origin. In particular, restricting to $(\xi,\eta)\in \R^2$, we obtain
\[
|F(\xi,\eta)|\le \frac{C''}{1+|\xi|^2+|\eta|^2},
\]
and therefore $F|_{\R^2}\in L^2(\R^2)$.

Let $u\in L^2(\R^2)$ be such that $\widehat{u}=F$ on $\R^2$. Then, up to the harmless Fourier-transform normalisation constant,
\[
\widehat{\Delta u}(\xi,\eta)=-(\xi^2+\eta^2)\widehat{u}(\xi,\eta)
= -(\xi^2+\eta^2)F(\xi,\eta),
\]
so after rescaling $u$ by a nonzero real constant if necessary, we may assume that
\[
\Delta u=\omega
\]
in the sense of distributions.

Since $\omega$ is a compactly supported Radon measure, $u$ agrees (up to an additive harmonic function, which must vanish by compact support) with the logarithmic potential of $\omega$:
\[
u(x,y)=c\, G*\omega(x,y),
\]
where $G(z)=\log|z|$ is the Green kernel in $\R^2$ and $c\in\R$. The absolutely continuous part of $\omega$ produces a continuous logarithmic potential, and the only possible singularities come from the finitely many atoms in the pure point part. Hence
\[
u\in C(\R^2\setminus P),
\]
where $P$ is the finite support of the atomic part of $\omega$.
\end{proof}

\subsection{Two perturbative auxiliary results}\label{app:auxiliary-results}

In this appendix we prove two perturbative results used in the proof of
Theorem 4. The first is a free-boundary stability
statement around a strictly negative obstacle-type solution. The second is a
local realization theorem for weighted quadrature domains near a reference
domain. Both statements are tailored to the setting of the present paper.

\subsubsection{A perturbative free-boundary theorem}\label{app:isakov-type}

We begin with a perturbative result for the one-phase problem
\[
-\Delta u = f\,\mathbf{1}_{\Omega},
\qquad
u=0 \ \text{on } \C\setminus \Omega,
\]
near a reference pair \((\Omega_0,u_0)\).

Let \(\Omega_0\subset \C\simeq \R^2\) be a bounded domain of class \(C^2\). For
\(\rho>0\) sufficiently small, every \(C^2\) domain \(\Omega\) which is
\(\rho\)-close to \(\Omega_0\) may be written uniquely as a normal graph over
\(\partial\Omega_0\):
\[
\partial\Omega
=
\{x+\varphi(x)\nu(x):x\in \partial\Omega_0\},
\]
where \(\nu\) denotes the outer unit normal to \(\partial\Omega_0\), and
\(\varphi\in C^2(\partial\Omega_0)\) has small norm. We shall write
\(\Omega=\Omega_\varphi\).

\begin{proposition}[\cite{Isakov}] \label{prop:isakov-tailored}
Let \(\Omega_0\subset \C\) be a bounded \(C^2\) domain, and let
\(f_0\in C^5(\overline{\Omega'})\) for some open neighborhood
\(\Omega'\supset \overline{\Omega_0}\), with
\[
f_0>0 \qquad \text{on } \partial\Omega_0.
\]
Assume there exists \(u_0\in C^{2,\alpha}(\overline{\Omega_0})\) for some
\(\alpha\in(0,1)\) such that
\begin{equation}\label{eq:appendix-reference-problem}
\left\{
\begin{aligned}
-\Delta u_0 &= f_0 &&\text{in }\Omega_0,\\
u_0&=0 &&\text{on }\partial\Omega_0,\\
u_0&<0 &&\text{in }\Omega_0.
\end{aligned}
\right.
\end{equation}
Then there exist \(\varepsilon>0\) and \(\rho>0\) such that the following holds.

Whenever \(f\in C^5(\overline{\Omega'})\) satisfies
\[
\|f-f_0\|_{C^5(\overline{\Omega'})}<\varepsilon,
\]
there exists a unique \(C^2\) function
\(\varphi\in C^2(\partial\Omega_0)\), with \(\|\varphi\|_{C^2}<\rho\), such that
for \(\Omega_f:=\Omega_\varphi\) the Dirichlet problem
\begin{equation}\label{eq:appendix-perturbed-problem}
\left\{
\begin{aligned}
-\Delta u_f &= f &&\text{in }\Omega_f,\\
u_f&=0 &&\text{on }\partial\Omega_f
\end{aligned}
\right.
\end{equation}
has a solution \(u_f\in C^{2,\alpha}(\overline{\Omega_f})\) satisfying
\[
u_f<0 \qquad \text{in }\Omega_f.
\]
Moreover, the map
\[
f\mapsto \varphi(f)
\]
is \(C^1\) from a neighborhood of \(f_0\) in \(C^5(\overline{\Omega'})\) into
\(C^2(\partial\Omega_0)\); in particular, \(\varphi\) depends continuously on
\(f\) in the indicated norms.
\end{proposition}

\begin{proof}
We divide the argument into several steps.

\medskip
\noindent\emph{Step 1: nondegeneracy of the reference solution.}
Since \(u_0<0\) in \(\Omega_0\), \(u_0=0\) on \(\partial\Omega_0\), and
\(-\Delta u_0=f_0\ge c_0>0\) in a neighborhood of \(\partial\Omega_0\), the Hopf
boundary lemma yields
\[
\partial_\nu u_0(x)>0
\qquad \text{for every }x\in \partial\Omega_0,
\]
where \(\nu\) is the outer unit normal. Since \(\partial\Omega_0\) is compact,
there exists \(\kappa_0>0\) such that
\begin{equation}\label{eq:appendix-hopf-lower-bound}
\partial_\nu u_0\ge \kappa_0
\qquad\text{on }\partial\Omega_0.
\end{equation}

\medskip
\noindent\emph{Step 2: reduction to a boundary equation.}
Fix a tubular neighborhood
\[
\mathcal N:=\{x+t\nu(x):x\in \partial\Omega_0,\ |t|<\delta\}
\]
on which the normal coordinates are well-defined. For \(\varphi\) sufficiently
small in \(C^2\), the perturbed boundary
\[
\Gamma_\varphi:=\{x+\varphi(x)\nu(x):x\in \partial\Omega_0\}
\]
bounds a \(C^2\) domain \(\Omega_\varphi\subset \Omega'\).

For each pair \((\varphi,f)\) with \(\varphi\) small and \(f\) close to \(f_0\),
let \(u_{\varphi,f}\) denote the unique solution of
\[
\left\{
\begin{aligned}
-\Delta u_{\varphi,f} &= f &&\text{in }\Omega_\varphi,\\
u_{\varphi,f}&=0 &&\text{on }\Gamma_\varphi.
\end{aligned}
\right.
\]
By pulling back the equation from \(\Omega_\varphi\) to \(\Omega_0\) through the
normal-graph diffeomorphism, standard Schauder theory implies that the map
\[
(\varphi,f)\mapsto u_{\varphi,f}\circ \Phi_\varphi
\]
is \(C^1\) from a neighborhood of \((0,f_0)\) in
\(C^2(\partial\Omega_0)\times C^5(\overline{\Omega'})\) to
\(C^{2,\alpha}(\overline{\Omega_0})\). Here \(\Phi_\varphi:\Omega_0\to \Omega_\varphi\)
denotes the graph diffeomorphism.

Now define the nonlinear boundary operator
\[
\mathcal F(\varphi,f)(x)
:=
u_{\varphi,f}(x+\varphi(x)\nu(x)),
\qquad x\in \partial\Omega_0.
\]
Since \(u_{\varphi,f}\) vanishes on \(\Gamma_\varphi\), the condition that
\(\Gamma_\varphi\) be the free boundary is precisely
\begin{equation}\label{eq:appendix-free-boundary-equation}
\mathcal F(\varphi,f)=0.
\end{equation}
By the regularity statement above, \(\mathcal F\) is \(C^1\) from a neighborhood
of \((0,f_0)\) into \(C^{2,\alpha}(\partial\Omega_0)\).

\medskip
\noindent\emph{Step 3: computation of the linearization in \(\varphi\).}
We compute the derivative of \(\mathcal F\) with respect to \(\varphi\) at
\((0,f_0)\). Let \(\psi\in C^2(\partial\Omega_0)\), and consider
\(\varphi_t=t\psi\). Since \(u_0=0\) on \(\partial\Omega_0\), the shape derivative
of the boundary trace is
\[
D_\varphi \mathcal F(0,f_0)[\psi]
=
\psi\,\partial_\nu u_0
\qquad \text{on }\partial\Omega_0.
\]
This is the standard Hadamard formula in the present Dirichlet setting: moving
the boundary in the normal direction by \(\psi\) changes the boundary value at
first order by \(\psi\,\partial_\nu u_0\).

In view of \eqref{eq:appendix-hopf-lower-bound}, the multiplication operator
\[
\psi\mapsto \psi\,\partial_\nu u_0
\]
is an isomorphism of \(C^2(\partial\Omega_0)\) onto itself, and likewise of
\(C^{2,\alpha}(\partial\Omega_0)\) onto itself.

\medskip
\noindent\emph{Step 4: application of the implicit function theorem.}
Since \(\mathcal F(0,f_0)=0\) and \(D_\varphi\mathcal F(0,f_0)\) is invertible, the
implicit function theorem yields \(\varepsilon>0\), \(\rho>0\), and a unique \(C^1\)
map
\[
f\mapsto \varphi(f)\in C^2(\partial\Omega_0),
\qquad
\|\varphi(f)\|_{C^2}<\rho,
\]
defined for \(\|f-f_0\|_{C^5}<\varepsilon\), such that
\[
\mathcal F(\varphi(f),f)=0.
\]
Equivalently, the corresponding domain \(\Omega_f:=\Omega_{\varphi(f)}\) satisfies
that the solution \(u_f:=u_{\varphi(f),f}\) of \eqref{eq:appendix-perturbed-problem}
vanishes on \(\partial\Omega_f\).

\medskip
\noindent\emph{Step 5: preservation of the sign.}
Since \(u_0<0\) in \(\Omega_0\), the continuity of the map
\((\varphi,f)\mapsto u_{\varphi,f}\) in \(C^0\) implies that, for \(f\) close
enough to \(f_0\), the solution \(u_f\) is still strictly negative in
\(\Omega_f\). Indeed, after pulling back to \(\Omega_0\), one has
\[
\|u_f\circ \Phi_{\varphi(f)}-u_0\|_{C^0(\overline{\Omega_0})}\to 0
\qquad \text{as } f\to f_0,
\]
but since \(u_0=0\) on \(\partial\Omega_0\), one needs a slight refinement near
the boundary. Fix \(\delta_0>0\) so small that the collar
\[
\mathcal N_{\delta_0}:=\{x-t\nu(x):x\in \partial\Omega_0,\ 0<t<\delta_0\}
\]
is contained in \(\Omega_0\). By \eqref{eq:appendix-hopf-lower-bound} and the
\(C^{1}\)-regularity of \(u_0\) up to the boundary, shrinking \(\delta_0\) if
necessary we obtain
\[
u_0(x-t\nu(x))\le -\frac{\kappa_0}{2}t
\qquad
\text{for }x\in \partial\Omega_0,\ 0<t<\delta_0.
\]
On the compact set
\[
K_{\delta_0}:=\{y\in \Omega_0:\operatorname{dist}(y,\partial\Omega_0)\ge \delta_0/2\}
\]
one has \(u_0\le -m_0\) for some \(m_0>0\). Now, for \(f\) sufficiently close
to \(f_0\), the graph map \(\Phi_{\varphi(f)}\) is \(C^1\)-close to the
identity and \(u_f\circ \Phi_{\varphi(f)}\) is \(C^1\)-close to \(u_0\) on
\(\overline{\Omega_0}\). It follows that \(u_f\circ \Phi_{\varphi(f)}<0\) on
\(K_{\delta_0}\), while on the collar \(\mathcal N_{\delta_0}\) its outward
normal derivative remains bounded below by \(\kappa_0/2\). Since
\(u_f=0\) on \(\partial\Omega_f\), integrating this derivative estimate along
the inward normal segments shows that \(u_f\circ \Phi_{\varphi(f)}<0\) on the
whole collar \(\mathcal N_{\delta_0}\) as well. Hence \(u_f<0\) throughout
\(\Omega_f\).

This proves existence. Uniqueness among sufficiently small \(C^2\) normal graphs
follows from the uniqueness statement in the implicit function theorem.
\end{proof}

\subsubsection{A perturbative weighted quadrature-domain theorem}\label{app:cherednichenko-type}

We now prove a local realization statement for weighted quadrature domains near a
reference domain.

Fix a positive weight \(\mu\in C^2(\overline{U_2})\), where
\(U_1\subset U_0\subset U_2\) are bounded \(C^2\) domains. For a bounded domain
\(U\subset U_2\), define its weighted Cauchy transform by
\[
g_U(z):=\int_U \frac{\mu(w)}{z-w}\,dA(w),
\qquad z\in \C\setminus \overline U.
\]
This defines a holomorphic function on $\C\setminus \overline U$, and, in the sense of distributions on $\C$,
\[
\partial_{\overline z} g_U = \pi\,\mu\,\mathbf{1}_U.
\]

As before, every domain \(U\) sufficiently close to \(U_0\) may be written as a
normal graph over \(\partial U_0\):
\[
\partial U
=
\{x+\varphi(x)\nu(x):x\in \partial U_0\},
\]
with \(\varphi\in C^2(\partial U_0)\) small; we denote this domain by \(U_\varphi\).

\begin{proposition}[\cite{Cherednichenko}]\label{prop:cherednichenko-tailored}
Let \(U_1\subset U_0\subset U_2\) be bounded \(C^2\) domains, and let
\(\mu\in C^2(\overline{U_2})\) satisfy
\[
\mu>0 \qquad \text{on }\overline{U_2}.
\]
Then there exist \(\rho>0\) and \(\varepsilon>0\) with the following property.

Suppose \(g\) is holomorphic on \(\C\setminus \overline{U_1}\), continuous up to
\(\partial U_1\), and satisfies
\begin{equation}\label{eq:appendix-g-hypotheses}
\lim_{z\to\infty} g(z)=0,
\qquad
\lim_{z\to\infty} zg(z)=c_0>0,
\end{equation}
together with
\[
\sup_{z\in \partial U_1}|g(z)-g_{U_0}(z)|<\varepsilon.
\]
Then there exists a unique \(C^2\) function
\(\varphi\in C^2(\partial U_0)\), with \(\|\varphi\|_{C^2}<\rho\), such that the
associated domain \(U_\varphi\) satisfies
\[
U_1\subset U_\varphi\subset U_2
\]
and
\[
g_{U_\varphi}(z)=g(z)
\qquad \text{for all } z\in \C\setminus \overline{U_\varphi}.
\]
Moreover, the map
\[
g\mapsto \varphi(g)
\]
is \(C^1\) with respect to the \(C^2(\partial U_1)\) norm on the trace data and
the \(C^2(\partial U_0)\) norm on the boundary graph.
\end{proposition}

\begin{proof}
The proof is again an application of the implicit function theorem.

\medskip
\noindent\emph{Step 1: reduction to a boundary matching problem.}
Since both \(g\) and \(g_{U_\varphi}\) are holomorphic on
\(\C\setminus \overline{U_1}\) and vanish at infinity, it is enough to impose
equality on the boundary \(\partial U_1\). More precisely, if
\[
g_{U_\varphi}=g
\qquad \text{on }\partial U_1,
\]
then by the identity theorem the equality extends to all of
\(\C\setminus \overline{U_\varphi}\), because the difference is holomorphic on the
common exterior domain and vanishes on the closed curve \(\partial U_1\), which
lies strictly inside that domain once \(U_1\subset U_\varphi\). Equivalently,
since the zero set of a nontrivial holomorphic function is discrete, vanishing
on the one-dimensional set \(\partial U_1\) forces the difference to be
identically zero on the connected component of \(\C\setminus \overline{U_\varphi}\)
containing infinity.

Thus we define
\[
\mathcal G(\varphi)(z)
:=
g_{U_\varphi}(z)\big|_{\partial U_1},
\qquad z\in \partial U_1.
\]
We regard \(\mathcal G\) as a map from a neighborhood of \(0\) in
\(C^2(\partial U_0)\) into \(C^2(\partial U_1)\) (or, equivalently, into the
space of traces on \(\partial U_1\) of holomorphic functions on
\(\C\setminus \overline{U_1}\) vanishing at infinity).

The dependence of \(g_{U_\varphi}\) on \(\varphi\) is \(C^1\): indeed, writing
\(U_\varphi\) as the image of \(U_0\) under the graph diffeomorphism
\(\Phi_\varphi\), one obtains
\[
g_{U_\varphi}(z)
=
\int_{U_0}
\frac{\mu(\Phi_\varphi(w))J_\varphi(w)}{z-\Phi_\varphi(w)}\,dA(w),
\]
and differentiation with respect to \(\varphi\) is justified by the smoothness of
\(\mu\), the smallness of the perturbation, and the fact that
\(z\in \partial U_1\) stays a positive distance from \(\overline{U_0}\).

\medskip
\noindent\emph{Step 2: computation of the linearization.}
Let \(\psi\in C^2(\partial U_0)\). The first variation of the weighted Cauchy
transform with respect to a normal deformation of the domain is
\begin{equation}\label{eq:appendix-shape-derivative-cauchy}
D\mathcal G(0)[\psi](z)
=
\int_{\partial U_0}
\frac{\mu(w)\psi(w)}{z-w}\,d\sigma(w),
\qquad z\in \partial U_1.
\end{equation}
This is the standard Hadamard variation formula: to first order, varying the
domain in the normal direction by \(\psi\) replaces the bulk integral by a
boundary layer integral with density \(\mu\psi\).

Therefore the linearized operator
\[
L\psi(z):=\int_{\partial U_0}\frac{\mu(w)\psi(w)}{z-w}\,d\sigma(w),
\qquad z\in \partial U_1,
\]
is precisely the weighted Cauchy transform of the boundary density \(\mu\psi\).

\medskip
\noindent\emph{Step 3: invertibility of the linearized operator.}
We claim that \(L\) is injective. Suppose \(L\psi=0\) on \(\partial U_1\). Since
\(L\psi\) is holomorphic on \(\C\setminus \overline{U_0}\), vanishes at infinity,
and vanishes on the closed curve \(\partial U_1\) contained in that exterior
region, the identity theorem yields
\[
L\psi\equiv 0
\qquad \text{on }\C\setminus \overline{U_0}.
\]
Taking non-tangential boundary values on \(\partial U_0\) and using the standard
jump relation for the Cauchy transform, we conclude that
\[
\mu\psi=0
\qquad \text{on }\partial U_0.
\]
Since \(\mu>0\) on \(\partial U_0\), it follows that \(\psi=0\). Thus \(L\) is
injective.

To see surjectivity onto the natural trace space, let \(h\) be the trace on
\(\partial U_1\) of a holomorphic function on \(\C\setminus \overline{U_1}\)
which vanishes at infinity and has the admissible principal part there. Since
\(\partial U_0\subset \C\setminus \overline{U_1}\), the same function is
holomorphic in a full neighborhood of \(\partial U_0\). Denote by \(H\) this
holomorphic function on \(\C\setminus \overline{U_1}\). By the Cauchy integral
formula applied on \(\partial U_1\), the values of \(H\) on the annular region
\(\C\setminus (\overline{U_0}\cup U_1)\) are uniquely determined by its trace
\(h\). In particular, the trace \(H|_{\partial U_0}\) depends continuously on
\(h\).

Now apply the standard Sokhotski--Plemelj inversion on \(\partial U_0\): there
exists a unique density \(\sigma\in C^2(\partial U_0)\) such that
\[
H(z)=\int_{\partial U_0}\frac{\sigma(w)}{z-w}\,d\sigma(w),
\qquad z\in \C\setminus \overline{U_0},
\]
and \(\sigma\) depends continuously on the boundary trace \(H|_{\partial U_0}\).
Since \(\mu>0\) on \(\partial U_0\), setting
\[
\psi:=\frac{\sigma}{\mu}
\]
gives \(\psi\in C^2(\partial U_0)\) and
\[
L\psi(z)=H(z)
\qquad \text{for }z\in \partial U_1.
\]
Thus \(L\) is onto the trace space under consideration. Combining injectivity
and surjectivity, \(L\) is an isomorphism between
\(C^2(\partial U_0)\) and that trace space. In particular, after fixing any
convenient Banach norm on the target, \(L^{-1}\) is bounded by the open mapping
theorem.

\medskip
\noindent\emph{Step 4: application of the implicit function theorem.}
Since \(\mathcal G\) is \(C^1\) and \(D\mathcal G(0)=L\) is invertible, the
implicit function theorem gives \(\rho>0\) and \(\varepsilon>0\) such that, for
every \(g\) with
\[
\|g-g_{U_0}\|_{C^2(\partial U_1)}<\varepsilon,
\]
there exists a unique \(\varphi\in C^2(\partial U_0)\), \(\|\varphi\|_{C^2}<\rho\),
such that
\[
\mathcal G(\varphi)=g|_{\partial U_1}.
\]
Moreover, the implicit function theorem yields that the solution map
\[
g\mapsto \varphi(g)
\]
is \(C^1\) on a neighborhood of \(g_{U_0}|_{\partial U_1}\) in the trace space.
By construction,
\[
g_{U_\varphi}=g
\qquad \text{on }\partial U_1,
\]
hence on the whole common exterior domain by holomorphic continuation.

Finally, by taking \(\rho\) smaller if necessary, every such \(U_\varphi\) still
satisfies
\[
U_1\subset U_\varphi\subset U_2.
\]
This proves the proposition.
\end{proof}

\begin{remark}
The only point in Proposition~\ref{prop:cherednichenko-tailored} that is slightly
nontrivial is the invertibility of the linearized map. In the present perturbative
regime, the key observation is that the Cauchy transform of a boundary density is
determined by its trace on any inner contour, and the standard jump relation then
recovers the density itself. This is exactly the mechanism needed in the proof of
Theorem~\ref{thm:perturb}.
\end{remark}

\section{Proof of Theorem \ref{thm:perturb}} 

\begin{proof}[Proof of Theorem \ref{thm:perturb}]
We again use the Bargmann transform. Let
\[
P(z)=1+\sum_{k=1}^N b_k z^k,
\]
where the coefficients $b_k$ are chosen so that the Bargmann transform of
\[
f_0=h_0+\sum_{k=1}^N a_k h_k
\]
is exactly $P$. Then the condition that $f_0$ be an eigenfunction of $\mathcal{L}_U$ with eigenvalue $\lambda$ is equivalent to
\begin{equation}\label{eq:eigenfunction-general-new}
\int_U P(z)e^{-\pi|z|^2}e^{\pi \overline z w}\,dA(z)=\lambda P(w),
\qquad w\in \C.
\end{equation}

\medskip
\noindent\textbf{Step 0: a first reduction.}
For a polynomial
\[
R(z)=\alpha_0+\alpha_1 z+\cdots+\alpha_n z^n,
\]
define the differential operator
\[
R(\partial_w/\pi):=\sum_{k=0}^n \alpha_k\Big(\frac{\partial_w}{\pi}\Big)^k.
\]
Applying $P(\partial_w/\pi)$ to both sides of \eqref{eq:eigenfunction-general-new}, and using
\[
\Big(\frac{\partial_w}{\pi}\Big)^j e^{\pi \overline z w}=\overline z^{\,j} e^{\pi \overline z w},
\]
which is legitimate because \(P(\partial_w/\pi)\) is a finite-order differential
operator and the left-hand side of \eqref{eq:eigenfunction-general-new} is an
entire function of \(w\),
we obtain
\begin{equation}\label{eq:almost-reduced-new}
\int_U |P(z)|^2 e^{-\pi|z|^2} e^{\pi \overline z w}\,dA(z)=\lambda Q(w),
\qquad w\in\C,
\end{equation}
where
\[
Q(w):=P(\partial_w/\pi)\big(P(w)\big).
\]
Thus, our first task is to construct a domain $U$ satisfying \eqref{eq:almost-reduced-new}.

\medskip
\noindent\textbf{Step 1: construction of an auxiliary domain.}
We first show that there exists a $C^2$ domain $U_0$ such that
\begin{equation}\label{eq:first-domain-new}
\int_{U_0} |P(z)|^2 e^{-\pi|z|^2} e^{\pi \overline z w}\,dA(z)=\lambda,
\qquad w\in\C.
\end{equation}

Let $\psi\in C_c^\infty(\C)$ be a nonnegative radial function such that
\[
\supp \psi \subset D_2(0)\setminus D_1(0),
\]
and define
\[
\psi_\delta(z):=\delta^{-2}\psi(z/\delta).
\]
Since $\psi_\delta$ is radial, a polar-coordinate computation shows that
\[
\int_{\C} \psi_\delta(z)e^{-\pi|z|^2}e^{\pi \overline z w}\,dA(z)=c_\delta,
\]
where
\[
c_\delta:=\int_{\C}\psi_\delta(z)e^{-\pi|z|^2}\,dA(z).
\]
Thus \eqref{eq:first-domain-new} is equivalent to
\begin{equation}\label{eq:second-domain-new}
\int_{U_0} |P(z)|^2 e^{-\pi|z|^2} e^{\pi \overline z w}\,dA(z)
=
\frac{\lambda}{c_\delta}
\int_{\C}\psi_\delta(z)e^{-\pi|z|^2}e^{\pi \overline z w}\,dA(z).
\end{equation}

We shall construct directly a compactly supported distributional solution of the
Poisson equation
\begin{equation}\label{eq:inverse-first-new}
\Delta u_0
=
-|P(z)|^2 e^{-\pi|z|^2}\mathbf{1}_{U_0}
+\frac{\lambda}{c_\delta}\psi_\delta(z)e^{-\pi|z|^2}.
\end{equation}
Once such a solution has been constructed, the moment identity
\eqref{eq:first-domain-new} follows by applying the same Fourier-analytic
divisibility argument used in the proof of Theorem~\ref{thm:eig-loc-U} to the
corresponding compactly supported distribution
\[
\omega_0:=
|P(z)|^2e^{-\pi|z|^2}\mathbf 1_{U_0}
-\frac{\lambda}{c_\delta}\psi_\delta(z)e^{-\pi|z|^2}.
\]
Indeed, \eqref{eq:inverse-first-new} says exactly that \(\Delta u_0=-\omega_0\),
so the vanishing of the Fourier transform of \(\omega_0\) on the two
characteristic lines is equivalent to \eqref{eq:first-domain-new}.

We first consider the unperturbed case $P\equiv 1$. By the theorem of Daubechies \cite{Daubechies}, there exists a disk $D_0$ such that
\[
\int_{D_0} e^{-\pi|z|^2}e^{\pi \overline z w}\,dA(z)=\lambda,
\qquad w\in\C.
\]
Hence there exists a compactly supported solution $\widetilde u_0$ of
\begin{equation}\label{eq:unperturbed-free-boundary-new}
\Delta \widetilde u_0
=
- e^{-\pi|z|^2}\mathbf{1}_{D_0}
+\frac{\lambda}{c_\delta}\psi_\delta(z)e^{-\pi|z|^2}.
\end{equation}
If $\delta>0$ is sufficiently small, then $\supp\psi_\delta\subset D_0$. Since all
data are radial, the solution $\widetilde u_0$ is radial as well. We claim that,
for $\delta$ sufficiently small,
\[
\widetilde u_0>0 \qquad \text{in } D_0.
\]
Indeed, let \(u_*\) denote the compactly supported radial solution of
\[
\Delta u_*=-e^{-\pi|z|^2}\mathbf{1}_{D_0}+\lambda \delta_0.
\]
By the defining property of \(D_0\), this is exactly the reference obstacle-type
potential associated with the disk, and \(u_*>0\) in \(D_0\): away from the origin,
\(-\Delta u_*=e^{-\pi|z|^2}>0\) in \(D_0\setminus\{0\}\), \(u_*=0\) on \(\partial
D_0\), and the logarithmic singularity created by the mass \(\lambda\delta_0\) is
positive. On the other hand, the family
\[
\frac{\lambda}{c_\delta}\psi_\delta(z)e^{-\pi|z|^2}
\]
converges to \(\lambda \delta_0\) in the sense of distributions as \(\delta\to 0\),
so the corresponding compactly supported radial solutions \(\widetilde u_0\)
converge locally uniformly on \(\overline D_0\setminus\{0\}\) to \(u_*\). Since
\(u_*\) is strictly positive on compact subsets of \(D_0\setminus\{0\}\), and the
support of \(\psi_\delta\) lies in an arbitrarily small neighborhood of the
origin, it follows that \(\widetilde u_0>0\) throughout \(D_0\) for \(\delta>0\)
sufficiently small.

We shall now invoke the perturbative free-boundary theorem of Isakov, given by Proposition \ref{prop:isakov-tailored}, verifying its hypotheses carefully. Set
\[
f_0(z):=e^{-\pi|z|^2}\Big(1-\frac{\lambda}{c_\delta}\psi_\delta(z)\Big).
\]
Since $\psi_\delta\in C_c^\infty(\C)$, we have $f_0\in C^\infty(\R^2)$, hence in particular $f_0\in C^5(\R^2)$. If $\delta$ is chosen so small that $\supp\psi_\delta\subset D_0$, then on $\partial D_0$ one has
\[
f_0(z)=e^{-\pi|z|^2}>0.
\]
Next, define
\[
v_0:=-\widetilde u_0.
\]
Then \eqref{eq:unperturbed-free-boundary-new} becomes
\[
-\Delta v_0=f_0\,\mathbf{1}_{D_0}
\qquad \text{in }\C,
\]
with \(v_0=0\) on \(\partial D_0\) and \(v_0<0\) in \(D_0\). Thus the hypotheses
of Proposition \ref{prop:isakov-tailored} are satisfied with \(\Omega_0=D_0\).
Now define
\[
f(z):=e^{-\pi|z|^2}\Big(|P(z)|^2-\frac{\lambda}{c_\delta}\psi_\delta(z)\Big).
\]
Since $P$ is a polynomial, $f\in C^\infty(\R^2)$ as well. Moreover, for any fixed disk $\Omega'\supset D_0$, the $C^5(\Omega')$ norm of $f-f_0$ is controlled by the coefficients of $P-1$:
\[
\|f-f_0\|_{C^5(\Omega')}\le C(\Omega',N,\lambda)\max_{1\le k\le N}|b_k|.
\]
Hence, if the coefficients $b_k$ are sufficiently small, we may arrange that
\[
\|f-f_0\|_{C^5(\Omega')}<\varepsilon_1.
\]
Applying Proposition \ref{prop:isakov-tailored}, we obtain a $C^2$ domain $U_0$ and a function $v$ such that
\[
-\Delta v=f\,\mathbf{1}_{U_0}
\qquad \text{in }\C,
\]
with $v=0$ outside $U_0$ and $v<0$ in $U_0$. Set
\[
u_0:=-v.
\]
Then
\[
\Delta u_0=f\,\mathbf{1}_{U_0}.
\]
If the perturbation is sufficiently small, then $U_0$ remains a small deformation of $D_0$, and since $\supp\psi_\delta\subset D_0$ we still have
\[
\supp\psi_\delta\subset U_0.
\]
Therefore
\[
\Delta u_0
=
-|P(z)|^2e^{-\pi|z|^2}\mathbf{1}_{U_0}
+\frac{\lambda}{c_\delta}\psi_\delta(z)e^{-\pi|z|^2},
\]
so \eqref{eq:inverse-first-new} holds. Hence \eqref{eq:first-domain-new} is proved.

\medskip
\noindent\textbf{Step 2: construction of a domain satisfying \eqref{eq:almost-reduced-new}.}
We now convert \eqref{eq:almost-reduced-new} into the weighted Cauchy-transform
form required by Proposition~\ref{prop:cherednichenko-tailored}. Restrict
\eqref{eq:almost-reduced-new} to real values \(w=t>0\), multiply by
\(e^{-\pi \zeta t}\), and integrate in \(t>0\). If
\[
\zeta>\sup_{z\in U}|z|,
\]
Fubini's theorem gives
\[
\int_U \left(\int_0^\infty e^{\pi(\overline z-\zeta)t}\,dt\right)|P(z)|^2e^{-\pi|z|^2}\,dA(z)
=
\lambda\int_0^\infty e^{-\pi\zeta t}Q(t)\,dt.
\]
Indeed, for real \(\zeta>\sup_{z\in U}|z|\) one has
\[
\Re(\overline z-\zeta)\le |z|-\zeta<0
\qquad \text{for every } z\in U,
\]
so the inner exponential is integrable uniformly in \(z\in U\).
Writing
\[
Q(w)=\alpha_0+\alpha_1 w+\cdots+\alpha_n w^n,
\]
we obtain
\[
\frac{1}{\pi}\int_U \frac{|P(z)|^2e^{-\pi|z|^2}}{\zeta-\overline z}\,dA(z)
=
\lambda \sum_{k=0}^n \alpha_k \int_0^\infty e^{-\pi \zeta t} t^k\,dt.
\]
Since
\[
\int_0^\infty e^{-\pi \zeta t} t^k\,dt
=
\frac{\Gamma(k+1)}{(\pi \zeta)^{k+1}},
\]
it follows that
\[
\frac{1}{\pi}\int_U \frac{|P(z)|^2e^{-\pi|z|^2}}{\zeta-\overline z}\,dA(z)
=
\lambda \sum_{k=0}^n \alpha_k \frac{\Gamma(k+1)}{\pi^{k+1}\zeta^{k+1}}.
\]

To put this into the form of a weighted Cauchy transform, let
\[
V:=\overline U,
\qquad
\mu(w):=|P(\overline w)|^2 e^{-\pi|w|^2}.
\]
Changing variables $w=\overline z$ yields
\[
\int_U \frac{|P(z)|^2e^{-\pi|z|^2}}{\zeta-\overline z}\,dA(z)
=
\int_V \frac{\mu(w)}{\zeta-w}\,dA(w)
=
g_V(\zeta).
\]
Hence
\begin{equation}\label{eq:analytic-construct-new}
g_V(\zeta)
=
\lambda \sum_{k=0}^n \alpha_k \frac{\Gamma(k+1)}{\pi^k \zeta^{k+1}},
\qquad \zeta\in \R_+, \ \zeta\gg 1.
\end{equation}
Since both sides are holomorphic for large $\zeta$, analytic continuation gives the desired identity on $\C\setminus V$.

Conversely, if a domain $V$ satisfies \eqref{eq:analytic-construct-new}, then by expanding at infinity and comparing coefficients one recovers the moment identities corresponding to \eqref{eq:almost-reduced-new} for the conjugate domain $U=\overline V$. Thus the conjugation is merely a bookkeeping device that places the kernel in the form \(1/(z-w)\) required by the weighted Cauchy transform formalism.

Let us denote by $g$ the rational function on the right-hand side of \eqref{eq:analytic-construct-new}. We now verify the hypotheses of Proposition \ref{prop:cherednichenko-tailored}. Since
\[
Q(w)=P(\partial_w/\pi)\big(P(w)\big),
\]
and $P(0)=1$, the constant term of $Q$ is
\[
Q(0)=1+\sum_{k=1}^N \frac{k!}{\pi^k} b_k^2.
\]
In particular, if the coefficients $b_k$ are sufficiently small, then
\[
Q(0)>0.
\]
Therefore the leading term in the expansion of $g$ at infinity is
\[
\frac{\lambda Q(0)}{z},
\]
so
\[
g(z)\to 0,
\qquad
z\,g(z)\to \lambda Q(0)>0
\qquad \text{as } z\to\infty.
\]
Thus \eqref{eq:appendix-g-hypotheses} holds, with \(c_0=\lambda Q(0)\).

We next verify the closeness condition on \(\partial U_1\). Since \(U_0\) was
obtained in Step 1 as a small perturbation of the Daubechies disk, we choose
\(U_1\subset U_0\subset U_2\) to be concentric disks with radii, say,
\((1-10^{-3})\) and \((1+10^{-3})\) times the radius of \(D_0\). By Step 1,
\[
\int_{U_0}|P(z)|^2e^{-\pi|z|^2}e^{\pi \overline zw}\,dA(z)=\lambda,
\qquad w\in \C.
\]
Repeating the Laplace-transform argument above therefore gives
\[
g_{U_0}(z)=\frac{\lambda}{z}
\qquad \text{for } z\in \C\setminus \overline{U_0},
\]
where \(g_{U_0}\) denotes the weighted Cauchy transform of the conjugate domain
\(\overline{U_0}\) with weight \(\mu(w)=|P(\overline w)|^2e^{-\pi|w|^2}\).
Consequently, on \(\partial U_1\),
\[
g(z)-g_{U_0}(z)
=
\lambda \sum_{k=1}^n \alpha_k \frac{\Gamma(k+1)}{\pi^k z^{k+1}}
+
\lambda\,(Q(0)-1)\,\frac1z .
\]
Since \(\partial U_1\) is a fixed circle staying a positive distance from the
origin, it follows that
\[
\sup_{z\in \partial U_1}|g(z)-g_{U_0}(z)|
\le
C(U_1,\lambda)\Big(|Q(0)-1|+\sum_{k=1}^n |\alpha_k|\Big).
\]
Each coefficient \(\alpha_k\) of \(Q-1\) depends polynomially on the
coefficients of \(P-1\), and therefore tends to zero with
\(\max_{1\le j\le N}|b_j|\). Hence, by choosing the coefficients sufficiently
small, we may ensure that
\[
\sup_{z\in \partial U_1}|g(z)-g_{U_0}(z)|<\varepsilon_0.
\]
Therefore Proposition \ref{prop:cherednichenko-tailored} applies and yields a \(C^2\) domain \(\widetilde U\) such that
\[
g_{\widetilde U}(z)=g(z),
\qquad z\in \C\setminus \widetilde U.
\]
Replacing \(\widetilde U\) by its conjugate if necessary, we obtain a domain \(U\) satisfying \eqref{eq:almost-reduced-new}.

Finally, since \(U\subset U_2\) and \(U_2\) is fixed once \(\lambda\) is fixed, by shrinking the admissible size of the coefficients of \(P-1\) once more we may ensure that
\[
P(z)\neq 0
\qquad \text{for all } z\in U_2.
\]
In particular,
\[
P(z)\neq 0
\qquad \text{for all } z\in U.
\]

\medskip
\noindent\textbf{Step 3: recovery of the original eigenvalue equation.} Let
\[
G(w):=\int_U P(z)e^{-\pi|z|^2}e^{\pi \overline z w}\,dA(z).
\]
By construction, \eqref{eq:almost-reduced-new} holds, so applying \(P(\partial_w/\pi)\) to \(G\) gives
\[
P(\partial_w/\pi)G(w)=P(\partial_w/\pi)(\lambda P)(w),
\qquad w\in \C.
\]
Hence
\[
H:=G-\lambda P
\]
satisfies
\[
P(\partial_w/\pi)H=0.
\]

Introduce the characteristic polynomial
\[
\widetilde P(\zeta):=P(\zeta/\pi).
\]
If \(\zeta_1,\dots,\zeta_\ell\) are the zeros of \(\widetilde P\), with multiplicities \(m_1,\dots,m_\ell\), then standard ODE theory implies that every entire solution of
\[
P(\partial_w/\pi)H=0
\]
of at most exponential growth is of the form
\[
H(w)=\sum_{j=1}^\ell R_j(w)e^{\zeta_j w},
\]
where each \(R_j\) is a polynomial of degree strictly less than \(m_j\). Since \(P\) has no zeros on \(U\), and in fact none on the larger fixed set \(U_2\), by choosing the coefficients of \(P-1\) sufficiently small we may also assume that all zeros of \(\widetilde P\) stay at distance at least \(10\sup_{z\in U}|z|\) from the origin:
\[
|\zeta_j|>10\sup_{z\in U}|z|,
\qquad 1\le j\le \ell.
\]
Suppose, for contradiction, that some \(R_j\) is not identically zero. Reorder so that
\[
|\zeta_1|=\max_j |\zeta_j|
\qquad\text{and}\qquad
R_1\not\equiv 0.
\]
Then along the ray
\[
w=t\,\overline{\zeta_1},
\qquad t>0,
\]
the dominant exponential term yields
\[
|H(t\overline{\zeta_1})|\ge c\,e^{t|\zeta_1|^2}
\]
for all sufficiently large \(t\). On the other hand, since \(P\) is a fixed polynomial and \(U\) is bounded, if
\[
M:=\sup_{z\in U}|z|,
\]
there exists \(C>0\) such that
\[
|G(w)|\le C(1+|w|)^N e^{\pi M|w|},
\qquad w\in \C.
\]
Indeed,
\[
|G(w)|
\le
\int_U |P(z)|e^{-\pi|z|^2}e^{\pi |z||w|}\,dA(z)
\le
C(1+|w|)^N e^{\pi M|w|},
\]
because \(P\) has degree at most \(N\) and \(U\) is bounded. Since every
polynomial is dominated by \(e^{\varepsilon |w|}\) for arbitrary
\(\varepsilon>0\), it follows that for every \(\varepsilon>0\) there exists
\(C_\varepsilon>0\) such that
\[
|G(w)|\le C_\varepsilon e^{(\pi+\varepsilon)M|w|},
\qquad w\in \C.
\]
As \(\lambda P(w)\) grows only polynomially, the same estimate holds for \(H\):
\[
|H(w)|\le C'_\varepsilon e^{(\pi+\varepsilon)M|w|}.
\]
Evaluating at \(w=t\overline{\zeta_1}\) gives
\[
|H(t\overline{\zeta_1})|
\le
C'_\varepsilon e^{(\pi+\varepsilon)Mt|\zeta_1|}.
\]
Comparing the two inequalities and letting \(t\to\infty\), we obtain
\[
|\zeta_1|\le (\pi+\varepsilon)M.
\]
Since \(\varepsilon>0\) is arbitrary, this implies
\[
|\zeta_1|\le \pi M,
\]
contradicting our choice
\[
|\zeta_1|>10M.
\]
Hence every \(R_j\) vanishes identically, so
\[
H\equiv 0.
\]
Equivalently,
\[
G=\lambda P,
\]
which is exactly \eqref{eq:eigenfunction-general-new}. This proves that \(P\), and therefore \(f_0\), is an eigenfunction of the localization operator associated with the domain \(U\) with eigenvalue \(\lambda\).
\end{proof}

\section{A Classical Rigidity Corollary}\label{sec:classical-rigidity}

\subsection{First proof}

\begin{proof}[First proof under the additional assumption that \(\partial U\) is \(C^2\)]
The Bargmann transform of a Hermite function $h_k$ is a nonzero multiple of the monomial $z^k$. Thus, arguing as in \cite{Abreu-Doerfler}, the assumption that $h_k$ is an eigenfunction of the localisation operator $\mathcal{L}_U$ is equivalent to the identity
\[
\int_U z^k e^{\pi w\overline{z}} e^{-\pi |z|^2}\,dA(z)=\lambda w^k,
\qquad w\in \C.
\]
Both sides are entire functions of $w$. Differentiating $k$ times with respect to $w$, we obtain
\begin{equation}\label{eq:fourier-delta}
\int_{\C} \mathbf{1}_U(z)\,|z|^{2k}e^{-\pi |z|^2} e^{\pi \overline{z} w}\,dA(z)=c(\lambda,k),
\end{equation}
for some real constant $c(\lambda,k)$.

Consider now the compactly supported distribution
\[
\eta:=|z|^{2k}e^{-\pi |z|^2}\mathbf{1}_U(z)-c(\lambda,k)\delta_0.
\]
Writing $z=x+iy$, we may rewrite the exponential factor as
\[
e^{\pi \overline{z}w}=e^{\pi(x-iy)w}=e^{\pi xw}e^{-i\pi yw}.
\]
Hence \eqref{eq:fourier-delta} implies that the Fourier transform $\widehat{\eta}$, viewed as an entire function on $\C^2$ by Paley--Wiener, satisfies
\[
\widehat{\eta}(w,-iw)=0,\qquad \forall w\in\C.
\]
Since $c(\lambda,k)\in \R$, taking complex conjugates in \eqref{eq:fourier-delta} after replacing $w$ by $\overline{w}$ also yields
\[
\widehat{\eta}(w,iw)=0,\qquad \forall w\in\C.
\]

Applying Lemma \ref{lemma:fourier-to-fboundary} to the distribution $\eta$, we obtain a function $u\in L^2(\R^2)$ such that
\begin{equation}\label{eq:potato}
\Delta u=f(|z|)\mathbf{1}_U-c(\lambda,k)\delta_0
\end{equation}
in the sense of distributions, where
\[
f(|z|)=|z|^{2k}e^{-\pi |z|^2}.
\]
Writing $u$ explicitly as a logarithmic potential, we conclude that $U$ satisfies
\begin{equation}\label{eq:potential}
\int_U \log|z-z'|\,f(|z'|)\,dA(z')=c(\lambda,k)\log|z|,
\qquad z\in \C\setminus (U\cup\{0\}).
\end{equation}

We now reduce to the case $0\notin U$. Indeed, if $0\in U$, then for every sufficiently small $\rho>0$ with $D_\rho(0)\subset U$, the potential generated by
\[
f(|z|)\mathbf{1}_{D_\rho(0)}
\]
is radial outside $D_\rho(0)$, hence equals a constant plus a multiple of $\log|z|$ there. Therefore, after subtracting such a disk from $U$, identity \eqref{eq:potential} is preserved outside the new set, up to replacing the constant $c(\lambda,k)$ by another real constant. Thus, without loss of generality, we may assume that
\[
0\notin U.
\]
Moreover, $0\notin \partial U$, because the left-hand side of \eqref{eq:potential} stays finite as $z\to 0$, whereas the right-hand side diverges.

Define
\[
\Psi_U(z):=\int_U \log|z-z'|\,f(|z'|)\,dA(z').
\]
We claim that $\Psi_U$ is \emph{locally radial}: for each $z_0\in U$, there exists a ball $B(z_0)\subset U$ such that $\Psi_U$ agrees on $B(z_0)$ with a radial function.

To prove this, let $B\supset U$ be a ball centered at $0$, and define
\[
\Phi(z):=\int_B \log|z-z'|\,f(|z'|)\,dA(z').
\]
Since both $B$ and the weight $f(|z'|)$ are radial, $\Phi$ is radial. In particular, $\nabla \Phi(z)$ is parallel to $z$ for all $z\in \R^2$. Also, for $z\in B$,
\[
\Delta \Phi(z)=c\,f(|z|)
\]
for a suitable constant $c>0$.

Now $\Psi_U$ satisfies \eqref{eq:potential}, so on $\partial U$ its gradient is parallel to $z$ by continuity. Therefore the harmonic function
\[
v:=\Psi_U-\Phi
\]
satisfies
\[
\nabla v(x)\ \text{is parallel to}\ x,\qquad x\in \partial U.
\]
The next lemma shows that this forces $v$, and hence $\Psi_U$, to be locally radial.

\begin{lemma}
Let $V\subset \R^2$ be a domain, and suppose there exists a harmonic function
$v\in C^1(\overline{V})$ such that $\nabla v(x)$ is parallel to $x$ for every $x\in \partial V$. Then for each circle $S\subset \R^2$, the function $v$ is constant on every connected component of $S\cap V$.
\end{lemma}

\begin{proof}
Consider the angular derivatives
\[
g_{i,j}(x)=x_i\,\partial_j v(x)-x_j\,\partial_i v(x).
\]
Since $v$ is harmonic, each $g_{i,j}$ is harmonic in $V$. By the boundary assumption,
\[
g_{i,j}=0 \qquad \text{on } \partial V.
\]
Hence, by the maximum principle,
\[
g_{i,j}\equiv 0 \qquad \text{in } V.
\]
Therefore $\nabla v(x)$ is parallel to $x$ for every $x\in V$. In particular, on any circle $S$, the tangential derivative of $v$ vanishes on $S\cap V$, so $v$ is constant on each connected component of $S\cap V$.
\end{proof}

\begin{lemma}\label{lem:local-to-global-radial}
Let \(W\subset \R^2\setminus\{0\}\) be a connected open set, and let
\(u\in C(W)\). Suppose that for every \(x\in W\) there exist \(r_x>0\) and a
scalar function \(\phi_x\) such that
\[
u(y)=\phi_x(|y|)
\qquad\text{for all } y\in B(x,r_x)\subset W.
\]
Then there exists a single scalar function \(\phi\), defined on the interval
\[
I_W:=\{|y|:y\in W\},
\]
such that
\[
u(y)=\phi(|y|)
\qquad\text{for all } y\in W.
\]
\end{lemma}

\begin{proof}
Choose for each \(x\in W\) a ball \(B_x:=B(x,r_x)\subset W\) on which \(u\) is a
radial function of \(|y|\). If \(B_x\cap B_{x'}\neq\emptyset\), then
\[
u(y)=\phi_x(|y|)=\phi_{x'}(|y|)
\qquad\text{for every }y\in B_x\cap B_{x'}.
\]
Since \(B_x\cap B_{x'}\) is open and avoids the origin, the set of radii attained
on \(B_x\cap B_{x'}\) contains a nontrivial open interval. Hence \(\phi_x\) and
\(\phi_{x'}\) agree on that interval.

Fix \(x_*\in W\). For any \(y\in W\), choose a polygonal path in \(W\) joining
\(x_*\) to \(y\). By compactness of the image of the path, it may be covered by
finitely many balls
\[
B_{x_1},\dots,B_{x_m}
\]
with successive overlaps nonempty, \(x_1=x_*\), and \(y\in B_{x_m}\). The previous
paragraph shows that the corresponding local radial profiles agree on the
overlaps, and therefore patch together uniquely along the chain. This defines a
single scalar function \(\phi\) on \(I_W\) such that \(u(z)=\phi(|z|)\) for every
\(z\in W\). The definition is independent of the chosen chain because any two
chains can be concatenated and compared in the same way.
\end{proof}

We now conclude the proof. Since \(U\) is simply connected, it suffices to rule
out the possibility that \(U\) is nonradial. Suppose therefore that \(U\) is not
a disk. Then \(U\) is neither a disk nor an annulus, so there exist radii
$0<r_1<r_2$ such that the annulus
\[
A:=\{z\in \R^2:\ r_1<|z|<r_2\}
\]
meets both $U$ and $\partial U$. Choose a point
\[
x_0\in \partial U\cap A
\]
and let
\[
\widetilde{U}_0:=U\cap B_r(x_0),
\]
for $r>0$ sufficiently small. Shrinking \(r\) if necessary, we may assume that
\(\widetilde U_0\) is connected and that \(0\notin \overline{\widetilde U_0}\).
Since \(\Psi_U\) is locally radial in \(U\), Lemma~\ref{lem:local-to-global-radial}
applied to \(W=\widetilde U_0\) yields a scalar function \(\psi_U\) such that
\[
\Psi_U(x)=\psi_U(|x|),\qquad x\in \widetilde{U}_0.
\]

On the other hand, by \eqref{eq:potential}, for every $z\in \partial \widetilde{U}_0$ we have
\[
\nabla \Psi_U(z)=c\,\frac{z}{|z|^2}
\]
for a suitable constant $c\in \R$. Since $\Psi_U(x)=\psi_U(|x|)$ on $\widetilde{U}_0$, this yields
\[
\psi_U'(|z|)\frac{z}{|z|}=c\,\frac{z}{|z|^2},
\qquad z\in \partial \widetilde{U}_0,
\]
hence
\[
\psi_U'(r)=\frac{c}{r}
\]
for all radii attained by points of $\partial \widetilde{U}_0$. Since
\(\partial U\) is \(C^2\) and \(x_0\in \partial U\), after shrinking \(r\) once
more we may assume that the set of radii attained by
\(\partial \widetilde U_0=\partial U\cap B_r(x_0)\) contains a nontrivial open
interval \(I\subset (0,\infty)\). Therefore
\[
\psi_U'(r)=\frac{c}{r}
\]
for every \(r\in I\), whence
\[
\psi_U(r)=c\log r+C_0
\qquad\text{for every }r\in I
\]
for some constant \(C_0\). Let
\[
W_I:=\{x\in \widetilde U_0:\ |x|\in I\}.
\]
Then \(W_I\) is a nonempty open subset of \(U\), and on \(W_I\) we have
\[
\Delta \Psi_U\equiv 0 \qquad \text{in } W_I.
\]

This is impossible, because by construction
\[
\Delta \Psi_U=c' f(|z|)
\]
inside $U$, where $f(|z|)=|z|^{2k}e^{-\pi |z|^2}$ vanishes only at the origin,
while \(0\notin \widetilde U_0\) and hence \(0\notin W_I\). Thus
\(\Delta\Psi_U\) cannot vanish on \(W_I\), a contradiction. This contradiction
shows that $U$ must be radial. Since $U$ is assumed bounded and simply
connected, it follows that $U$ is a disk.
\end{proof}

\subsection{Second proof}

\begin{proof}[Second proof of Theorem \ref{thm:eig-loc-U}]
We keep the notation
\[
f(|z|):=|z|^{2k}e^{-\pi|z|^2},
\qquad
c:=c(\lambda,k),
\]
from the first proof.  Starting again from \eqref{eq:fourier-delta} and applying
Lemma~\ref{lemma:fourier-to-fboundary}, we obtain a compactly supported distributional solution
of
\begin{equation}\label{eq:potato-second-proof}
\Delta u=f(|z|)\mathbf{1}_U-c\,\delta_0
\qquad\text{in }\mathcal D'(\R^2).
\end{equation}
Using the planar fundamental solution \((2\pi)^{-1}\log|z|\), we may write
\begin{equation}\label{eq:explicit-u-second-proof}
u(z):=\frac{1}{2\pi}\int_U \log|z-z'|\,f(|z'|)\,dA(z')-\frac{c}{2\pi}\log|z|.
\end{equation}
By \eqref{eq:potential}, this function satisfies
\[
u\equiv 0
\qquad\text{on }\R^2\setminus (U\cup\{0\}).
\]

We first show that \(0\in U\). Since the singular term in
\eqref{eq:potato-second-proof} is \(-c\delta_0\), one must have \(0\in\overline U\). If
\(0\in\partial U\), then there exist exterior points \(z_j\to 0\) with
\[
u(z_j)=0.
\]
On the other hand, the first term in \eqref{eq:explicit-u-second-proof} remains finite as
\(z\to 0\), because \(f\in L^\infty(U)\) and \(\log|z-z'|\) is locally integrable, whereas the
second term equals \(-\frac{c}{2\pi}\log|z|\to +\infty\). This contradicts the exterior identity
\(u(z_j)=0\). Hence \(0\notin\partial U\), and therefore
\[
0\in U.
\]

We next show directly from \eqref{eq:explicit-u-second-proof} that
\[
u\in C^1(\R^2\setminus\{0\}).
\]
Indeed, if we write
\[
g(z):=f(|z|)\mathbf{1}_U(z)\in L^\infty_c(\R^2),
\]
then
\[
u(z)=\frac{1}{2\pi}\int_{\R^2}\log|z-z'|\,g(z')\,dA(z')-\frac{c}{2\pi}\log|z|.
\]
Since the kernel \(K(\zeta)=\zeta/|\zeta|^2\) belongs to \(L^1_{\mathrm{loc}}(\R^2)\), the
translation map \(z\mapsto K(z-\cdot)\) is continuous in \(L^1_{\mathrm{loc}}\). It follows that
\[
\nabla\!\left(\int_{\R^2}\log|z-z'|\,g(z')\,dA(z')\right)
=
\int_{\R^2}\frac{z-z'}{|z-z'|^2}\,g(z')\,dA(z')
\]
defines a continuous function of \(z\in\R^2\). Hence \(u\in C^1(\R^2\setminus\{0\})\).

Now let \(z_*\in\partial U\). Since \(0\in U\), we have \(z_*\neq 0\), and since
\[
u\equiv 0
\qquad\text{on }\R^2\setminus \overline U,
\]
both \(u\) and \(\nabla u\) vanish on the exterior side of \(\partial U\). By continuity of
\(\nabla u\) at \(z_*\), we conclude that
\begin{equation}\label{eq:dz-zero-second-proof}
\partial_z u=0
\qquad\text{on }\partial U.
\end{equation}

Define
\[
\Phi(t):=\int_0^t s^k e^{-\pi s}\,ds
\qquad\text{and}\qquad
\Psi(z):=\frac{\Phi(|z|^2)}{z},
\quad z\neq 0.
\]
Then
\[
\partial_{\bar z}\Psi(z)=f(|z|)
\qquad (z\neq 0),
\]
and \(\Phi(0)=0\), so the same identity holds distributionally on \(U\). Now set
\[
F(z):=4\partial_z u(z)-\Psi(z)+\frac{c}{\pi z},
\qquad z\in U\setminus\{0\}.
\]
Using \(\partial_{\bar z}(1/z)=\pi\delta_0\), \(\Delta=4\partial_z\partial_{\bar z}\), and
\eqref{eq:potato-second-proof}, we obtain in \(\mathcal D'(U)\)
\[
\partial_{\bar z}F
=
\Delta u-f(|z|)+c\delta_0
=
\bigl(f(|z|)-c\delta_0\bigr)-f(|z|)+c\delta_0
=0.
\]
Hence \(F\) extends holomorphically across \(0\) and is holomorphic on all of \(U\).

Define
\[
W(z):=\pi z\,F(z).
\]
Then \(W\) is holomorphic on \(U\) and \(W(0)=0\). Since \(\partial_z u\) is continuous up to
\(\partial U\), so is \(W\), and by \eqref{eq:dz-zero-second-proof},
\[
W|_{\partial U}
=
\pi z\left(4\partial_z u-\frac{\Phi(|z|^2)}{z}+\frac{c}{\pi z}\right)\Big|_{\partial U}
=
c-\pi\Phi(|z|^2).
\]
In particular, \(W|_{\partial U}\) is real-valued. The imaginary part of \(W\) is therefore
harmonic in \(U\), continuous up to \(\partial U\), and vanishes on \(\partial U\). By the
maximum principle,
\[
\Im W\equiv 0
\qquad\text{in }U.
\]
Since a holomorphic function with zero imaginary part is constant, there exists \(\kappa\in\R\)
such that
\[
W\equiv \kappa
\qquad\text{in }U.
\]

Restricting to the boundary gives
\[
c-\pi\Phi(|z|^2)=\kappa
\qquad\text{for all }z\in\partial U.
\]
Since \(\Phi\) is strictly increasing, \(|z|\) is constant on \(\partial U\). Thus
\(\partial U\) is contained in a circle centered at the origin. Let \(R>0\) be such that
\(\partial U\subset \{|z|=R\}\). Since \(U\) is open and connected, it cannot meet both
\(\{|z|<R\}\) and \(\{|z|>R\}\), because these are disjoint open sets and \(U\cap \{|z|=R\}=\varnothing\).
As \(U\) is bounded, we must therefore have \(U\subset \{|z|<R\}\). If this inclusion were strict,
then \(\{|z|<R\}\setminus U\) would contain a point, and hence, being open, would contribute a
boundary point of \(U\) lying strictly inside the disk \(\{|z|<R\}\), contradicting
\(\partial U\subset \{|z|=R\}\). Therefore
\[
U=\{|z|<R\}
\]
that is, \(U\) is a disk.
\end{proof}

\begin{remark}
The two proofs use the boundary assumptions in rather different ways.

In the first proof, once the logarithmic potential \(\Psi_U\) is written explicitly, no elliptic
boundary regularity is needed to conclude that \(\Psi_U\in C^1(\R^2\setminus\{0\})\): this follows
directly from convolution with the locally integrable kernel \(\nabla\log|z|\). Thus the local
radiality argument itself only uses continuity of the gradient up to the boundary. The genuine
extra input appears only in the final geometric step, where one needs that, near a suitable point
\(x_0\in \partial U\), the set \(\partial U\cap B_r(x_0)\) contains a nontrivial boundary arc, so
that the radii attained there form an interval. Accordingly, the first proof does not really need
\(C^2\) regularity; it works under any hypothesis guaranteeing such a local arc structure, for
instance when \(U\) is a Jordan domain. In particular, any Lipschitz, \(C^1\), or \(C^2\) boundary
is more than sufficient.

By contrast, the second proof does not use any boundary regularity at all. The explicit formula
\eqref{eq:explicit-u-second-proof} already yields \(u\in C^1(\R^2\setminus\{0\})\), and since
\(u\equiv 0\) on the exterior open set \(\R^2\setminus \overline U\), continuity implies that both
\(u\) and \(\nabla u\) vanish on \(\partial U\). This is exactly the boundary information needed to
construct the holomorphic function \(W\) and conclude that \(|z|\) is constant on \(\partial U\).
Hence the second proof recovers the theorem under the original assumptions of Abreu and D\"orfler,
namely for arbitrary bounded simply connected domains, with no boundary smoothness hypothesis. 
\end{remark}

\section{Proof of Theorem \ref{thm:GGRT-sharp}}

We prove that the estimate in Corollary~1.4 of \cite{Gomez-Guerra-Ramos-Tilli}
is sharp by constructing a family of domains which are perturbations of the disk,
whose first localization eigenvalue has the same first-order behaviour as the
disk, but whose distance to the class of disks is of order $\varepsilon$.

\begin{proof}[Proof of Theorem~\ref{thm:GGRT-sharp}]
Fix
\[
\lambda_0:=1-e^{-\pi},
\]
which is the eigenvalue of the localization operator of the unit disk $D_1$
associated with the ground state $h_0$. For $\varepsilon$ small, define
\[
P_\varepsilon(w):=1+\varepsilon w^2.
\]
By Theorem~\ref{thm:perturb}, applied with \(\lambda=\lambda_0\) and
\(P_\varepsilon(w)=1+\varepsilon w^2\), for all sufficiently small
\(|\varepsilon|\) there exists a domain \(U_\varepsilon\) such that
\begin{equation}\label{eq:functiona-P-eps-new}
\int_{U_\varepsilon} P_\varepsilon(z)e^{\pi \overline zw}e^{-\pi|z|^2}\,dA(z)
=
\lambda_0 P_\varepsilon(w),
\qquad w\in\C.
\end{equation}
Moreover, by construction, $U_\varepsilon$ is a $C^2$ perturbation of the unit
disk. Hence, after possibly shrinking the range of admissible $\varepsilon$,
its boundary can be written as
\begin{equation}\label{eq:def-u-eps-new}
\partial U_\varepsilon
=
\{(1+\varphi_\varepsilon(y))y:\ y\in \mathbb S^1\},
\end{equation}
where
\[
\varphi_\varepsilon\in C^2(\mathbb S^1),
\qquad
\|\varphi_\varepsilon\|_{C^2(\mathbb S^1)}\lesssim |\varepsilon|.
\]
\[
\varepsilon\mapsto \varphi_\varepsilon
\]
may in fact be chosen \(C^1\) as a map from a neighborhood of \(0\) into
\(C^1(\mathbb S^1)\): this follows from the proof of
Theorem~\ref{thm:perturb}, since both the Isakov step and the weighted
quadrature-domain step are obtained by implicit-function-theorem arguments with
\(C^1\) dependence on the data, and the polynomial
\[
P_\varepsilon(z)=1+\varepsilon z^2
\]
depends smoothly on \(\varepsilon\). We therefore write
\[
\dot\varphi_0(y):=
\left.\partial_\varepsilon\right|_{\varepsilon=0}\varphi_\varepsilon(y).
\]

\medskip
\noindent
{\bf Step 1: first variation of the eigenvalue identity.} We differentiate \eqref{eq:functiona-P-eps-new} at $\varepsilon=0$.
Since \(P_\varepsilon\to 1\) in \(C^\infty\) and
\(\varepsilon\mapsto \varphi_\varepsilon\) is \(C^1\) into \(C^1(\mathbb S^1)\),
all terms in \eqref{eq:functiona-P-eps-new} are differentiable at
\(\varepsilon=0\). Since $U_\varepsilon$ is given by a normal graph over $\mathbb S^1$, Reynolds'
transport formula yields, for every smooth test function $H$,
\[
\left.\frac{d}{d\varepsilon}\right|_{\varepsilon=0}
\int_{U_\varepsilon} H(z)\,dA(z)
=
\int_{\mathbb S^1} \dot\varphi_0(z)H(z)\,d\mathcal H^1(z).
\]
Applying this to
\[
H(z)=e^{\pi \overline zw}e^{-\pi|z|^2},
\]
and differentiating the factor $P_\varepsilon(z)=1+\varepsilon z^2$, we obtain
\begin{equation}\label{eq:first-variation-identity-new}
\int_{\mathbb S^1}
\dot\varphi_0(z)e^{\pi \overline zw}e^{-\pi|z|^2}\,d\mathcal H^1(z)
+
\int_{D_1} z^2 e^{\pi \overline zw}e^{-\pi|z|^2}\,dA(z)
=
\lambda_0 w^2.
\end{equation}
Since $D_1$ is radial, the monomial $z^2$ is an eigenfunction of the
localization operator of $D_1$; therefore there exists $\lambda_2>0$ such that
\[
\int_{D_1} z^2 e^{\pi \overline zw}e^{-\pi|z|^2}\,dA(z)=\lambda_2 w^2.
\]
Hence \eqref{eq:first-variation-identity-new} becomes
\begin{equation}\label{eq:boundary-identity-new}
\int_{\mathbb S^1}
\dot\varphi_0(z)e^{\pi \overline zw}e^{-\pi|z|^2}\,d\mathcal H^1(z)
=
(\lambda_0-\lambda_2)w^2.
\end{equation}
In fact, \(\lambda_2\) can be computed explicitly. Expanding
\(e^{\pi \overline zw}\) into its power series and using orthogonality on the
disk,
\[
\lambda_2
=
\frac{\pi^2}{2}\int_{D_1}|z|^4e^{-\pi|z|^2}\,dA(z)
=
1-\Bigl(1+\pi+\frac{\pi^2}{2}\Bigr)e^{-\pi}.
\]
Therefore
\[
\lambda_0-\lambda_2
=
\pi\Bigl(1+\frac{\pi}{2}\Bigr)e^{-\pi}>0,
\]
so the coefficient of \(w^2\) on the right-hand side is indeed nonzero.
In particular, since $|z|=1$ on $\mathbb S^1$, we may rewrite this as
\begin{equation}\label{eq:boundary-identity-circle-new}
e^{-\pi}
\int_{\mathbb S^1}
\dot\varphi_0(z)e^{\pi \overline zw}\,d\mathcal H^1(z)
=
(\lambda_0-\lambda_2)w^2.
\end{equation}

Setting $w=0$ gives
\begin{equation}\label{eq:cancellation-f-0-new}
\int_{\mathbb S^1}\dot\varphi_0\,d\mathcal H^1=0.
\end{equation}

\medskip
\noindent
{\bf Step 2: identification of the leading mode.} Expanding
\[
e^{\pi \overline zw}
=
\sum_{n\geq 0}\frac{\pi^n}{n!}\overline z^{\,n}w^n,
\]
and comparing coefficients in \eqref{eq:boundary-identity-circle-new}, we find
\[
\int_{\mathbb S^1}\dot\varphi_0(z)\overline z^{\,n}\,d\mathcal H^1(z)=0
\qquad \text{for all } n\neq 2,
\]
whereas the coefficient corresponding to $n=2$ is nonzero. Since
$\dot\varphi_0$ is real-valued, it follows that
\begin{equation}\label{eq:mode-decomposition-new}
\dot\varphi_0(e^{it})
=
a\cos(2t)+b\sin(2t)
\end{equation}
for some pair $(a,b)\neq (0,0)$. In particular,
\begin{equation}\label{eq:l1-nontrivial-new}
\int_{\mathbb S^1} |\dot\varphi_0|\,d\mathcal H^1>0.
\end{equation}

\medskip
\noindent
{\bf Step 3: normalization of the area.} The area of $U_\varepsilon$ is given by
\begin{equation}\label{eq:area-u-eps-new}
|U_\varepsilon|
=
\frac12\int_{\mathbb S^1} (1+\varphi_\varepsilon(y))^2\,d\mathcal H^1(y).
\end{equation}
Since
\[
\varphi_\varepsilon
=
\varepsilon \dot\varphi_0 + O(\varepsilon^2)
\quad \text{in } C^1(\mathbb S^1),
\]
and \eqref{eq:cancellation-f-0-new} gives the vanishing of the first-order
term, we obtain from \eqref{eq:area-u-eps-new}
\begin{equation}\label{eq:measure-close-u-eps-new}
|U_\varepsilon|=\pi+O(\varepsilon^2).
\end{equation}

Let
\[
r_\varepsilon:=\sqrt{\frac{\pi}{|U_\varepsilon|}},
\qquad
\widetilde U_\varepsilon:=r_\varepsilon U_\varepsilon.
\]
Then $|\widetilde U_\varepsilon|=\pi$, and by
\eqref{eq:measure-close-u-eps-new},
\[
r_\varepsilon=1+O(\varepsilon^2).
\]
Hence $\widetilde U_\varepsilon$ still admits a normal graph representation
\[
\partial \widetilde U_\varepsilon
=
\{(1+\widetilde\varphi_\varepsilon(y))y:\ y\in \mathbb S^1\},
\]
with
\[
\widetilde\varphi_\varepsilon
=
\varepsilon \dot\varphi_0 + O(\varepsilon^2)
\quad \text{in } C^1(\mathbb S^1).
\]
Thus the area normalization does not alter the first-order deformation.

\medskip
\noindent
{\bf Step 4: distance from the centered disk.} Since both $\widetilde U_\varepsilon$ and $D_1$ are star-shaped with respect to
the origin, the symmetric difference may be computed in polar coordinates:
\[
|\widetilde U_\varepsilon \triangle D_1|
=
\frac12\int_{\mathbb S^1}
\left|(1+\widetilde\varphi_\varepsilon(y))^2-1\right|
\,d\mathcal H^1(y).
\]
Using
\[
(1+\widetilde\varphi_\varepsilon)^2-1
=
2\widetilde\varphi_\varepsilon+\widetilde\varphi_\varepsilon^2
=
2\varepsilon \dot\varphi_0 + O(\varepsilon^2)
\quad \text{in } L^\infty(\mathbb S^1),
\]
we obtain
\[
|\widetilde U_\varepsilon \triangle D_1|
=
\frac12\int_{\mathbb S^1}
\left|2\varepsilon \dot\varphi_0 + O(\varepsilon^2)\right|
\,d\mathcal H^1.
\]
Since \eqref{eq:l1-nontrivial-new} gives
\[
\int_{\mathbb S^1} |\dot\varphi_0|\,d\mathcal H^1>0,
\]
it follows that there exist constants $c,C>0$ such that
\begin{equation}\label{eq:distance-centered-disk-new}
c|\varepsilon|
\leq
|\widetilde U_\varepsilon \triangle D_1|
\leq
C|\varepsilon|
\end{equation}
for all sufficiently small $\varepsilon$.

\medskip
\noindent
{\bf Step 5: comparison with arbitrary disks.} Let $D$ be any disk of area $\pi$. Then $D=D_1+a$ for some $a\in \C$. We first consider disks whose centers are not too close to the origin.
Since
\[
|D_1\triangle (D_1+a)|
\gtrsim \min\{|a|,1\},
\]
the triangle inequality and \eqref{eq:distance-centered-disk-new} imply that
there exists $K\gg 1$ such that, if $|a|\geq K|\varepsilon|$, then
\begin{equation}\label{eq:far-disks-new}
|\widetilde U_\varepsilon \triangle D|
\geq
|D_1\triangle D|-|\widetilde U_\varepsilon\triangle D_1|
\gtrsim
|a|-C|\varepsilon|
\gtrsim
|\varepsilon|.
\end{equation}

It remains to consider disks whose centers satisfy $|a|\leq K|\varepsilon|$.
Write
\[
a=\eta \varepsilon e^{i\theta},
\qquad
|\eta|\leq K,
\quad
\theta\in [0,2\pi).
\]
We claim that, for $\varepsilon$ sufficiently small, the translated domain
\[
\widetilde U_\varepsilon-a
\]
is still a normal graph over $\mathbb S^1$:
\[
\partial(\widetilde U_\varepsilon-a)
=
\{(1+\psi_\varepsilon^{\theta,\eta}(y))y:\ y\in \mathbb S^1\},
\]
with
\[
\psi_\varepsilon^{\theta,\eta}\in C^1(\mathbb S^1),
\qquad
\|\psi_\varepsilon^{\theta,\eta}\|_{C^1}\lesssim |\varepsilon|,
\]
uniformly in $(\theta,\eta)$ with $|\eta|\leq K$. This follows because
$\widetilde U_\varepsilon$ is already an $O(\varepsilon)$-$C^1$ perturbation of
$D_1$, and the translation vector $a$ is also of order $\varepsilon$. We now compute the first-order term of $\psi_\varepsilon^{\theta,\eta}$.
A translation by $a$ contributes, to first order in $\varepsilon$, the normal
displacement
\[
-\eta \Re(e^{-i\theta}y).
\]
Indeed, for \(y\in \mathbb S^1\),
\[
|y-a|
=
\bigl(1-2\Re(\overline y a)+|a|^2\bigr)^{1/2}
=
1-\Re(\overline y a)+O(|a|^2),
\]
uniformly in \(y\), and here \(a=\eta\varepsilon e^{i\theta}\).
Therefore
\begin{equation}\label{eq:psi-expansion-new}
\psi_\varepsilon^{\theta,\eta}(y)
=
\varepsilon\Big(\dot\varphi_0(y)-\eta \Re(e^{-i\theta}y)\Big)
+
O(\varepsilon^2)
\quad \text{in } C^1(\mathbb S^1),
\end{equation}
uniformly for $|\eta|\leq K$.

Define
\[
g^{\theta,\eta}(y)
:=
\dot\varphi_0(y)-\eta \Re(e^{-i\theta}y).
\]
Using \eqref{eq:mode-decomposition-new}, we see that $g^{\theta,\eta}$ belongs
to the finite-dimensional space spanned by the first and second harmonics.
More precisely,
\[
g^{\theta,\eta}(e^{it})
=
a\cos(2t)+b\sin(2t)-\eta \cos(t-\theta).
\]
The second-harmonic part is independent of $(\theta,\eta)$ and nonzero, because
$(a,b)\neq (0,0)$. Hence the family
\[
\{g^{\theta,\eta}:\ \theta\in[0,2\pi),\ |\eta|\leq K\}
\]
is a compact subset of a finite-dimensional vector space which does not contain
the zero function. It follows that there exists $\delta_K>0$ such that
\begin{equation}\label{eq:uniform-l1-lower-new}
\int_{\mathbb S^1}|g^{\theta,\eta}(y)|\,d\mathcal H^1(y)\geq \delta_K
\qquad
\text{for all }\theta\in[0,2\pi),\ |\eta|\leq K.
\end{equation}

Now,
\[
|\widetilde U_\varepsilon \triangle D|
=
|(\widetilde U_\varepsilon-a)\triangle D_1|.
\]
Since $\widetilde U_\varepsilon-a$ is star-shaped with radial graph
$1+\psi_\varepsilon^{\theta,\eta}$, the same polar-coordinate computation as in
Step~4 gives
\[
|(\widetilde U_\varepsilon-a)\triangle D_1|
=
\frac12\int_{\mathbb S^1}
\left|(1+\psi_\varepsilon^{\theta,\eta}(y))^2-1\right|
\,d\mathcal H^1(y).
\]
Using \eqref{eq:psi-expansion-new},
\[
(1+\psi_\varepsilon^{\theta,\eta})^2-1
=
2\varepsilon g^{\theta,\eta}+O(\varepsilon^2)
\quad \text{in } L^\infty(\mathbb S^1),
\]
uniformly for $|\eta|\leq K$. Thus, by \eqref{eq:uniform-l1-lower-new},
there exists $c_K>0$ such that
\begin{equation}\label{eq:near-disks-new}
|\widetilde U_\varepsilon \triangle D|
=
|(\widetilde U_\varepsilon-a)\triangle D_1|
\geq
c_K |\varepsilon|
\end{equation}
whenever $|a|\leq K|\varepsilon|$. Combining \eqref{eq:far-disks-new} and \eqref{eq:near-disks-new}, we conclude
that
\[
\inf_{\substack{D \text{ disk}\\ |D|=\pi}}
|\widetilde U_\varepsilon \triangle D|
\gtrsim |\varepsilon|.
\]
Since \(r_\varepsilon=1+O(\varepsilon^2)\), the dilation from \(U_\varepsilon\)
to \(\widetilde U_\varepsilon\) changes the set only by
\[
|\widetilde U_\varepsilon\triangle U_\varepsilon|=O(\varepsilon^2),
\]
and therefore does not affect the first-order distance estimate. Thus the
area-normalized family \(\widetilde U_\varepsilon\) remains at
symmetric-difference distance \(\gtrsim |\varepsilon|\) from the class of
disks, while its spectral data differ from those of the disk only at the same
perturbative scale as before the normalization. This shows that the estimate in
Corollary~1.4 of \cite{Gomez-Guerra-Ramos-Tilli} is sharp.

\end{proof}

\section{Local maximizers for the Faber--Krahn problem in the Gaussian STFT setting}

We work throughout in the Bargmann--Fock representation. Thus, if \(F\in \mathcal F^2(\C)\)
is the Bargmann transform of \(f\in L^2(\R)\), then
\[
|\mathcal V_\varphi f(z)|^2 = |F(z)|^2 e^{-\pi |z|^2},
\]
up to the normalization already fixed in the previous sections. Accordingly, for
\(F\in \mathcal F^2(\C)\) we set
\[
u_F(z):=|F(z)|^2 e^{-\pi |z|^2}.
\]

For a measurable set \(\Omega\subset \C\), let
\[
\lambda_{1,\varphi}(\Omega)
:=
\sup_{\|f\|_2=1}\int_\Omega |\mathcal V_\varphi f(z)|^2\,dz.
\]
We say that \(\Omega\) is a \emph{local maximizer} for \(\lambda_{1,\varphi}\) under the
measure constraint if there exists \(\delta_0>0\) such that, whenever
\(E\subset \C\) is measurable, \(|E|=|\Omega|\), and
\[
\int_\C |\mathbf 1_\Omega-\mathbf 1_E|\,e^{-\pi |z|^2}\,dz<\delta_0,
\]
one has
\[
\lambda_{1,\varphi}(\Omega)\ge \lambda_{1,\varphi}(E).
\]

We now prove the local rigidity statement announced in the introduction.

\begin{proof}[Proof of the local-maximizer theorem]
Let \(\Sigma\subset \C\) be a local maximizer for \(\lambda_{1,\varphi}\) under the
measure constraint, and let \(F\in \mathcal F^2(\C)\) be a normalized first
eigenfunction for the localization operator associated with \(\Sigma\).
Equivalently,
\[
\int_\Sigma H(z)\overline{F(z)}\,d\mu(z)
=
\lambda_{1,\varphi}(\Sigma)\int_\C H(z)\overline{F(z)}\,d\mu(z)
\]
for every \(H\in \mathcal F^2(\C)\), where
\[
d\mu(z):=e^{-\pi |z|^2}\,dA(z).
\]

\medskip
\noindent
{\bf Step 1: \(\Sigma\) is a superlevel set of \(u_F\).}

We claim that there exists \(t_0\ge 0\) such that, up to sets of measure zero,
\[
\Sigma=\{u_F>t_0\}.
\]

Indeed, suppose this fails. Then there exists \(t_0\) such that
\[
A:=\{u_F>t_0\}
\qquad\text{satisfies}\qquad
|A|=|\Sigma|,
\]
but \(A\) and \(\Sigma\) do not agree a.e. Since the weighted symmetric-difference
distance is absolutely continuous with respect to Lebesgue measure on bounded
sets, we may choose measurable sets
\[
S\subset A\setminus \Sigma,
\qquad
T\subset \Sigma\setminus A,
\qquad
|S|=|T|>0,
\]
with
\[
\int_\C |\mathbf 1_S+\mathbf 1_T|e^{-\pi |z|^2}\,dz<\delta_0.
\]
Define
\[
\Sigma':=(\Sigma\setminus T)\cup S.
\]
Then \(|\Sigma'|=|\Sigma|\) and
\[
\int_\C |\mathbf 1_\Sigma-\mathbf 1_{\Sigma'}|e^{-\pi |z|^2}\,dz<\delta_0.
\]
On the other hand, by construction,
\[
u_F(z)>t_0 \quad \text{on } S,
\qquad
u_F(z)\le t_0 \quad \text{on } T,
\]
hence
\[
\int_{\Sigma'} u_F\,dA
=
\int_\Sigma u_F\,dA + \int_S u_F\,dA-\int_T u_F\,dA
>
\int_\Sigma u_F\,dA.
\]
Since \(F\) is admissible for the variational problem defining
\(\lambda_{1,\varphi}(\Sigma')\), this implies
\[
\lambda_{1,\varphi}(\Sigma')
\ge \int_{\Sigma'} u_F\,dA
>
\int_\Sigma u_F\,dA
=
\lambda_{1,\varphi}(\Sigma),
\]
contradicting the local maximality of \(\Sigma\). This proves the claim.

In particular, \(u_F\) is constant on \(\partial \Sigma\), in the sense that
there exists \(c_\Sigma>0\) with
\begin{equation}\label{eq:boundary-level-set-new}
u_F(z)=c_\Sigma
\qquad \text{for } z\in \partial \Sigma.
\end{equation}

\medskip
\noindent
{\bf Step 2: simplicity of the first eigenvalue.}

We claim that the first eigenvalue \(\lambda_{1,\varphi}(\Sigma)\) is simple.

Suppose instead that \(G\in \mathcal F^2(\C)\) is another first eigenfunction,
linearly independent of \(F\). By the previous step, \(\Sigma\) is also a
superlevel set of \(u_G\), hence there exists \(c'_\Sigma>0\) such that
\[
u_G(z)=c'_\Sigma
\qquad\text{for } z\in \partial \Sigma.
\]
Equivalently,
\[
|F(z)|^2 e^{-\pi |z|^2}=c_\Sigma,
\qquad
|G(z)|^2 e^{-\pi |z|^2}=c'_\Sigma,
\qquad z\in \partial \Sigma.
\]
Dividing, we obtain
\[
\frac{|F(z)|^2}{|G(z)|^2}=\frac{c_\Sigma}{c'_\Sigma}
\qquad \text{for } z\in \partial \Sigma.
\]

Now we use the standard fact that first eigenfunctions may be chosen without zeros
in \(\Sigma\): by Step 1, \(\Sigma\) is a superlevel set of \(u_F\), hence \(u_F>0\) in the
interior of \(\Sigma\). Since \(u_F=|F|^2e^{-\pi |z|^2}\) and the Gaussian factor
never vanishes, it follows that \(F\) has no zeros in \(\Sigma\). The same
argument applies to \(G\). Therefore \(F/G\) is holomorphic and nonvanishing in
\(\Sigma\), and the function
\[
\log \left|\frac{F}{G}\right|
=
\frac12 \log \frac{|F|^2}{|G|^2}
=
\frac12 \log \frac{u_F}{u_G}
\]
is harmonic in \(\Sigma\). Since the Gaussian factors cancel, this function is
constant on \(\partial \Sigma\). By the maximum principle, it is constant on each
connected component of \(\Sigma\). Hence
\[
\left|\frac{F}{G}\right| \equiv \text{constant}
\qquad \text{on each connected component of } \Sigma.
\]
Since \(F/G\) is holomorphic and has constant modulus on each connected component,
it is constant there. Thus
\[
F=\gamma_j G
\qquad \text{on each connected component of } \Sigma.
\]
By analyticity, the same relation holds on all of \(\C\). By unique continuation, $F$ and $G$
are linearly dependent, a contradiction. This proves simplicity.

\medskip
\noindent
{\bf Step 3: a first-variation identity.}

Since \(\Sigma\) is a level set of the smooth function \(u_F\), its boundary is
Lipschitz (indeed \(C^1\) away from the critical set of \(u_F\), which is enough
for the divergence argument below). Let \(\psi\) be holomorphic on an open
neighborhood \(V\) of \(\overline \Sigma\), and write
\[
\psi=u+iv
\]
with \(u,v\) real-valued. Consider the vector field
\[
X_\psi:=(u,-v).
\]
Since \(\psi\) is holomorphic, the Cauchy--Riemann equations give
\[
u_x=v_y,
\qquad
u_y=-v_x,
\]
hence
\[
\operatorname{div}X_\psi=u_x-v_y=0
\qquad \text{in }V.
\]
Because \(u_F\) is constant on \(\partial \Sigma\), we obtain
\[
0=\int_{\partial \Sigma} u_F\, X_\psi\cdot \nu\, d\mathcal H^1
 =\int_\Sigma \operatorname{div}(u_F X_\psi)\,dA
 =\int_\Sigma \langle \nabla u_F, X_\psi\rangle\,dA.
\]
Since \(u_F=|F|^2e^{-\pi |z|^2}\), this becomes
\begin{equation}\label{eq:vector-field-identity-real-new}
0=
\int_\Sigma
\left(
\langle \nabla |F|^2,X_\psi\rangle
-
2\pi |F|^2 \langle z,X_\psi\rangle
\right)
\,d\mu.
\end{equation}

Let
\[
A(z):=F'(z)\overline{F(z)}.
\]
Since
\[
\partial_x |F|^2 = A+\overline A = 2\Re A,
\qquad
\partial_y |F|^2 = iA-i\overline A = -2\Im A,
\]
a direct computation gives
\[
\langle \nabla |F|^2,X_\psi\rangle
=
2\,\Re\!\big(\psi\,F'\overline{F}\big).
\]
Likewise, writing \(z=x+iy\), we have
\[
\langle z,X_\psi\rangle
=
\Re(z\psi).
\]
Substituting these expressions into
\eqref{eq:vector-field-identity-real-new}, we obtain
\[
0=
\Re\!\int_\Sigma
\psi(z)\Big(F'(z)\overline{F(z)}-\pi z\,|F(z)|^2\Big)\,d\mu(z).
\]
Applying the same identity with \(i\psi\) in place of \(\psi\), we also obtain
\[
0=
\Im\!\int_\Sigma
\psi(z)\Big(F'(z)\overline{F(z)}-\pi z\,|F(z)|^2\Big)\,d\mu(z).
\]
Hence
\begin{equation}\label{eq:orthogonality-main-new}
0=
\int_\Sigma
\psi(z)\Big(F'(z)\overline{F(z)}-\pi z\,|F(z)|^2\Big)\,d\mu(z).
\end{equation}
Taking complex conjugates, we may equivalently write
\[
0=
\int_\Sigma
\psi(z)\,F(z)\Big(\overline{F'(z)}-\pi z\,\overline{F(z)}\Big)\,d\mu(z)
\]
for every holomorphic \(\psi\) on a neighborhood of \(\overline \Sigma\).

\medskip
\noindent
{\bf Step 4: \(F'\) is again a first eigenfunction.}

We now show that \(F'\) is an eigenfunction associated with the same first
eigenvalue \(\lambda_{1,\varphi}(\Sigma)\).

For every \(w\in \C\), set
\[
K_w(z):=e^{\pi \overline wz}.
\]
By Step 1, after replacing \(\Sigma\) by the corresponding superlevel-set
representative if necessary, we may suppose that
\[
\Sigma=\{u_F>c_\Sigma\}.
\]
In particular, \(u_F\ge c_\Sigma>0\) on \(\overline \Sigma\). Since
\[
u_F(z)=|F(z)|^2e^{-\pi |z|^2},
\]
it follows that \(F\) has no zeros on \(\overline \Sigma\). By continuity and the
compactness of \(\overline \Sigma\), there exists an open neighborhood
\(V\supset \overline \Sigma\) on which \(F\) has no zeros. Hence, for every
\(w\in \C\), the function
\[
\psi_w(z):=\frac{K_w(z)}{F(z)}
\]
is holomorphic on \(V\). Applying the identity from Step 3 with \(\psi=\psi_w\),
we get
\[
0=
\int_\Sigma
e^{\pi \overline wz}
\Big(\overline{F'(z)}-\pi z\,\overline{F(z)}\Big)\,d\mu(z).
\]
Hence
\begin{equation}\label{eq:key-kernel-identity-new}
\int_\Sigma
e^{\pi \overline wz}\overline{F'(z)}\,d\mu(z)
=
\int_\Sigma
\pi z\,e^{\pi \overline wz}\overline{F(z)}\,d\mu(z).
\end{equation}

On the other hand, from the eigenvalue equation for \(F\), taking
\[
G_0(z):=z e^{\pi \overline wz},
\]
we obtain
\[
\int_\Sigma z e^{\pi \overline wz}\overline{F(z)}\,d\mu(z)
=
\lambda_{1,\varphi}(\Sigma)\int_\C z e^{\pi \overline wz}\overline{F(z)}\,d\mu(z).
\]
Using the reproducing property of the Fock space,
\[
\int_\C e^{\pi \overline wz}\overline{F(z)}\,d\mu(z)=\overline{F(w)},
\]
and differentiating with respect to \(\overline w\), we get
\[
\pi\int_\C z e^{\pi \overline wz}\overline{F(z)}\,d\mu(z)=\overline{F'(w)}.
\]
Therefore
\[
\int_\Sigma \pi z e^{\pi \overline wz}\overline{F(z)}\,d\mu(z)
=
\lambda_{1,\varphi}(\Sigma)\,\overline{F'(w)}.
\]
Combining this with \eqref{eq:key-kernel-identity-new}, we conclude that
\[
\int_\Sigma e^{\pi \overline wz}\overline{F'(z)}\,d\mu(z)
=
\lambda_{1,\varphi}(\Sigma)\,\overline{F'(w)},
\qquad w\in \C.
\]
This is exactly the eigenvalue equation for \(F'\). Thus \(F'\) is a first
eigenfunction as claimed.

By simplicity of the first eigenvalue, proved in Step 2, it follows that
\[
F'=\theta F
\]
for some constant \(\theta\in \C\). Solving this ODE gives
\[
F(z)=C e^{\theta z},
\]
for some nonzero constant \(C\).

\medskip
\noindent
{\bf Step 5: reduction to the constant eigenfunction.}

Up to translation of the set \(\Sigma\), we may absorb the exponential factor
\(e^{\theta z}\) and reduce to the case
\[
F\equiv 1.
\]
Indeed, in the Bargmann--Fock model, translating the underlying domain corresponds
to composing the eigenfunction with the appropriate Weyl factor; thus, after a
translation, we may normalize away the linear exponential term.

Therefore, after translation, the ground state \(h_0\) is an eigenfunction of the
localization operator of \(\Sigma\). Since \(\Sigma\) is a level set of \(u_F(z) = e^{-\pi |z|^2},\) it follows at once that $\Sigma$ is either an annulus or a disk. 

\medskip
\noindent
{\bf Step 6: exclusion of annuli and conclusion.}

It remains to show that among radial sets of prescribed measure, the only local
maximizer is the disk. If \(\Sigma\) is radial, write
\[
\widetilde \Sigma:=\{r>0:\ re^{it}\in \Sigma \text{ for some } t\}.
\]
Since \(\Sigma\) has Lipschitz boundary, \(\widetilde \Sigma\) is a finite union
of disjoint intervals:
\[
\widetilde \Sigma=\bigcup_j (a_j,b_j).
\]
Then
\[
\int_\Sigma e^{-\pi |z|^2}\,dz
=
2\pi \int_{\widetilde \Sigma} r e^{-\pi r^2}\,dr
=
\sum_j 2\pi \int_{a_j}^{b_j} r e^{-\pi r^2}\,dr
=
\sum_j \big(e^{-\pi a_j^2}-e^{-\pi b_j^2}\big).
\]
Under the constraint
\[
|\Sigma|
=
2\pi \int_{\widetilde \Sigma} r\,dr
=
\pi \sum_j (b_j^2-a_j^2)
\]
fixed, this quantity is maximized by placing all the radial mass as close as
possible to the origin, namely by taking
\[
\widetilde \Sigma=(0,R)
\]
for the unique \(R>0\) such that \(|D_R|=|\Sigma|\). Therefore the only radial
maximizer is the centered disk.

We conclude that every local maximizer of \(\lambda_{1,\varphi}\) under the
measure constraint is, up to translation, a disk.
\end{proof}

\section{Global maximizers for the Faber--Krahn problem in the Gaussian STFT setting}

In this section we prove existence of global maximizers for the Gaussian
time-frequency concentration problem under a measure constraint, and then combine
that existence result with the local rigidity theorem from the previous section
to recover the disk as the unique maximizer, up to translation.


Since the Bargmann transform is unitary, the variational problem may be written as
\[
M(s):=\sup_{|\Omega|=s}\lambda_{1,\varphi}(\Omega)
=
\sup_{\|F\|_{\mathcal F^2}=1} I_F(s).
\]



\begin{lemma}\label{lem:weyl-orthogonality-global}
Let \(F,G\in \mathcal F^2(\C)\), and let \(\{a_k\}\subset \C\) satisfy
\(|a_k|\to \infty\). Then
\[
T_{a_k}G \rightharpoonup 0
\qquad\text{weakly in }\mathcal F^2(\C),
\]
and in particular
\[
\langle F,T_{a_k}G\rangle_{\mathcal F^2}\to 0.
\]
Consequently,
\[
\|F+T_{a_k}G\|_{\mathcal F^2}^2
=
\|F\|_{\mathcal F^2}^2+\|G\|_{\mathcal F^2}^2+o(1).
\]
More generally, if \(|a_k-b_k|\to\infty\), then
\[
\langle T_{a_k}F,T_{b_k}G\rangle_{\mathcal F^2}\to 0.
\]
\end{lemma}

\begin{proof}
Using unitarity,
\[
\langle F,T_{a_k}G\rangle
=
\langle T_{-a_k}F,G\rangle.
\]
Now \(T_{-a_k}F\) is bounded in \(\mathcal F^2\), and
\eqref{eq:u-translation-global} shows that its associated densities drift to
infinity. Equivalently, for every compact \(K\subset \C\),
\[
\sup_{z\in K}|T_{-a_k}F(z)|e^{-\pi |z|^2/2}\to 0.
\]
Hence \(T_{-a_k}F\to 0\) locally uniformly, and therefore weakly in
\(\mathcal F^2\). This proves the first claim. The norm decoupling follows by
expanding the square. The general case reduces to the first by writing
\[
\langle T_{a_k}F,T_{b_k}G\rangle
=
\langle F,T_{b_k-a_k}G\rangle.
\]
\end{proof}

\begin{lemma}\label{lem:pointwise-decoupling-global}
Fix \(J\ge 1\), let \(G^{(1)},\dots,G^{(J)}\in \mathcal F^2(\C)\), and let
\(\{a_k^{(j)}\}_{k\ge 1}\subset \C\) for \(1\le j\le J\) satisfy
\[
|a_k^{(j)}-a_k^{(\ell)}|\to\infty
\qquad\text{whenever }j\neq \ell.
\]
For each \(k\), set
\[
S_k:=\sum_{j=1}^J T_{a_k^{(j)}}G^{(j)}.
\]
Then the following hold.
\begin{enumerate}
\item[(i)] For every \(R>0\) and every \(1\le j\le J\),
\[
\sup_{z\in D_R(a_k^{(j)})}
\big|u_{S_k}(z)-u_{T_{a_k^{(j)}}G^{(j)}}(z)\big|
\to 0
\qquad (k\to\infty).
\]

\item[(ii)] For every \(R>0\), if
\[
E_{k,R}:=\C\setminus \bigcup_{j=1}^J D_R(a_k^{(j)}),
\]
then
\[
\limsup_{k\to\infty}\sup_{z\in E_{k,R}}u_{S_k}(z)
\le
\left(
\sum_{j=1}^J
\sup_{\zeta\in \C\setminus D_R} u_{G^{(j)}}(\zeta)^{1/2}
\right)^2.
\]
In particular, given \(\varepsilon>0\), there exists \(R_\varepsilon>0\) such that
for every \(R\ge R_\varepsilon\),
\[
\sup_{z\in E_{k,R}}u_{S_k}(z)\le \varepsilon
\]
for all sufficiently large \(k\).
\end{enumerate}
\end{lemma}

\begin{proof}
We repeatedly use the identity
\[
u_{T_a F}(z)=u_F(z-a),
\qquad z,a\in \C,
\]
proved earlier.

\medskip
\noindent\emph{Proof of (i).}
Fix \(R>0\) and \(j\in\{1,\dots,J\}\). Write
\[
S_k=T_{a_k^{(j)}}G^{(j)}+H_{k,j},
\qquad
H_{k,j}:=\sum_{\ell\neq j} T_{a_k^{(\ell)}}G^{(\ell)}.
\]
Then for every \(z\in \C\),
\begin{align*}
\big|u_{S_k}(z)-u_{T_{a_k^{(j)}}G^{(j)}}(z)\big|
&=
e^{-\pi |z|^2}
\big||T_{a_k^{(j)}}G^{(j)}(z)+H_{k,j}(z)|^2-|T_{a_k^{(j)}}G^{(j)}(z)|^2\big| \\
&\le
2\,u_{T_{a_k^{(j)}}G^{(j)}}(z)^{1/2}u_{H_{k,j}}(z)^{1/2}
+u_{H_{k,j}}(z).
\end{align*}
By Lemma~\ref{lem:Fock-basic}(ii),
\[
u_{T_{a_k^{(j)}}G^{(j)}}(z)\le \|G^{(j)}\|_{\mathcal F^2}^2
\qquad \text{for all }z,k.
\]
Thus it is enough to show that
\[
\sup_{z\in D_R(a_k^{(j)})} u_{H_{k,j}}(z)\to 0.
\]

For \(z\in D_R(a_k^{(j)})\) and \(\ell\neq j\), we have
\[
z-a_k^{(\ell)}\in D_R(a_k^{(j)}-a_k^{(\ell)}).
\]
Since \(|a_k^{(j)}-a_k^{(\ell)}|\to\infty\), it follows that
\[
\operatorname{dist}\big(D_R(a_k^{(j)}-a_k^{(\ell)}),0\big)\to\infty.
\]
Because \(u_{G^{(\ell)}}\in C_0(\C)\), we obtain
\[
\sup_{z\in D_R(a_k^{(j)})} u_{T_{a_k^{(\ell)}}G^{(\ell)}}(z)
=
\sup_{\zeta\in D_R(a_k^{(j)}-a_k^{(\ell)})} u_{G^{(\ell)}}(\zeta)
\to 0.
\]
Now
\[
u_{H_{k,j}}(z)^{1/2}
\le
\sum_{\ell\neq j} u_{T_{a_k^{(\ell)}}G^{(\ell)}}(z)^{1/2},
\]
so the preceding convergence implies
\[
\sup_{z\in D_R(a_k^{(j)})} u_{H_{k,j}}(z)^{1/2}\to 0.
\]
This proves (i).

\medskip
\noindent\emph{Proof of (ii).}
Fix \(R>0\) and \(z\in E_{k,R}\). Then \(z\notin D_R(a_k^{(j)})\) for every
\(j\), hence \(|z-a_k^{(j)}|\ge R\). Therefore
\[
u_{T_{a_k^{(j)}}G^{(j)}}(z)
=
u_{G^{(j)}}(z-a_k^{(j)})
\le
\sup_{\zeta\in \C\setminus D_R}u_{G^{(j)}}(\zeta).
\]
Using again the elementary inequality
\[
u_{\sum_{j=1}^J F_j}(z)^{1/2}
\le \sum_{j=1}^J u_{F_j}(z)^{1/2},
\]
we get
\[
u_{S_k}(z)^{1/2}
\le
\sum_{j=1}^J u_{T_{a_k^{(j)}}G^{(j)}}(z)^{1/2}
\le
\sum_{j=1}^J \sup_{\zeta\in \C\setminus D_R}u_{G^{(j)}}(\zeta)^{1/2}.
\]
Squaring yields
\[
u_{S_k}(z)
\le
\left(
\sum_{j=1}^J \sup_{\zeta\in \C\setminus D_R}u_{G^{(j)}}(\zeta)^{1/2}
\right)^2.
\]
Taking the supremum over \(z\in E_{k,R}\) proves the displayed estimate in (ii).
Since each \(u_{G^{(j)}}\in C_0(\C)\), the right-hand side tends to \(0\) as
\(R\to\infty\). The final claim follows.
\end{proof}

\begin{lemma}\label{lem:density-perturbation-global}
Let \(\{S_k\}\subset \mathcal F^2(\C)\) be bounded and let \(\{R_k\}\subset \mathcal
F^2(\C)\) satisfy
\[
\|u_{R_k}\|_{L^\infty(\C)}\to 0.
\]
Then
\[
\|u_{S_k+R_k}-u_{S_k}\|_{L^\infty(\C)}\to 0.
\]
More precisely, if \(\sup_k \|S_k\|_{\mathcal F^2}\le M\), then
\[
\|u_{S_k+R_k}-u_{S_k}\|_{L^\infty(\C)}
\le
2M\,\|u_{R_k}\|_{L^\infty(\C)}^{1/2}
+\|u_{R_k}\|_{L^\infty(\C)}.
\]
\end{lemma}

\begin{proof}
For every \(z\in \C\),
\begin{align*}
\big|u_{S_k+R_k}(z)-u_{S_k}(z)\big|
&=
e^{-\pi |z|^2}\big||S_k(z)+R_k(z)|^2-|S_k(z)|^2\big| \\
&\le
2\,u_{S_k}(z)^{1/2}u_{R_k}(z)^{1/2}
+u_{R_k}(z).
\end{align*}
By Lemma~\ref{lem:Fock-basic}(ii),
\[
u_{S_k}(z)\le \|S_k\|_{\mathcal F^2}^2\le M^2.
\]
Hence
\[
\big|u_{S_k+R_k}(z)-u_{S_k}(z)\big|
\le
2M\,\|u_{R_k}\|_{L^\infty}^{1/2}
+\|u_{R_k}\|_{L^\infty}.
\]
Taking the supremum in \(z\) gives the claim.
\end{proof}

\begin{lemma}\label{lem:I-continuity-global}
Let \(F\in \mathcal F^2(\C)\). Then the function
\[
s\mapsto I_F(s)
\]
is finite, nondecreasing, and continuous on \([0,\infty)\).
\end{lemma}

\begin{proof}
By Lemma~\ref{lem:basic-properties-global}, \(u_F\in C_0(\C)\cap L^1(\C)\), so
\(I_F(s)\) is finite for every \(s\ge 0\). Monotonicity is immediate from the
definition.

To prove continuity, let \(u:=u_F\), and for \(t\ge 0\) define
\[
\mu_u(t):=\big|\{z\in \C:\ u(z)>t\}\big|.
\]
Since \(u\in C_0(\C)\), the bathtub principle shows that for each \(s\ge 0\),
\[
I_F(s)=\int_0^s u^*(\tau)\,d\tau,
\]
where \(u^*\) is the decreasing rearrangement of \(u\). In particular,
\(u^*\in L^1_{\mathrm{loc}}([0,\infty))\), and therefore \(I_F\) is continuous on
\([0,\infty)\).
\end{proof}

The next result is the profile decomposition adapted to the present setting.
Since it is the main compactness input of the section, we record it explicitly.

\begin{proposition}[Profile decomposition in the Fock space]\label{prop:profile-global}
Let \(\{F_k\}_{k\ge 1}\subset \mathcal F^2(\C)\) be a bounded sequence. Then,
after passing to a subsequence, there exist profiles
\[
G^{(j)}\in \mathcal F^2(\C),
\qquad j\ge 1,
\]
and sequences of points
\[
a_k^{(j)}\in \C,
\qquad j\ge 1,\ k\ge 1,
\]
such that:
\begin{enumerate}
    \item[(i)] for each \(j\neq \ell\),
    \[
    |a_k^{(j)}-a_k^{(\ell)}|\to \infty
    \qquad (k\to \infty);
    \]
    \item[(ii)] for every \(J\ge 1\), if
    \[
    R_k^{(J)}
    :=
    F_k-\sum_{j=1}^J T_{a_k^{(j)}}G^{(j)},
    \]
    then
    \[
    T_{-a_k^{(j)}}R_k^{(J)}\to 0
    \qquad\text{locally uniformly on }\C
    \]
    for every \(1\le j\le J\);
    \item[(iii)] for every \(J\ge 1\),
    \[
    \|F_k\|_{\mathcal F^2}^2
    =
    \sum_{j=1}^J \|G^{(j)}\|_{\mathcal F^2}^2
    +\|R_k^{(J)}\|_{\mathcal F^2}^2
    +o_k(1);
    \]
    \item[(iv)]
    \[
    \lim_{J\to\infty}\left(\limsup_{k\to\infty}
    \|u_{R_k^{(J)}}\|_{L^\infty(\C)}\right)=0.
    \]
\end{enumerate}
\end{proposition}

\begin{proof}
We argue by the standard iterative extraction scheme, but we spell out the points
used later.

If
\[
\limsup_{k\to\infty}\|u_{F_k}\|_{L^\infty(\C)}=0,
\]
there is nothing to prove: take all profiles equal to zero. Otherwise choose
\(a_k^{(1)}\in \C\) such that
\[
u_{F_k}(a_k^{(1)})\ge \frac12\|u_{F_k}\|_{L^\infty(\C)}.
\]
By \eqref{eq:I-translation-global}, replacing \(F_k\) by \(T_{-a_k^{(1)}}F_k\)
if necessary, we may suppose \(a_k^{(1)}=0\). Since \(\{F_k\}\) is bounded in
\(\mathcal F^2\), the reproducing-kernel bound implies that \(\{F_k\}\) is a
normal family. Passing to a subsequence, \(F_k\to G^{(1)}\) locally uniformly,
with \(G^{(1)}\neq 0\) by the choice of the centering. Set
\[
R_k^{(1)}:=F_k-G^{(1)}.
\]

If
\[
\limsup_{k\to\infty}\|u_{R_k^{(1)}}\|_{L^\infty(\C)}=0,
\]
we stop. Otherwise choose \(a_k^{(2)}\) so that
\[
u_{R_k^{(1)}}(a_k^{(2)})\ge \frac12\|u_{R_k^{(1)}}\|_{L^\infty(\C)}.
\]
The local uniform convergence \(R_k^{(1)}\to 0\) on compact sets forces
\(|a_k^{(2)}|\to\infty\). Then \(T_{-a_k^{(2)}}R_k^{(1)}\) is again bounded in
\(\mathcal F^2\), so after passing to a subsequence it converges locally
uniformly to some nonzero profile \(G^{(2)}\). We set
\[
R_k^{(2)}:=R_k^{(1)}-T_{a_k^{(2)}}G^{(2)}.
\]

Proceeding inductively, we obtain profiles \(G^{(j)}\), translations
\(a_k^{(j)}\), and remainders \(R_k^{(J)}\). Let us verify the properties in the
statement.

\medskip
\noindent\emph{Proof of (i).}
Suppose \(j<\ell\). At the \(\ell\)-th step, the sequence
\[
T_{-a_k^{(\ell)}}R_k^{(\ell-1)}
\]
converges locally uniformly to the nonzero profile \(G^{(\ell)}\). On the other
hand, for every \(m<\ell\), property (ii) already proved up to stage \(\ell-1\)
gives
\[
T_{-a_k^{(m)}}R_k^{(\ell-1)}\to 0
\qquad\text{locally uniformly on }\C.
\]
If \(|a_k^{(\ell)}-a_k^{(m)}|\) failed to tend to infinity along a subsequence,
then after passing to a further subsequence the difference
\(a_k^{(\ell)}-a_k^{(m)}\) would stay bounded, and therefore local uniform
convergence of \(T_{-a_k^{(m)}}R_k^{(\ell-1)}\) to \(0\) would force
\(T_{-a_k^{(\ell)}}R_k^{(\ell-1)}\to 0\) locally uniformly as well, a
contradiction. Thus
\[
|a_k^{(j)}-a_k^{(\ell)}|\to\infty
\qquad\text{for }j\neq \ell.
\]

\medskip
\noindent\emph{Proof of (ii).}
Fix \(J\ge 1\) and \(1\le j\le J\). By construction, at the \(j\)-th extraction
step we have
\[
T_{-a_k^{(j)}}R_k^{(j-1)}\to G^{(j)}
\qquad\text{locally uniformly on }\C.
\]
Subtracting the \(j\)-th extracted profile yields
\[
T_{-a_k^{(j)}}R_k^{(j)}
=
T_{-a_k^{(j)}}R_k^{(j-1)}-G^{(j)}
\to 0
\qquad\text{locally uniformly on }\C.
\]
Now let \(J>j\). Since
\[
R_k^{(J)}
=
R_k^{(j)}-\sum_{\ell=j+1}^J T_{a_k^{(\ell)}}G^{(\ell)},
\]
we get
\[
T_{-a_k^{(j)}}R_k^{(J)}
=
T_{-a_k^{(j)}}R_k^{(j)}
-\sum_{\ell=j+1}^J T_{a_k^{(\ell)}-a_k^{(j)}}G^{(\ell)}.
\]
The first term tends to \(0\) locally uniformly by the previous paragraph. For
each fixed \(\ell>j\), property (i) and the identity
\[
u_{T_b G}(z)=u_G(z-b)
\]
imply that \(T_{a_k^{(\ell)}-a_k^{(j)}}G^{(\ell)}\to 0\) locally uniformly, since
\(|a_k^{(\ell)}-a_k^{(j)}|\to\infty\). Thus
\[
T_{-a_k^{(j)}}R_k^{(J)}\to 0
\qquad\text{locally uniformly on }\C,
\]
which is exactly (ii).

\medskip
\noindent\emph{Proof of (iii).}
For fixed \(J\), write
\[
F_k=\sum_{j=1}^J T_{a_k^{(j)}}G^{(j)}+R_k^{(J)}.
\]
Expanding the \(\mathcal F^2\)-norm squared and using Lemma
~\ref{lem:weyl-orthogonality-global}, together with the pairwise divergence of the
translation parameters from (i), we obtain
\[
\|F_k\|_{\mathcal F^2}^2
=
\sum_{j=1}^J \|G^{(j)}\|_{\mathcal F^2}^2
+\|R_k^{(J)}\|_{\mathcal F^2}^2
+o_k(1).
\]

\medskip
\noindent\emph{Proof of (iv).}
Suppose, on the contrary, that (iv) fails. Then there exist \(\delta>0\), a
sequence \(J_m\to\infty\), and indices \(k_m\to\infty\) such that
\[
\|u_{R_{k_m}^{(J_m)}}\|_{L^\infty(\C)}\ge \delta
\qquad\text{for every }m.
\]
Choose \(b_m\in \C\) so that
\[
u_{R_{k_m}^{(J_m)}}(b_m)\ge \delta/2.
\]
We claim that for every fixed \(j\ge 1\),
\[
|b_m-a_{k_m}^{(j)}|\to\infty
\qquad (m\to\infty).
\]
Indeed, if this failed for some fixed \(j\), then after passing to a subsequence
we would have \(|b_m-a_{k_m}^{(j)}|\le C\). Since \(J_m\ge j\) for all large \(m\),
property (ii) gives
\[
T_{-a_{k_m}^{(j)}}R_{k_m}^{(J_m)}\to 0
\qquad\text{locally uniformly on }\C.
\]
Evaluating at \(z=b_m-a_{k_m}^{(j)}\), which remains in a fixed compact set, we
would obtain
\[
u_{R_{k_m}^{(J_m)}}(b_m)\to 0,
\]
contradicting \(u_{R_{k_m}^{(J_m)}}(b_m)\ge \delta/2\). Thus the claim holds.

Now the sequence
\[
T_{-b_m}R_{k_m}^{(J_m)}
\]
is bounded in \(\mathcal F^2(\C)\), so after passing to a subsequence it converges
locally uniformly to some \(H\in \mathcal F^2(\C)\). By construction,
\[
u_H(0)=\lim_{m\to\infty}u_{T_{-b_m}R_{k_m}^{(J_m)}}(0)
=
\lim_{m\to\infty}u_{R_{k_m}^{(J_m)}}(b_m)\ge \delta/2,
\]
so \(H\neq 0\).

Moreover, for every fixed \(j\),
\[
|b_m-a_{k_m}^{(j)}|\to\infty,
\]
so the new centers \(b_m\) are asymptotically disjoint from all previously
extracted profiles. Therefore \(H\) could be extracted as one more nontrivial
profile, contradicting the maximality of the decomposition procedure. This proves
(iv).
\end{proof}

We now turn to the global problem.

\begin{theorem}\label{thm:existence-global-maximizer}
For every \(s>0\), there exists \(F_*\in \mathcal F^2(\C)\) with
\(\|F_*\|_{\mathcal F^2}=1\) such that
\[
I_{F_*}(s)=M(s).
\]
Equivalently, there exists a measurable set \(\Omega_*\subset \C\) with
\(|\Omega_*|=s\) such that
\[
\lambda_{1,\varphi}(\Omega_*)=M(s).
\]
\end{theorem}

\begin{proof}
Let \(\{F_k\}\subset \mathcal F^2(\C)\) be a maximizing sequence:
\[
\|F_k\|_{\mathcal F^2}=1,
\qquad
I_{F_k}(s)\to M(s).
\]
By replacing \(F_k\) with a Weyl translate if necessary, and using the
translation invariance \eqref{eq:I-translation-global}, we may suppose that
\[
u_{F_k}(0)=\|u_{F_k}\|_{L^\infty(\C)}.
\]
Applying Proposition~\ref{prop:profile-global}, after passing to a subsequence we
may write
\[
F_k=\sum_{j=1}^J T_{a_k^{(j)}}G^{(j)}+R_k^{(J)}
\]
for every fixed \(J\), with the conclusions of that proposition.

Fix \(J\ge 1\) and let \(\Omega_k^{(J)}\) be a measurable set of measure \(s\)
such that
\[
I_{F_k}(s)=\int_{\Omega_k^{(J)}} u_{F_k}\,dA.
\]
By the bathtub principle, we may choose \(\Omega_k^{(J)}\) to be a superlevel set
of \(u_{F_k}\).

Let \(\varepsilon>0\). Choose \(R>0\) so large that
\[
\int_{\C\setminus D_R} u_{G^{(j)}}\,dA<\varepsilon
\qquad\text{for all } 1\le j\le J.
\]
Since the translation parameters are pairwise divergent, for \(k\) large enough
the disks
\[
D_R(a_k^{(1)}),\dots,D_R(a_k^{(J)})
\]
are pairwise disjoint. Write
\[
B_{k,j}:=D_R(a_k^{(j)}),
\qquad
E_{k,j}:=\Omega_k^{(J)}\cap B_{k,j},
\qquad
E_{k,0}:=\Omega_k^{(J)}\setminus \bigcup_{j=1}^J B_{k,j}.
\]
Set
\[
s_{k,j}:=|E_{k,j}|,
\qquad 0\le j\le J.
\]
Then
\[
\sum_{j=0}^J s_{k,j}=s.
\]

Fix
\[
S_k^{(J)}:=\sum_{j=1}^J T_{a_k^{(j)}}G^{(j)}.
\]
By Proposition~\ref{prop:profile-global}(iii), the sequence \(\{S_k^{(J)}\}_k\) is
bounded in \(\mathcal F^2(\C)\). By Proposition~\ref{prop:profile-global}(iv), if
\(J\) is taken sufficiently large then
\[
\limsup_{k\to\infty}\|u_{R_k^{(J)}}\|_{L^\infty(\C)}<\varepsilon^2.
\]
Lemma~\ref{lem:density-perturbation-global} therefore gives
\[
\|u_{F_k}-u_{S_k^{(J)}}\|_{L^\infty(\C)}\to 0
\qquad (k\to\infty),
\]
and in particular, after increasing \(k\) if necessary,
\begin{equation}\label{eq:density-approximation-global}
\|u_{F_k}-u_{S_k^{(J)}}\|_{L^\infty(\C)}<\varepsilon.
\end{equation}

Now let
\[
E_{k,0}
=
\Omega_k^{(J)}\setminus \bigcup_{j=1}^J B_{k,j}
\subset
\C\setminus \bigcup_{j=1}^J D_R(a_k^{(j)}).
\]
Lemma~\ref{lem:pointwise-decoupling-global}(i) gives, for each fixed
\(j\in\{1,\dots,J\}\),
\[
\sup_{z\in B_{k,j}}
\big|u_{S_k^{(J)}}(z)-u_{T_{a_k^{(j)}}G^{(j)}}(z)\big|\to 0
\qquad (k\to\infty),
\]
while Lemma~\ref{lem:pointwise-decoupling-global}(ii) yields, after taking \(R\)
sufficiently large and then \(k\) sufficiently large,
\[
\sup_{z\in E_{k,0}}u_{S_k^{(J)}}(z)\le \varepsilon.
\]
Combining these bounds with \eqref{eq:density-approximation-global}, we conclude
that for all sufficiently large \(k\),
\begin{enumerate}
\item[(a)] on each \(B_{k,j}\),
\[
u_{F_k}(z)\le u_{T_{a_k^{(j)}}G^{(j)}}(z)+2\varepsilon;
\]
\item[(b)] on \(E_{k,0}\),
\[
u_{F_k}(z)\le 2\varepsilon.
\]
\end{enumerate}

Therefore
\begin{align*}
I_{F_k}(s)
&= 
\int_{\Omega_k^{(J)}} u_{F_k}\,dA \le
\sum_{j=1}^J \int_{E_{k,j}} u_{T_{a_k^{(j)}}G^{(j)}}\,dA
+
+2\varepsilon \sum_{j=0}^J |E_{k,j}| \\
&=
\sum_{j=1}^J \int_{E_{k,j}-a_k^{(j)}} u_{G^{(j)}}\,dA
+
+2\varepsilon s \le
\sum_{j=1}^J I_{G^{(j)}}(s_{k,j})
+2\varepsilon s.
\end{align*}
Passing to a further subsequence if necessary, we may suppose
\[
s_{k,j}\to \sigma_j\in [0,s]
\qquad (k\to\infty)
\]
for each \(1\le j\le J\). By continuity of \(I_{G^{(j)}}(\cdot)\) in the measure
parameter, given by Lemma~\ref{lem:I-continuity-global}, we obtain
\[
M(s)\le \sum_{j=1}^J I_{G^{(j)}}(\sigma_j)+2\varepsilon s.
\]
Since \(\varepsilon>0\) is arbitrary,
\begin{equation}\label{eq:M-upper-profiles-global}
M(s)\le \sum_{j=1}^J I_{G^{(j)}}(\sigma_j).
\end{equation}

Now use Lemma~\ref{lem:basic-properties-global}. For each \(j\),
\[
I_{G^{(j)}}(\sigma_j)
=
\|G^{(j)}\|_{\mathcal F^2}^2
I_{\widehat G^{(j)}}(\sigma_j),
\qquad
\widehat G^{(j)}:=\frac{G^{(j)}}{\|G^{(j)}\|_{\mathcal F^2}}
\]
if \(G^{(j)}\neq 0\), and
\[
I_{\widehat G^{(j)}}(\sigma_j)\le M(s)
\]
because \(\sigma_j\le s\) and \(I_{\widehat G^{(j)}}(\cdot)\) is nondecreasing. Hence
\[
I_{G^{(j)}}(\sigma_j)\le \|G^{(j)}\|_{\mathcal F^2}^2\, M(s).
\]
Summing and using Proposition~\ref{prop:profile-global}(iii),
\[
\sum_{j=1}^J I_{G^{(j)}}(\sigma_j)
\le
M(s)\sum_{j=1}^J \|G^{(j)}\|_{\mathcal F^2}^2
\le M(s).
\]
Together with \eqref{eq:M-upper-profiles-global}, this gives
\[
M(s)\le \sum_{j=1}^J I_{G^{(j)}}(\sigma_j)\le M(s)
\]
for every \(J\). Therefore equality holds throughout, and in particular
\[
\sum_{j\ge 1}\|G^{(j)}\|_{\mathcal F^2}^2=1.
\]
Moreover, for every \(j\) with \(G^{(j)}\neq 0\),
\[
I_{\widehat G^{(j)}}(s)\ge I_{\widehat G^{(j)}}(\sigma_j)=M(s).
\]
Since \(M(s)\) is the supremum over all unit-norm functions, it follows that
\[
I_{\widehat G^{(j)}}(s)=M(s).
\]
Thus any nonzero normalized profile is a maximizer. Taking one such profile and
denoting it by \(F_*\), we obtain
\[
\|F_*\|_{\mathcal F^2}=1,
\qquad
I_{F_*}(s)=M(s).
\]

Finally, let \(\Omega_*\) be a superlevel set of \(u_{F_*}\) with \(|\Omega_*|=s\)
and
\[
\int_{\Omega_*}u_{F_*}\,dA=I_{F_*}(s)=M(s).
\]
Then
\[
\lambda_{1,\varphi}(\Omega_*)\ge \int_{\Omega_*}u_{F_*}\,dA=M(s),
\]
while by definition \(\lambda_{1,\varphi}(\Omega_*)\le M(s)\). Hence
\[
\lambda_{1,\varphi}(\Omega_*)=M(s),
\]
and \(\Omega_*\) is a global maximizer.
\end{proof}

We may now identify the global maximizers.

\begin{theorem}\label{thm:global-maximizers-are-disks}
For every \(s>0\), every global maximizer for
\[
\sup_{|\Omega|=s}\lambda_{1,\varphi}(\Omega)
\]
is, up to translation, a disk.
\end{theorem}

\begin{proof}
Let \(\Omega_*\) be a global maximizer of measure \(s\). Then \(\Omega_*\) is
automatically a local maximizer under the same measure constraint. Therefore the
rigidity theorem proved in the previous section applies: every local maximizer is,
up to translation, a disk.

Hence every global maximizer is a disk up to translation. Since the problem is
translation invariant, this completely characterizes the set of global maximizers.
\end{proof}

As an immediate consequence, we recover the classical Nicola--Tilli theorem.

\begin{corollary}
Among all measurable sets of prescribed measure, the disk maximizes Gaussian
time-frequency concentration:
\[
\lambda_{1,\varphi}(\Omega)\le \lambda_{1,\varphi}(D)
\]
whenever \(|\Omega|=|D|\) and \(D\) is a disk.
\end{corollary}

\begin{proof}
Existence of a global maximizer is given by
Theorem~\ref{thm:existence-global-maximizer}, and
Theorem~\ref{thm:global-maximizers-are-disks} shows that every such maximizer is a
disk up to translation. Since translation preserves both measure and the value of
\(\lambda_{1,\varphi}\), the claim follows.
\end{proof}

\bibliography{biblio}
\bibliographystyle{amsplain}

\end{document}

\section{Results}

\noindent{\bf Question:} The results from \cite{Abreu-Doerfler} deal with the \emph{simply connected} case. Can a similar result hold for more general domains?  

\begin{theorem}\label{thm:inverse-disk} Suppose $U$ is a connected, bounded domain of class $C^2,$ and that one of the eigenfunctions of the localisation operator associated to $U$ is a Hermite function. Then $U$ is an annulus. 
\end{theorem}

%\noindent{\bf Question:} Does Stability hold for the main result in \cite{Abreu-Doerfler}? 

%\begin{theorem}\label{thm:stability-inverse} For each $N>0,$ there exists $\delta_0(N)>0$ such that the following holds. 

%Let $U$ be a simply connected domain. Suppose $f$ is an eigenfunction of $\mathcal{L}_U,$ where $f$ is a finite linear combination of the first $N$ Hermite functions. If $\|f-h_i\|_2 \le \delta<\delta_0(N),$ then we have that $|\Omega \, \triangle \,D| \lesssim \omega(\delta).$
%\end{theorem}

\noindent{\bf Question:} Can we construct domains with prescribed eigenfunctions, as long as the eigenfunction is sufficiently close to a Hermite function? 

\begin{theorem}
\label{thm:local-perturbative}
Given $\lambda \in (0,1)$ and $N \in \mathbb{N},$ there is $\varepsilon_0(\lambda,N)>0$ such that, if a function $f_0 = h_0 + \sum_{k=0}^N a_k h_k$ satisfies $\max |a_k| < \varepsilon_0,$ then there is a domain $U_{\lambda}$ such that 
\[
\mathcal{L}_{U_{\lambda}} f_0 = \lambda f_0. 
\]
\end{theorem} 

\begin{theorem} Corollary 1.4 from \cite{Gomez-Guerra-Ramos-Tilli} is sharp. 
\end{theorem}

\section{Proofs}

\begin{proof}[Proof of the first result:] The Bargmann transform of a Hermite function $h_k$ is just a multiple of the monomial $z^k;$ thus, we may rephrase the statement that $h_k$ is an eigenfunction of the localisation operator $\mathcal{L}_U$ as the following equality, as in \cite{Abreu-Doerfler}: 
\[
\int_{U} z^k e^{\pi w \overline{z}} e^{-\pi |z|^2} \, dA(z) = \lambda w^k.
\]
Both sides clearly define holomorphic functions of the variable $w \in \C$. Thus, differentiating $k$ times in $w,$ we obtain that
\begin{equation}\label{eq:fourier-delta} 
\int_{\C} 1_U(z) |z|^{2k}e^{-\pi |z|^2} e^{\pi \overline{z} w} \, dA(z) = c(\lambda,k).
\end{equation} 
On the other hand, consider now the (compactly supported) distribution 
$$\eta :=  |z|^{2k} e^{-\pi |z|^2}1_U(z) - c(\lambda,k) \delta_0.$$
If we write $e^{\pi \overline{z} w} = e^{\pi (x-iy)w} = e^{\pi x \cdot w} \cdot e^{- i \pi y \cdot w}$ and consider the Fourier transform $\widehat{\eta}$ -- which, by the Paley-Wiener result, is an entire function on $\C^2$ --, then we have that $\widehat{\eta}(w,\pm i w) = 0, \, \forall w \in \C.$\footnote{The fact that $\widehat{\eta}(w,-iw) = 0$ follows directly from \eqref{eq:fourier-delta}. The other relation follows from the observation that $c(\lambda,k) \in \R,$ together with substituting $w \mapsto \overline{w}$ in \eqref{eq:fourier-delta} and taking conjugates on both sides} 

We now have to resort to an abstract result of Ehrenpreis-Malgrange flavour. 

In order to state it, note that we denote a distribution $\omega$ to be a \emph{singular-continuous free measure} if $\omega$ defines a Radon measure on $\R^2,$ whose Lebesgue decomposition with respect to the Lebesgue measure has only absolutely continuous and pure point parts, and only finitely many atoms compose the pure point part.

\begin{lemma} Let $\omega \in \mathcal{S}'(\R^2)$ be a compactly supported singular-continuous free distribution  such that its Fourier transform satisfies $\widehat{\omega}(w,\pm i w) = 0, \, \forall w \in \C.$ Then there is a finite set $P$ and a compactly supported function $u \in L^2(\R^2)$ such that $\Delta u = \omega,$ and $u \in C(\R^2 \setminus P).$ 
\end{lemma}

\begin{proof} By the Paley--Wiener theorem and the hypotheses, the function $\widehat{\omega}(w,z)$ is analytic of order $1$ on $\C^2.$ Thus, it admits the power series expansion 
\[
\widehat{\omega}(w,z) = \sum_{k,l \ge 0} a_{k,l} z^k w^l.
\]
By the hypothesis of vanishing on $\{(w,\pm iw), w \in \C\},$ then we get that 
\[
0 = \widehat{\omega}(w,\pm iw) = \sum_{n \ge 0} \left(\sum_{k,l \ge 0, k+l=n} a_{k,l} (\pm i)^k \right) w^n,
\]
and thus, by comparing coefficients for all $n \ge 0,$ we obtain that 
\[
\sum_{k,l \ge 0, k+l=n} a_{k,l} (\pm i)^k = \sum_{k=0}^n a_{k,n-k} (\pm i)^k =  0, \, \forall n \ge 0. 
\]
 Thus, rewriting $\widehat{\omega}$ as
\begin{align*} 
\widehat{\omega}(w,z) &= \sum_{n \ge 0} \left( \sum_{k=0}^n a_{k,n-k} \left(\frac{z}{w}\right)^k \right) w^n,  
\end{align*} 
the polynomials $p_n(t) = \sum_{k=0}^n a_{k,n-k} t^n$ always vanish at $t = \pm i,$ and thus we may write $p_n(t) = (t-i)(t+i) q_n(t),$ where the degree of $q_n$ is less than or equal to $n-2.$ In other words, 
\[
\widehat{\omega}(z,w) = \sum_{n \ge 0} w^{n-2} (z^2 + w^2) q_n(z/w).
\]
Direct inspection shows that $q_n \equiv 0$ if $n=0,1.$ As $q_n(z/w) w^{n-2}$ is again a polynomial in the variables $z,w,$ we see that there is a holomorphic function $F$ defined over $\C^2$ such that $\widehat{\omega}(z,w) = (z^2+w^2)F(z,w).$ By the conditions on $\omega$ again,  one has that $|\widehat{\omega}(z,w)| \le C e^{c(|\text{Im}(z)| + |\text{Im}(w)|)}.$ Thus, $F$ satisfies $|F(z,w)| \le C' \frac{e^{c(|\text{Im}(z)|+|\text{Im}(w)|)}}{(1+|w|+|z|)^2}.$ For $z,w \in \R,$ then we obtain that $|F(z,w)|\le \frac{C''}{1+|w|^2+|z|^2},$ and thus $F \in L^2(\R^2).$ This shows that $\widehat{F} =: u \in L^2(\R^2)$ satisfies $\Delta u = \omega$ in the distributional sense, but this is equivalent to 
\[
u(x,y) = c \cdot G * \omega (x,y), \text{ for some } c \in \R. 
\]
where $G(z) = \log|z|$ denotes the Green function on $\R^2.$ As $\omega$ is singular-free, then $G * \omega \in C(\R^2 \setminus P),$ where $P$ is the support of the point masses in the pure point part of $\omega.$  This concludes the proof. 
\end{proof}

Using this result in our case, we manage find some $u \in L^2(\R^2)$ which satisfies
\begin{equation}\label{eq:potato}
\Delta u = f(|z|)1_U - c(\lambda,k) \delta_0
\end{equation}
in the sense of distributions, where $f(|z|) = |z|^{2k} e^{-\pi |z|^2}.$ Writing $u$ in explicit form as above, we obtain that the domain $U$ satisfies 
\begin{equation}\label{eq:potential}  
\int_U \log|z-z'| f(|z'|) \, dA(z') = c(\lambda,k) \log|z|, \, \forall \, z \in \C\setminus (U \cup \{0\}).
\end{equation} 
We note the following: if $0 \in U,$ we consider a small disk centered at $0$ and contained in $U,$ and remove it from $U$. This yields a new domain $U', 0 \not\in U',$ which satisfies the property \eqref{eq:potential}. So, we may always suppose $0$ is not in the set $U$. Moreover, $0 \not\in \partial U,$ as the left-hand side of \eqref{eq:potential} converges to a finite quantity as $z \to 0,$ while the right-hand side diverges. 

Now, consider the function $\int_U \log|z-z'| f(|z'|) \, dA(z') = \Psi_U(z).$ We claim that $\Psi_U$ is \emph{locally radial}: for each $z \in U,$ there is a ball $B(z) \subset U$ such that $\Psi_U$ agrees with a radial function on $B(z).$ The proof of this fact will follow closely the ideas from Aharonov, Schiffer and Zalcman in \cite{Aharonov-Schiffer-Zalcman}.

Indeed, consider the auxiliary function $\Phi(z) = \int_B \log|z-z'|f(|z|) \, dA(z'),$ where $B \supset U$ is a ball centered at $0.$ We have immediately that $\Phi$ is radial, and as such $\nabla \Phi(z)$ is parallel to $z$, for all $z \in \R^2$. Moreover, if $z \in B,$ we have $\Delta \Phi (z) = c f(|z|).$ By the properties of $\Psi_U,$ difference $\Psi_U - \Phi$ still satisfies that $\nabla (\Psi_U - \Phi)(z) $ is parallel to $z$ if $z \in \partial U,$ by continuity of the gradient, and $\Delta(\Psi_U - \Phi) \equiv 0$ inside $U.$ The following lemma shows that this can only happen if $\Psi_U - \Phi$ is \emph{locally radial}, and thus $\Psi_U$ is locally radial:

\begin{lemma} Let $V \subset \R^2$ be a domain, and suppose that there exists a harmonic function $v \in C^1(\overline{V})$ with the property that $\nabla v(x)$ is parallel to $x$ at any $x \in \partial V.$ Then for each sphere $S \subset \R^2,$ $v$ is constant on the connected components of the intersection $S \cap V.$ 
\end{lemma}
 
\begin{proof} Consider the functions $g_{i,j} = x_i \partial_j u - x_j \partial_i u.$ By the hypotheses, $g_{i,j}$ are harmonic functions in $V,$ with $g_{i,j} \equiv 0$ on $\partial V.$ Thus, $g_{i,j} \equiv 0$ accross $V$, for each pair of indices $i,j.$ This implies, in particular, that $\nabla u $ is parallel to $x$ for \emph{all} $x \in V.$ That is, $\nabla u(x) \in T_x^{\perp}S, \forall x \in S\cap V.$ This immediately implies that $u$ is \emph{constant} on each connected component of $S\cap V,$ as desired.
 \end{proof}

We may now conclude: indeed, if $U$ is not an annulus, then there are $r_1,r_2>0$ such that the annulus $A = \{z \in \R^2 \colon r_1 < |z| < r_2\}$ intersects $\partial U$ and $U$. Let $x_0 \in (\partial U) \cap A$ be an arbitrary point, and consider the set $\tilde{U}_0 = U \cap B_r(x_0),$ for $r >0$ sufficiently small. We claim that $\Psi_U $ agrees with a radial function on $\tilde{U}_0:$ the set $\tilde{U}_0$ is connected for $r$ sufficiently small, and by the property of local radiality, one concludes the claim. Let then $\psi_U(|x|) = \Psi_U(x)$ over the set $\tilde{U}_0.$

By continuity of the gradient, one has that $c \frac{z}{|z|^2} =\nabla \Psi_U(z) = \psi_U'(|z|) \cdot \frac{z}{|z|},$ which implies $\psi_U'(|z|) = \frac{c}{|z|}$ for $z \in \partial \tilde{U}_0,$ and thus also for $z \in \tilde{U}_0.$ But then a simple computation of the laplacian of $\Psi_U$ on $\tilde{U}_0$ implies $\Delta \Psi_U \equiv 0$ in $\tilde{U}_0,$ which contradicts the fact that $\Delta \Psi_U = c' \cdot f(|z|),$ as $f$ only vanishes at the origin. 
 \end{proof}

\section{Proof of Theorem 4} 

 \begin{proof}[Proof of the second result] Again, we employ the same Bargmann transform method used in the previous result: let $P(z) = 1 + \sum_{k =1}^N b_k z^k,$ where $b_k$ is such that the Bargmann transform of $a_k h_k$ equals $b_k z^k$. The problem of finding a domain $U$ with $f_0$ as eigenfunction of $\mathcal{L}_U$ is equivalent to the following relationship:
\begin{equation}\label{eq:eigenfunction-general}  
\int_U P(z) e^{-\pi|z|^2} e^{\pi \overline{z} w} \, dA(z) = \lambda P(w). 
\end{equation} 
\noindent{\bf Step 0:} First reduction. 

For a polynomial $R \in \C[z],$ given by $R(z) = \alpha_0 + \alpha_1 z + \cdots + \alpha_n z^n,$ define the differential operator $R(\partial_w/\pi) = \sum_{k=0}^n \alpha_k \left( \frac{\partial_w}{\pi}\right)^k.$ We then apply the differential operator $P\left(\frac{\partial_w}{\pi}\right)$ to both sides in \eqref{eq:eigenfunction-general}. Denoting the new polynomial $P(\partial_w/\pi)(P(w)) =: Q(w),$ we need to find $U$ such that 
\begin{equation}\label{eq:almost-reduced} 
\int_U |P(z)|^2 e^{-\pi |z|^2} e^{\pi \overline{z}w} \, dA(z) = \lambda Q(w).
\end{equation}

\noindent \textbf{Step 1:} There is a domain $U_0$ whose boundary is of class $C^2,$ such that 
\begin{equation}\label{eq:first-domain}
\int_{U_0} |P(z)|^2 e^{-\pi |z|^2} e^{\pi \overline{z} w} dA(z) = \lambda.
\end{equation}
In order to prove this step, we shall resort to a perturbative result by Isakov \cite{Isakov}. Before we are able to use it, let us provide the setting for its application.

Let $\psi \in C^{\infty}_c(\C)$ be a bump function with the property that $\psi \equiv 0$ in $D_1(0),$ $\psi$ radial, monotone, and $\psi \equiv 0$ outside $D_2(0).$ Let then $\psi_{\delta}(z) = \frac{1}{\delta^2} \psi(z/\delta).$ Consider then the integral 
\[
\int \psi_{\delta}(z) e^{-\pi|z|^2} e^{\pi \overline{z} w} dA(z) = F_{\delta}(w).
\]
As $\psi$ is radial, a simple computation with polar coordinates shows that $F_{\delta}(w) = c_{\delta},$ where $c_{\delta} = \int \psi_{\delta}(z) e^{-\pi |z|^2} dA(z).$ Thus, we may restate \eqref{eq:first-domain} as 
\begin{equation}\label{eq:second-domain} 
\int_{U_0} |P(z)|^2 e^{-\pi |z|^2} e^{\pi \overline{z} w} dA(z) = \frac{\lambda}{c_{\delta}} \int \psi_{\delta}(z) e^{-\pi|z|^2} e^{\pi \overline{z}w} dA(z).
\end{equation}
Using Proposition \ref{prop:auxiliary}, finding $U_0$ is equivalent to finding a compactly supported solution $u_0$ to 
\begin{equation}\label{eq:inverse-first} 
\Delta u_0 = - |P(z)|^2 e^{-\pi |z|^2}1_{U_0} + \frac{\lambda}{c_{\delta}}\psi_{\delta}(z) e^{-\pi|z|^2}. 
\end{equation}
In order to find such solution, first notice that a similar problem admits a solution if we replace $P$ by $1$: indeed, a certain disk $D_0$ satisfies, by the results by Daubechies \cite{Daubechies}, that 
\[
\int_{D_0} e^{-\pi|z|^2} e^{\pi \overline{z}w} \, dA(z) = \lambda,
\]
and thus we may find a compactly supported solution $\tilde{u}_0$ of 
\begin{equation}
\Delta \tilde{u}_0 = -e^{-\pi|z|^2}1_{D_0} + \frac{\lambda}{c_{\delta}}\psi_{\delta}(z) e^{-\pi|z|^2}.
\end{equation}
Moreover, if $\delta >0$ is chosen sufficiently small, then the support of $\psi_{\delta}$ is contained in $D_0.$ This solution is radial and, furthermore, if $\delta >0$ is small, $\tilde{u}_0 <0$ on its support. Now we are able to use the machinery developped in \cite[Chapter~5]{Isakov}.

\begin{theorem}[Corollary 5.1.4 in \cite{Isakov}]\label{thm:isakov} Consider a $C^2$ domain $\Omega_0 \subset \C$ and fix $f_0 \in C^5(\R^2)$ with $f_0 > 0$ on $\partial \Omega_0.$ Suppose there exists a solution $v$ to 
\begin{enumerate}
    \item[(i)] $-\Delta v = f_0 1_{\Omega_0}$ in $\C,$ with $v = 0$ outside $\Omega_0;$
    \item[(ii)] $v <0$ in $\Omega_0.$
\end{enumerate}
Then there is a neighborhood $\Omega' \supset \Omega_0$ and $\varepsilon_1 (\Omega_0) > 0$ such that, if 
\[
\|f-f_0\|_{C^5(\Omega')} < \varepsilon_1,
\]
there is a $C^2$ domain $\Omega_f$ --- which is a perturbation of $\Omega_0$ --- for which there is a solution $\tilde{v}$ of 
\begin{enumerate}
    \item[(i)] $-\Delta \tilde{v} = f1_{\Omega_f}$ in $\C,$ with $v = 0$ outside $\Omega_f;$
    \item[(ii)] $\tilde{v} <0$ in $\Omega_f.$
\end{enumerate}
Moreover, the solution is unique among the class of regular perturbations of the domain $\Omega_0.$
\end{theorem} 
Let $f_0(z) = e^{-\pi|z|^2}\left(1 - \frac{\lambda}{c_{\delta}}\psi_{\delta}(z)\right).$ By the previous considerations, $f_0$ satisfies the conditions of Theorem \ref{thm:isakov}. We showed also that the function $\tilde{u_0}$ also matches the conditions in that result. Now, if we let $f(z) = e^{-\pi|z|^2}\left(|P(z)|^2 - \frac{\lambda}{c_{\delta}}\psi_{\delta}(z)\right),$ a direct computation shows that the condition that $\|f-f_0\|_{C^5(\Omega')}$ is small can be achieved for any disk $\Omega' \supset D_0,$ as long as the coefficients of $P$ are chosen sufficiently small (note that this choice is independent of $\delta>0$). Thus, Theorem \ref{thm:isakov} gives us the existence of a function $u_0$ with $\Delta u_0 = f(z)1_{U_0},$ where $U_0$ is the perturbed domain given in the statement of that result. If $\varepsilon_1>0$ is taken sufficiently small, then $U_0$ contains the support of $\psi_{\delta}$ still, and hence $u_0$ satisfies \eqref{eq:inverse-first}. \\

\noindent{\bf Step 2:} There is a domain $U$ which satisfies \eqref{eq:almost-reduced}. 

In order to complete this step, we need to invoke a result on existence of quadrature domains, originally attributed to V. K. Ivanov in the context of harmonic functions, and proved in its present form by Cherednichenko \cite{Cherednichenko}. In order to state it, we must introduce some notation: given a positive measurable function $f:\C \to \R_+$ and a domain $U \subset \C,$ we define the following Cauchy-like transform with weight $f$ and associated with $U:$
\begin{equation}\label{eq:cauchy-domain} 
g_U(z) = \int_U \frac{\mu(w)}{z-w} \, dA(w).
\end{equation}
It is easy to see that $g_U$ is \emph{analytic} on the complement of $U,$ and moreover it satisfies $\partial_{\overline{z}} g_U = \pi \cdot \mu$ on $U.$ 

%Note furthermore that, if we take $f$ radial and $U$ to be a disk, then $g_U$ satisfies $g_U(z) = \frac{\lambda_U}{z}, \, z \in \C \setminus U,$ where $\lambda_U = \int_U f.$ 

\begin{proposition}[Theorem 2.1 in \cite{Cherednichenko}]\label{prop:auxiliary} Fix a domain $U_0 \subset \C$ of class $C^2,$ and two other domains $U_1 \subset U_0 \subset U_2$ possessing the same regularity. Suppose one is given an analytic function $g:\C \setminus U_1 \to \C,$ with the property that 
\begin{equation}\label{eq:first-conditions-g}
\lim_{z \to \infty} g(z) = 0, \,\,\, |g(z)| \le C_0 \text{ for } z \in \C \setminus U_1, \,\,\, \lim_{z \to \infty} z g(z) = c_0 > 0.
\end{equation} 
There is $\varepsilon_0 = \varepsilon_0(U_0,U_1,C_0,\mu) > 0$ such that, if 
\begin{equation}\label{eq:closeness-condition-g}
\sup_{z \in \partial U_1} |g(z) - g_{U_0}(z)| \le \varepsilon_0,
\end{equation} 
then there is a $C^2$ domain $\tilde{U}$ with $U_1 \subset \tilde{U} \subset U_2$ which satisfies moreover 
\[
g_{\tilde{U}}(z) = g(z), \, \forall \, z \in \C \setminus \tilde{U}. 
\]
\end{proposition}

Returning to the task at hand, restricting \eqref{eq:almost-reduced} to $t \in \R_+,$ multiplying both sides agains $e^{-\pi \zeta t}$ and integrating over $t > 0,$ we obtain 
\[
\int_{U} \left(\int_0^{\infty} e^{\pi (\overline{z} - \zeta) t} dt \right) |P(z)|^2e^{-\pi|z|^2} \, dA(z) = \lambda \int_0^{\infty} e^{-\pi \zeta t} Q(t) \, dt. 
\]
Write $Q(w) = \alpha_0 + \alpha_1 w + \cdots + \alpha_n w^n.$ We may rewrite the equation above as 
\[
\int_{U} \frac{1}{\pi(\zeta-\overline{z})} |P(z)|^2e^{-\pi|z|^2} \, dA(z) = \lambda_0 \left( \sum_{k=0}^n \alpha_k \int_0^{\infty} e^{-\pi \zeta t} t^k \, dt \right) = \lambda \left( \sum_{k=0}^n \alpha_k \zeta^{-k} \frac{\Gamma(k+1)}{\pi^{k+1}} \right),
\]
where we are taking $\zeta \in \R_+, \, \zeta > \sup_{w \in U} |w|.$ Replacing $U$ by its conjugate then yields that 
\begin{equation}\label{eq:analytic-construct} g_{U}(\zeta) = \pi \cdot \lambda \left( \sum_{k=0}^n \alpha_k \zeta^{-k} \frac{\Gamma(k+1)}{\pi^{k+1}} \right)
\end{equation} 
for all $\zeta \in \R_+$ sufficiently large, where $g_U$ is as in \eqref{eq:cauchy-domain} with $\mu(w) = |P(\overline{w})|^2 e^{-\pi|w|^2}.$ By analytic continuation, we obtain that both functions have to be equal over $\C \setminus U.$ Thus, in order to construct $U$, we must construct a conjugate domain satisfying \eqref{eq:analytic-construct}. Conversely, by series expansion, if a domain $U$ satisfies \eqref{eq:analytic-construct}, then its conjugate satisfies \eqref{eq:almost-reduced}, by a series expansion together with an analytic continuation argument. 

We denote the function on the right-hand side of \eqref{eq:analytic-construct} by $g(z).$ A quick analysis shows that, as $Q(w) = P(\partial_w/\pi) (P(w)),$ then, as $P(0) = 1,$ $Q(0) = 1 + \sum_{k=1}^N b_k^2.$  We wish to use Proposition \ref{prop:auxiliary} with $U_0$ and the function $g$ as above. 

Theorem \ref{thm:isakov} yields a natural choice for the domains $U_1$ and $U_2$. Indeed, as $U_0$ is a small perturbation of the disk $D_0,$ we may take $U_1$ and $U_2$ to be disks with, say, $1-10^{-3}$ and $1+10^{-3}$ times the radius of $D_0,$ respectively. If the coefficients $b_k$ of $P$ are sufficiently small --- in terms of $\lambda$ ---, it is clear that \eqref{eq:first-conditions-g} and \eqref{eq:closeness-condition-g} are satisfied by $g,$ as the leading coefficient of $Q$ is much larger than the other ones. Hence a domain with the property \eqref{eq:analytic-construct}, and thus also a domain with the property \eqref{eq:almost-reduced}, exists, as long as the coefficients of $P$ are taken to be small enough as a function of $\lambda.$

The domain $U$ so obtained is such that $U \subset U_2 \subset 2D_0$ in particular. As $\text{diam}(D_0) \to 0$ as $\lambda \to 0,$ and as the minimal norm of zeros of $P$ converges to $\infty$ as $\lambda \to 0,$ then the set $U$ does not contain zeros of $P.$ \\ 

\noindent{\bf Step 3:} The domain $U$ constructed above also satisfies \eqref{eq:eigenfunction-general}.  

Let $\int_{U} P(z)e^{-\pi |z|^2} e^{\pi \overline{z}w} dA(z) = G(w).$ By \eqref{eq:almost-reduced}, we obtain that 
\begin{equation}
P(\partial_w/\pi) G(w) = P(\partial_w/\pi)(\lambda P)(w), \, \forall w \in \C. 
\end{equation}
 Then it follows that $G - \lambda P =: H$ satisfies the equation 
\[
P(\partial_w/\pi) H = 0.
\]
Letting
\[
\tilde{P}(t) = 1 + \sum_{k=0}^n \frac{b_k}{\pi^k} t^k,
\]
suppose $\tilde{P}$ has zeros $z_1,\dots,z_l$ with multiplicities $m_1,\dots,m_l.$  Since $H$ has at most exponential growth in $\C,$ it follows from standard ODE theory that 
\[
H(w) = \sum_{i=1}^l R_i(w) e^{z_i w},
\]
where each $R_i$ is a polynomial of degree $\text{deg}(R_i)<m_i.$ By the observations in the previous part, we may suppose that $|z_i| > 10 \cdot \sup_{z \in U} |z|.$ 

Suppose without loss of generality that $z_1$ has the maximum modulus among all $z_i,$ that is, $|z_1| \ge |z_j|, \, \forall j \in \{1,\dots,l\}.$ We shall assume that $R_1 \not \equiv 0,$ and derive a contradiction. 

In order to do so, let then $w = t \overline{z_i}, \, t > 0.$ On the one hand, we have 
$$|H(t \overline{z_1})| \ge c \cdot e^{t|z_1|^2}$$
for $t$ sufficiently large. On the other hand, since we know that $P$ is a polynomial, and we have the bound 
\[
|G(w)| \le C_{\varepsilon} |U| e^{(\pi+\varepsilon) M |w|},
\]
for each $\varepsilon>0,$ where $M = \sup_{z \in U} |z|,$ then we conclude that 
\[
|H(t\overline{z_1})| \le C|U| e^{2\pi M t |z_1|}.
\]
This shows that $|z_1| \le 2 \pi M,$ which is an obvious contradiction to the fact that the zeros of $P$ are far away from the domain $U.$ Hence, $R_1 \equiv 0.$

Iterating such an argument for the remaining zeros of $\tilde{P},$ we conclude that $R_i \equiv 0$ for all 
$i \in \{1,\dots,l\}.$ In other words, we have concluded that $G = \lambda P,$ which shows that $P$ is an eigenfunction of the localization operator of the domain $U$ associated with the eigenvalue $\lambda.$ Since this was exactly what we aimed to prove, we have concluded the proof. 

\section{Proof of Stability of Corollary 1.4 in \cite{Gomez-Guerra-Ramos-Tilli}}

We start by defining $P_{\varepsilon}(w) = 1+ \varepsilon w^2.$ Fix $\lambda_0 = 1 - e^{-\pi}.$ 

According to Theorem 4, we know that, for each $\varepsilon > 0$ sufficiently small, there is a domain $U_{\varepsilon},$ which moreover may be taken to be a perturbation of the form 
\begin{equation}\label{eq:def-u-eps}
\partial U_{\varepsilon} = \{ (1+\varphi_{\varepsilon}(y)) \cdot y, y \in \mathbb{S}^1\}
\end{equation}
of the unit disk, such that 
\begin{equation}\label{eq:functiona-P-eps}
\int_{U_{\varepsilon}} P_{\varepsilon}(z) e^{\pi \overline{z} w} e^{-\pi |z|^2} dA(z) = \lambda P_{\varepsilon}(w), \, \forall \, w \in \mathbb{C}. 
\end{equation}
The assertion about $U_{\varepsilon}$ being a perturbation of the unit ball of the form \eqref{eq:def-u-eps} may be proved as follows: first, notice that the method of proof of Theorem \ref{thm:isakov} yields directly that the domains so obtained are of the desired form. A thorough inspection of the proof of Proposition \ref{prop:auxiliary} shows that the same holds for the domains constructed there - that is, the domain $\tilde{U}$ from Proposition \ref{prop:auxiliary} may be taken to be a normal perturbation of the domain $U_0$ in that result. If both perturbations are small enough in absolute norm, we readily conclude that the domains constructed in Theorem 4 are still a smooth perturbation over the unit disk. 

Let us suppose that $\varphi_{\varepsilon}$, defined as in \eqref{eq:def-u-eps}, are \emph{differentiable} in the parameter $\varepsilon \in (-\delta,\delta).$ In that case, let $f_0(y) := \partial_{\varepsilon}|_{\varepsilon = 0} \varphi_{\varepsilon}.$ A simple use of Reynold's formula in \eqref{eq:functiona-P-eps} shows, through differentiation in $\varepsilon$ and taking $\varepsilon \to 0,$ that 

\begin{align*} 
\int_{\mathbb{S}^1} f_0(z) e^{\pi \overline{z} w}\,e^{-\pi|z|^2}\, d\mathcal{H}^1(z) & = -\int_{U_0} z^2  e^{\pi \overline{z}w} e^{-\pi |z|^2} \, dA(z) + \lambda w^2 \cr 
    & = (\lambda -\tilde{\lambda})w^2,
\end{align*} 
where $\lambda > \tilde{\lambda} > 0.$ It readily follows then by setting $w=0$ that
\begin{equation}\label{eq:cancellation-f-0}
0 = \int_{\mathbb{S}^1} f_0(y) \, d \mathcal{H}^1(y).
\end{equation}

On the other hand, we note that a simple use of polar coordinates with \eqref{eq:def-u-eps} implies that 
\begin{equation}\label{eq:area-u-eps} 
|U_{\varepsilon}| = \frac{1}{2} \int_{\mathbb{S}^1} (1+\varphi_{\varepsilon}(y))^2 \, d \mathcal{H}^1(y),
\end{equation}
for all $\varepsilon \in (-\delta,\delta).$ Since $\varepsilon \mapsto \varphi_{\varepsilon}$ is smooth, we may conclude from \eqref{eq:cancellation-f-0} and the formula above that
\begin{equation}\label{eq:measure-close-u-eps} 
\left||U_{\varepsilon}| - |\mathbb{S}^1| \right| \le C \varepsilon^2,
\end{equation}
since we have $|\varphi_{\varepsilon}(y) - \varepsilon \cdot f_0(y)| = O(\varepsilon^2)$ by a simple Taylor expansion. We then have that 
\[
\left| 1 - \frac{1-e^{-1}}{1-e^{-|U_{\varepsilon}|}} \right| \le C' \varepsilon^2. 
\]
Finally, a similar polar coordinates argument to that used in order to obtain \eqref{eq:area-u-eps} can be used in order to deduce that 
\[
C'\varepsilon \ge C \varepsilon \int_{\mathbb{S}^1} |f_0(y)| \, d \mathcal{H}^1(y)  \ge  \left| (r \cdot U_{\varepsilon}) \, \triangle \,\, \mathbb{S}^1\right| \ge c \varepsilon \int_{\mathbb{S}^1} |f_0(y)| \, d \mathcal{H}^1(y) \ge c_0 \varepsilon,
\]
where $r = |U_{\varepsilon}|^{-1/2}.$ This follows by the fact that, as we established in \eqref{eq:measure-close-u-eps}, 
$$|U_{\varepsilon}| = 1 + O(\varepsilon^2) \, \Rightarrow |U_{\varepsilon}|^{1/2} = 1 + O(\varepsilon^2).$$
In order to show that the unit disk is the best possible approximating disk to $U_{\varepsilon},$ we note the following: by triangle inequality for the symmetric difference, we have that, if $D'$ is a different disk of unit radius, 
\[
| (r \cdot U_{\varepsilon}) \, \triangle \,\, D'| \ge \left| D' \, \triangle \,\, \mathbb{S}^1\right|-\left| (r \cdot U_{\varepsilon}) \, \triangle \,\, \mathbb{S}^1\right|.
\]
If $|D' \, \triangle \,\, \mathbb{S}^1| \ge 2C'\varepsilon,$ then we are done. On the other hand, by the triangle inequality again, we get that if $|D' \, \triangle \, \mathbb{S}^1| \le (c_0/2) \varepsilon,$ then we also get the desired lower bound. Hence, we may assume without loss of generality that 
\[
\frac{|D' \, \triangle \, \mathbb{S}^1|}{\varepsilon} \in (c_0/2,2C'). 
\]
We then fix an angle $\theta \in (0,2\pi)$ and a constant $\eta \in (c_0/2,2C'),$ and let $D(\theta,\eta)$ denote the disk of unitary volume, such that $c(D(\eta,\theta)) = \eta e^{i\theta} \cdot \varepsilon,$ where $c(D')$ denote the center of the disk $D'.$

Note hence that for $\varepsilon>0$ sufficiently small, then we know that $U_{\varepsilon}$ is still given by a graph over the boundary of each disk $D(\theta,\eta).$  Define then $\psi_{\varepsilon}^{\theta,\eta}:\mathbb{S}^1 \to \R$ to be such that 
\[
\partial(U_{\varepsilon} - c(D(\theta,\eta))) = \{ (1+\psi^{\theta,\eta}_{\varepsilon}(y)) \cdot y, y \in \mathbb{S}^1\}.
\]
We then repeat the same argument we used for $\varphi_{\varepsilon}$: let $\partial_{\varepsilon}|_{\varepsilon=0} \psi_{\varepsilon}^{\theta,\eta} = g^{\theta,\eta}.$ Moreover, note that the velocity of the domain derivative at $\varepsilon=0$ equals $\eta e^{i\theta} + g^{\theta,\eta}(y)\cdot y.$ Hence, by using these facts in \eqref{eq:functiona-P-eps}, differentiating that equation with respect to $\varepsilon,$ taking $\varepsilon\to 0,$ we obtain in the same fashion as before that 
\begin{align*} 
\int_{\mathbb{S}^1} g^{\theta,\eta}(z) e^{\pi \overline{z} w}\,e^{-\pi|z|^2}\, d\mathcal{H}^1(z) & = -\int_{U_0} z^2  e^{\pi \overline{z}w} e^{-\pi |z|^2} \, dA(z) + \lambda w^2 \cr 
& - \eta \int_{\mathbb{S}^1} \text{Re}(z\cdot e^{-i\theta}) e^{\pi \overline{z}w} e^{-\pi|z|^2} \, d\mathcal{H}^1(z) \cr 
    & = (\lambda -\tilde{\lambda})w^2 + \alpha(\theta,\eta) \cdot w,
\end{align*} 
where $\tilde{\lambda} < \lambda$ and $\alpha \in \C.$ This readily implies that $g^{\theta,\eta}(z) = \beta_1(\theta,\eta) z + \beta_2 z^2,$ with $\beta_2 > 0$ independent of $\eta,\theta.$ Since the space of such functions is two-dimensional, the $L^2$ norm on that space is equivalent to the $L^1$ norm, which implies that there is $\delta>0$ such that 
\[
\int_{\mathbb{S}^1} |g^{\theta,\eta}(z)| \, d\mathcal{H}^1(z) \ge \delta. 
\]
It then follows that    
\[
\int_{\mathbb{S}^1} |\psi_{\varepsilon}^{\theta,\eta}(y)| \, d \mathcal{H}^1(y) \gtrsim \varepsilon,
\]
where the implicit constant holds independently of $\theta,\eta$ as above. These considerations imply readily that 
$$\inf_{|D| =1} |(r \cdot U_{\varepsilon}) \, \triangle \,\, D| \gtrsim \varepsilon.$$
We hence obtain the desired result on the sharpness of Corollary 1.4 in \cite{Gomez-Guerra-Ramos-Tilli}, as desired. 


%We now replace our task of finding $U \subset \C$ with the property \eqref{eq:almost-reduced} with that of finding $U$ such that 
%\begin{equation}
%\int_U |P(z)|^2 e^{-\pi|z|^2} e^{\pi \overline{z} w} \, dA(z) = \frac{\lambda}{\lambda_0} \int_{D_0} e^{\pi \overline{z} w} \, dA(z).
%\end{equation}
%Reorganizing and using the same idea as in the proof of Theorem \ref{thm:inverse-disk}, it suffices then to find a set $U$ for which there exists a solution $u \in C_c(\R^2)$ of 
%\begin{equation}\label{eq:free-boundary}  
%\Delta u = -|P(z)|^2 e^{-\pi |z|^2} 1_U(z) + \frac{\lambda}{\lambda_0} 1_D(z). 
%\end{equation} 
%By taking the coefficients of order at least one of $P$ to be even smaller if needed, then we may find that $P$ does not have zeroes over $D_{10}(0).$ As we shall take $\lambda_0$ small at the end of the proof, we readily see that $U_2 \subset D_{10}.$ In order to construct the solution to \eqref{eq:free-boundary}, we shall resort to the following result from \cite{Gustafsson-Shahgholian} on existence of quadrature domains for general measures: 

%\begin{theorem}[Theorem 4.7 in \cite{Gustafsson-Shahgholian}]\label{thm:free-boundary} Let $\mu$ be a positive measure on $\C.$ Suppose one is given a function $h \in L^{\infty}(\C),$ with $0 \le h(z) \le b$ for all $z \in \C,$ and moreover assume $h$ is bounded from below outside a compact set. Suppose, moreover, that the measure $\mu$ has support on a disk $D_r(0),$ and satisfies 
%\[
%\mu(D_r(0)) \ge 40b |D_r(0)|. 
%\]
%Then there exists a domain $\Omega \subset \C,$ such that $\Omega \supset D_{3r}(0),$ and a compactly supported distribution $u \in \mathcal{S}'(\C)$ which satisfies 
%%\begin{equation}
%\Delta u = -h\cdot 1_{\Omega} + \mu
%\end{equation}
%in the sense of distributions. Moreover, if $\mu$ is absolutely continuous with respect to Lebesgue measure on $\C,$ then $u$ may be taken to be in $H^1(\C),$ $u \ge0,$ and $\Omega = \{u < 0\}.$ 
%\end{theorem}

%Let us now set the stage for the application of Theorem \ref{thm:free-boundary}. Let $f(z) : = |P(z)|^2 e^{-\pi |z|^2},$ and define $f_N$ suitably truncated versions of $f;$ in particular, we ask that $f_N$ is smooth over $\C,$ that $f_N$ is bounded from below, and that $f_N \equiv f$ on $D_N(0).$ 

%Let $\mu = \frac{\lambda}{\lambda_0} 1_{D_0},$ and $h = f_N$ in Theorem \ref{thm:free-boundary}. By the way we chose the approximations, there is an absolute constant $c_0 > 0$ such that $0 \le f_N \le c_0.$ Moreover, if we choose $\lambda_0$ sufficiently small, then $D_0 \subset D_{10 \sqrt{\lambda_0}}(0),$ and thus $\mu(D_0) = \lambda \ge 10^4 c_0 \lambda_0.$ Thus, we are able to find a solution $u_N \in H^1(\C), u_N$ compactly supported with $\Omega_N = \{u_N < 0\},$ and 
%\begin{equation}\label{eq:quadrature} 
%\Delta u_N = - f_N(z) 1_{\Omega_N} + \mu. 
%\end{equation} 
%In order to deduce that $u_N$ is a solution to our problem, for large enough $N,$ we only need to verify that $\Omega_N \subset D_N(0).$ We employ an argument originally contained in the works of Gustafsson \cite{Gustafsson} and Shahgholian \cite{Shahgholian}, but first we need to build the setting for it. 

%Notice first that, if $u_N \in \mathcal{S}'(\C)$ which solves \eqref{eq:quadrature} belongs to $C_c^1(\C)$ then $\Omega_N$ is a \emph{quadrature domain} for the measure $\mu$ with weight $f_N$: indeed, if $\varphi \in C(\overline{\Omega_N})$ is an arbitrary function, harmonic inside $\Omega_N,$ then 
%\[
%0 = \int_{\Omega_N} (\varphi \Delta u_N - u_N \Delta \varphi) = \int_{\Omega_N} (-f_N + \mu)\varphi,
%\]
%where the first equality is due to Green's identity (which holds, as $u_N \in C^1_c(\C)$) and the fact that $u = |\nabla u| = 0$ on $\partial \Omega_N$. On the other hand, if $\Omega_N$ is a quadrature domain for the measure $\mu$ with weight $f_N,$ if we let $K(w,z)$ denote the \emph{generalized Newton potential} in $\C$ (cf. \cite{Shahgholian}), then the function 
%\[
%\tilde{u}_N(z) = \int_{\Omega_N} K(w,z) f_N(w) \, dA(w) - \int_{\Omega_N} K(w,z) d\mu(w)
%\]
%still satisfies \eqref{eq:quadrature}  -- this is a simple computation with the definition of $K;$ we refer the reader to \cite{Shahgholian} or \cite{Karp}. Moreover, as $\Omega_N$ is a quadrature domain associated with $\mu$ and $f_N,$ it follows that, if $z \not \in \Omega_N,$ then $\tilde{u}_N(z) = 0.$ Thus, we may replace $u_N$ with $\tilde{u}_N$ everywhere above. 

%The previous step seems pointless at first, but 


%For that purpose, let $z_0 \in \partial \Omega_N.$ Take then an arbitrary sequence of points $z_j \in \Omega_N.$ By Theorem \ref{thm:free-boundary}, $u_N(z_j) > 0.$ Now consider the following function: 
%\[
%v_{N,j}(z) = u_N(z) - u_N(z_j) + c_N |z-z_j|^2,
%\]
%where $c_N >0$ is taken such that $f_N(z) \ge 20 c_N.$ Let $D_r(z_0)$ be a disk that does \emph{not} intersect the set $D_0.$ Then $v_{N,j}$ is \emph{subharmonic} on $D_r(z_0) \cap \Omega_N.$ We now use the maximum principle on that domain: if $z_j \in D_r(z_0),$ then we know that $v_{N,j}(z_j) = 0.$ Moreover, $v_{N,j} \ge 0$ on $(\partial \Omega_N)\cap D_r(z_0).$ As the boundary of $D_r(z_0) \cap \Omega_N$ is composed of two components - namely, $(\partial \Omega_N)\cap D_r(z_0)$ and $\Omega_N \cap (\partial D_r(z_0))$ - , then we conclude that 
%\[
%\inf_{z \in \Omega_N \cap (\partial D_r(z_0))} v_{N,j} (z) < 0.
%\]
%By taking $z_j \to z_0,$ this plainly implies that 
%\begin{equation} \label{eq:lower-bound-u_N}
%u_N(z) \le - c_N|z-z_0|^2 \text{ for some } z \in \Omega_N \cap (\partial D_r(z_0)),
%\end{equation} 
%\emph{as long as} $D_r(z_0)$ does not intersect $D_0.$

%Now let $D_1 \supset D_0$ be a disk. Suppose $\Omega_N \setminus D_1$ is nonempty. We again observe that $u_N$ is superharmonic on $\Omega_N \setminus D_1,$ $u_N = 0$ on $\partial \Omega_N,$ and $u_N < 0$ in $\Omega_N \setminus D_1,$ by \eqref{eq:lower-bound-u_N}. Thus, by the minimum principle for superharmonic functions, $u_N$ takes its minimum on $\Omega_N \cap (\partial D_1).$ But then: 
%\begin{align*} 
%c_N \text{dist}(z,D_1)^2 &\le \sup_{\Omega_N \cap (\partial D_{\text{dist}(z,D_1)}(z))} %(-u_N(z)) \cr 
 %                       &\le - \inf_{\Omega_N \setminus D_1} u_N(z) \le - \inf_{\partial D_1} u_N(z). 
%\end{align*} 
%We now wish to bound $|u_N(z)|$:the relevant construction seems to be that which removes some terms on the taylor expansion of the potential; as noted by Karp, the functions $u_N$ which arise that way satisfy the bound 
\end{proof}
